% =============================================================================
% Chapter 3 — Proposed Solution
% Include from main.tex via:  % =============================================================================
% Chapter 3 — Proposed Solution
% Include from main.tex via:  % =============================================================================
% Chapter 3 — Proposed Solution
% Include from main.tex via:  % =============================================================================
% Chapter 3 — Proposed Solution
% Include from main.tex via:  \input{chap3}
% Required: amsmath, booktabs, graphicx, tikz, backgrounds, algorithm, algpseudocode
%   \usetikzlibrary{positioning, arrows.meta, fit, calc, backgrounds}
% =============================================================================

%---- Colour palette for the diagrams ----------------------------------------
\definecolor{sosblue}{RGB}{38,84,124}
\definecolor{sosgreen}{RGB}{61,133,92}
\definecolor{sosorange}{RGB}{200,120,40}
\definecolor{sosgrey}{RGB}{90,90,90}
\definecolor{soslight}{RGB}{234,240,245}
\definecolor{soscream}{RGB}{248,244,232}

%---- Reusable TikZ node styles ----------------------------------------------
\tikzset{
  sensor/.style   = {rectangle, rounded corners=2pt, draw=sosblue, very thick,
                     fill=soslight, align=center, minimum width=2.5cm,
                     minimum height=1cm, font=\small},
  actuator/.style = {rectangle, rounded corners=2pt, draw=sosorange, very thick,
                     fill=soscream, align=center, minimum width=2.5cm,
                     minimum height=1cm, font=\small},
  core/.style     = {rectangle, rounded corners=4pt, draw=sosgreen, line width=1.2pt,
                     fill=sosgreen!12, align=center, minimum width=3.2cm,
                     minimum height=1.3cm, font=\small\bfseries},
  mission/.style  = {ellipse, draw=sosgrey, very thick, fill=white, align=center,
                     font=\small\itshape, minimum width=3cm},
  flow/.style     = {-{Stealth[length=2.5mm]}, line width=0.9pt, sosgrey},
  ctrl/.style     = {-{Stealth[length=2.5mm]}, line width=0.9pt, sosorange, dashed},
  block/.style    = {rectangle, rounded corners=2pt, draw=sosblue, thick,
                     fill=soslight, align=center, font=\small,
                     minimum width=2.6cm, minimum height=0.9cm},
}


\chapter{Proposed Solution}
\label{chap:proposed}


%------------------------------------------------------------------------------
\section{Introduction}
\label{sec:ch3_intro}
%------------------------------------------------------------------------------

Chapter~\ref{ch:background} positioned DSSAT as the crop-simulation subsystem
inside a smart-farming System-of-Systems. This chapter turns that view into a
concrete decision framework.

The problem is no longer a single irrigation rule. It is a \emph{multi-objective}
management task in which yield, water use, and nitrogen efficiency must be
balanced over a full growing season. The proposed framework combines three
elements:
\begin{enumerate}
  \item a DSSAT-based daily simulation environment (\texttt{gym-DSSAT});
  \item a hierarchical reward design with preference conditioning for
        multi-objective learning;
  \item learning algorithms---scalar PPO baselines, Preference-Conditioned PPO
        (PC-PPO), and Concave-Augmented Pareto Q-Learning (CAPQL).
\end{enumerate}

The chapter is organised as follows. Section~\ref{sec:sos} formalises the
control problem and the smart-farming System-of-Systems architecture.
Section~\ref{sec:env} describes the simulation platform, observation and action
spaces, and software stack. Section~\ref{sec:reward} presents the hierarchical
reward design and scalar reward variants. Section~\ref{sec:algos} defines the
learning algorithms. Section~\ref{sec:overall} assembles the integrated
architecture and traceability from requirements to implementation. Experimental
training protocols and numerical results are reported in the next chapter.


%------------------------------------------------------------------------------
\section{Problem Formulation and System-of-Systems Architecture}
\label{sec:sos}
%------------------------------------------------------------------------------

\subsection{Overview and Rationale}

A modern smart farm is not a single monolithic system but a \emph{System of
Systems}: a collection of heterogeneous, independently operated subsystems---
weather data, soil and crop observations, irrigation and fertilisation
controllers, and crop-modelling platforms---integrated to pursue missions such
as maximising yield, conserving water, and improving nitrogen-use efficiency.
Following Maier~\cite{maier1998}, a true SoS is distinguished from a merely
complex monolithic system by five traits: operational independence, managerial
independence, geographic distribution, emergent behaviour, and evolutionary
development. The smart farm exhibits all five, which is precisely what makes
mission-level decision-making difficult: an irrigation pulse today changes soil
moisture, crop stress, and nitrogen dynamics over the rest of the season.
Section~\ref{sec:env} formalises this task as a Markov decision process; here
we focus on the architectural view.


\subsection{Components of the Smart-Farming SoS}

The constituent systems are grouped by their role in the closed decision loop:
observe field conditions, act on the crop, and decide daily management. At the
centre of this loop sits the \textbf{crop-modelling platform (DSSAT)}, which
integrates weather, soil, and management inputs and produces crop growth, stress
indicators, and environmental fluxes. Monitoring systems read from it; control
systems write into it; the decision layer issues the commands that drive both.

\vspace{0.25cm}
\noindent\textbf{Monitoring systems.}
These constituents supply the observation stream the agent uses each day:
\begin{itemize}
  \item \textbf{Weather data} --- supplies exogenous drivers (temperature,
        rainfall, solar radiation) that the agent observes but cannot control.
  \item \textbf{Soil and crop state} --- reported through DSSAT variables such
        as soil-water content, water-stress factor, nitrogen-stress factor,
        leaf area, and grain weight.
\end{itemize}

\vspace{0.25cm}
\noindent\textbf{Control systems.}
These constituents execute management actions issued by the decision layer:
\begin{itemize}
  \item \textbf{Irrigation controller} --- the actuator through which the agent
        applies a daily water amount $a_t^{W}$ (mm/day).
  \item \textbf{Fertilisation unit} --- the actuator through which the agent
        applies a daily nitrogen amount $a_t^{N}$ (kg N/ha/day).
\end{itemize}

\vspace{0.25cm}
\noindent\textbf{Decision and mission layer.}
Mission objectives---yield, water efficiency, and nitrogen efficiency---condition
the agent through the preference vector $\boldsymbol{\omega}$. The decision
constituent that closes the loop is:
\begin{itemize}
  \item \textbf{Mission manager / RL agent} --- reads observations, respects the
        farmer's preference vector, and issues daily irrigation and fertilisation
        commands to the control systems.
\end{itemize}

Figure~\ref{fig:sos} depicts the architecture. Monitoring systems feed
observations into the mission manager (solid arrows); the mission manager issues
control commands to the actuators (dashed arrows); DSSAT closes the loop by
producing the next crop and soil state.


%---- FIGURE : SoS architecture ------------------------------------------------
\begin{figure}[ht]
\centering
\includegraphics[width=\textwidth]{sos_architecture.png}
\caption{Architecture of the proposed smart-farming System-of-Systems. Icons denote
constituent roles; solid arrows carry observations and simulated state, dashed
orange arrows carry management commands. DSSAT integrates inputs and closes the
daily decision loop.}
\label{fig:sos}
\end{figure}


\subsection{System-of-Systems Characteristics in Our Setting}

Mapping Maier's five traits~\cite{maier1998} onto our design clarifies why a
learning-based formulation is appropriate. \emph{Operational independence}
holds because weather data, DSSAT, and actuators can each be described on
their own. \emph{Managerial independence} holds because these parts are
developed and maintained separately yet must cooperate in practice.
\emph{Geographic distribution} is intrinsic to field deployment.
\emph{Emergent behaviour}---yield and resource efficiency arising from the
interaction of weather, soil, crop, and management---is exactly what the agent
must learn to shape. \emph{Evolutionary development} motivates a framework in
which the crop model, reward design, or learning algorithm can be upgraded
without rebuilding the entire platform.


\subsection{Mission Objectives and Measures of Effectiveness}
\label{sec:moe}

Measures of Effectiveness (MoEs) translate mission objectives into quantities
that can be logged in simulation and encoded in the reward function.
Table~\ref{tab:moe} links each MoE to its mission objective and to the reward
term that implements it. \emph{Seasonal} MoEs evaluate end-of-season success;
\emph{daily} indicators supply the dense feedback used in reward shaping
(Section~\ref{sec:reward}).

\begin{table}[ht]
\centering
\caption{Measures of Effectiveness and their alignment with mission objectives.}
\label{tab:moe}
\small
\renewcommand{\arraystretch}{1.25}
\begin{tabular}{>{\raggedright\arraybackslash}p{4.4cm}
                >{\raggedright\arraybackslash}p{3.2cm}
                >{\raggedright\arraybackslash}p{4.6cm}}
\toprule
\textbf{Measure of Effectiveness} & \textbf{Mission objective} & \textbf{Indicator / reward link} \\
\midrule
Final grain yield ($\mathrm{GRNWT}_{\text{final}}$) & Yield optimisation & $w_y\,\mathrm{GRNWT}_{\text{final}}$ \\
Harvest index ($\mathrm{HIAD}_{\text{final}}$) & Yield quality & $w_h\,\mathrm{HIAD}_{\text{final}}$ \\
Agronomic N efficiency ($\mathrm{ANE}_{\text{final}}$) & N-use efficiency & $w_a\,\mathrm{ANE}_{\text{final}}$ \\
Total irrigation ($\mathrm{IRRC}_{\text{season}}$) & Water-use efficiency & $w_w\,\mathrm{IRRC}_{\text{season}}$ \\
Total nitrogen applied ($\mathrm{FERN}_{\text{season}}$) & N-use efficiency & $w_f\,\mathrm{FERN}_{\text{season}}$ \\
Soil-moisture balance ($\mathrm{SMR}_t$) & Crop health / water stress & $R_{\text{water},t}$ \\
Water-stress factor ($\mathrm{SWFAC}_t$) & Crop health & $R_{\text{water},t}$ \\
Nitrogen-stress factor ($\mathrm{NSTRES}_t$) & Crop health & $R_{\text{fert},t}$ \\
Cumulative N \& water losses & Environmental impact & $R_{\text{losses},t}$ \\
\bottomrule
\end{tabular}
\end{table}


%------------------------------------------------------------------------------
\section{Simulation Platform and Environment Design}
\label{sec:env}
%------------------------------------------------------------------------------

\subsection{The DSSAT Cropping System Model}

The Decision Support System for Agrotechnology Transfer (DSSAT) is a suite of
mechanistic crop models developed over more than three decades and widely
validated by the agronomic community~\cite{jones2003dssat,hoogenboom2019}. At
its heart is the \emph{Cropping System Model} (CSM), which advances the state of
a land unit on a daily time step. Given weather, soil, genetics, and management
inputs, the CSM couples soil-water balance, soil--plant--atmosphere exchange,
and plant growth modules. For maize, this work uses CERES-Maize. Its outputs
include grain yield, biomass, harvest index, and water and nitrogen fluxes such
as runoff, leaching, and denitrification. Figure~\ref{fig:dssat} sketches the
modular structure.


\begin{figure}[ht]
\centering
\begin{tikzpicture}[node distance=0.7cm and 1.2cm]
  \node[core, fill=sosblue!10, draw=sosblue] (w) {Weather};
  \node[core, fill=sosblue!10, draw=sosblue, below=of w] (s) {Soil profile};
  \node[core, fill=sosgreen!15, draw=sosgreen, right=2.2cm of w, yshift=-0.4cm] (csm)
        {DSSAT-CSM\\daily integrator};
  \node[core, fill=soscream, draw=sosorange, right=2.2cm of csm, yshift=0.9cm] (crop)
        {CERES-Maize};
  \node[core, fill=soscream, draw=sosorange, right=2.2cm of csm, yshift=-0.9cm] (mgmt)
        {Management\\(N + irrigation)};
  \node[block, below=1.6cm of csm, minimum width=4.2cm] (out)
        {Outputs: yield, stress, losses, state vector};

  \draw[flow] (w) -- (csm);
  \draw[flow] (s) -- (csm);
  \draw[flow] (crop) -- (csm);
  \draw[flow] (mgmt) -- (csm);
  \draw[flow] (csm) -- (out);
\end{tikzpicture}
\caption{Schematic of the DSSAT Cropping System Model used in this work.}
\label{fig:dssat}
\end{figure}


\subsection{From DSSAT to an RL Environment: gym-DSSAT}

To use DSSAT for reinforcement learning, the simulator must expose a step-wise
interface in which an agent observes a state, applies an action, and receives
the resulting next state and reward. We rely on \texttt{gym-DSSAT}~\cite{gautron2022gymdssat},
which wraps the Fortran simulator in a Gymnasium-compatible API allowing daily
interaction throughout the growing season~\cite{gautron2022review}. A
\texttt{SmartFarmSoSEnv} layer maps DSSAT outputs to the observation vector,
translates the agent's action into DSSAT management events
(\texttt{anfer}, \texttt{amir}), and computes the reward defined in
Section~\ref{sec:reward}.


\subsection{Decision-Process Formulation}

Formally, the control task is a Markov Decision Process
$\langle \mathcal{S}, \mathcal{A}, P, R, \gamma \rangle$. In the
multi-objective setting, the terminal step also exposes a reward vector
$\mathbf{R}_T$ that is scalarised through a preference vector
$\boldsymbol{\omega}$ (Section~\ref{sec:reward}).

\begin{itemize}
  \item \textbf{State / observation $s_t$:} an 11-dimensional vector of
        normalised crop and SoS features: biomass, grain weight, days after
        planting, vegetative stage, leaf area index, cumulative nitrogen and
        irrigation, mean soil water, energy budget, communication quality, and
        fraction of healthy virtual sensors. PC-PPO and CAPQL append the
        3-dimensional preference vector, giving a 14-dimensional observation.
  \item \textbf{Action $a_t$:} the daily management decision
        $a_t = (\mathrm{FERN}_t, \mathrm{IRRC}_t) = (a_t^{N}, a_t^{W})$ with
        $a_t^{N}\in[0,200]$~kg N/ha/day and $a_t^{W}\in[0,50]$~mm/day.
  \item \textbf{Transition $P(s_{t+1}\mid s_t,a_t)$:} realised by the DSSAT
        daily integration; stochasticity enters through random weather resampling
        across episodes.
  \item \textbf{Horizon:} one growing season (${\approx}161$ days); the episode
        terminates at harvest, when the seasonal reward is granted.
\end{itemize}

Irrigation and fertilisation are coupled because both alter the same soil-water
and nitrogen pools that drive crop stress and final yield. A sequential
formulation is therefore more faithful than optimising either input in isolation.


\subsection{Observation and Action Spaces}

Table~\ref{tab:obs} lists the 11-dimensional observation vector exposed to
scalar PPO agents through the \texttt{pc\_env.py} adapter. Each component is
linearly scaled to $[0,1]$ from agronomically meaningful bounds. PC-PPO and
CAPQL append the three-dimensional preference vector
$\boldsymbol{\omega}=(\omega_y,\omega_a,\omega_w)$, yielding a
14-dimensional input.

\begin{table}[ht]
\centering
\caption{Observation vector for scalar PPO (\texttt{pc\_env.py}).}
\label{tab:obs}
\small
\renewcommand{\arraystretch}{1.2}
\begin{tabular}{clll}
\toprule
\textbf{Idx} & \textbf{DSSAT / SoS variable} & \textbf{Scale} & \textbf{Role} \\
\midrule
0  & \texttt{topwt}  & $[0, 20\,000]$ kg ha$^{-1}$ & Above-ground biomass \\
1  & \texttt{grnwt}  & $[0, 12\,000]$ kg ha$^{-1}$ & Grain weight \\
2  & \texttt{dap}    & $[0, 200]$ days             & Days after planting \\
3  & \texttt{vstage} & $[0, 18]$                   & Vegetative stage \\
4  & \texttt{xlai}   & $[0, 7]$                    & Leaf area index \\
5  & Cumulative N    & $[0, 300]$ kg ha$^{-1}$     & Seasonal fertiliser \\
6  & Cumulative W    & $[0, 1000]$ mm              & Seasonal irrigation \\
7  & Mean soil water & $[0, 1]$                    & Profile moisture \\
8  & Energy budget   & $[0, 100]$                  & SoS energy state \\
9  & Comm.\ quality  & $[0, 1]$                    & Link reliability \\
10 & Sensors alive   & $[0, 1]$                    & Virtual sensor health \\
\bottomrule
\end{tabular}
\end{table}

The daily action vector maps to DSSAT management events:
$a_t=(a_t^{N},a_t^{W})=(\texttt{anfer},\texttt{amir})$ with
$a_t^{N}\in[0,200]$~kg N ha$^{-1}$ day$^{-1}$ and
$a_t^{W}\in[0,50]$~mm day$^{-1}$.


\subsection{Episode Protocol and Calibration Anchors}

Each episode simulates one maize growing season of approximately 161 days.
Weather is resampled randomly across episodes (\texttt{random\_weather=True}),
injecting the main source of transition stochasticity. Yield calibration is
anchored to the stage-based baseline at ${\approx}7\,620$~kg ha$^{-1}$
(\texttt{BASELINE\_YIELD}); normalisation upper bounds follow
\texttt{MAX\_EXPECTED\_YIELD}$=1.5\times$ baseline. Fault injection
(\texttt{enable\_faults=False} in the main experiments) is disabled so that
learning focuses on crop-management trade-offs; resilience extensions are left
for future work.


\subsection{Software Architecture and Module Map}

The implementation stacks four layers: the DSSAT core accessed through
\texttt{gym-DSSAT}; a \texttt{SmartFarmSoSEnv} reward wrapper; a
Gymnasium-compatible adapter; and the training script. Table~\ref{tab:software}
maps each module to its role in the proposed solution.

\begin{table}[ht]
\centering
\caption{Software modules implementing the proposed solution.}
\label{tab:software}
\small
\renewcommand{\arraystretch}{1.25}
\begin{tabular}{>{\raggedright\arraybackslash}p{3.6cm}
                >{\raggedright\arraybackslash}p{4.2cm}
                >{\raggedright\arraybackslash}p{5.5cm}}
\toprule
\textbf{Layer} & \textbf{Module} & \textbf{Function} \\
\midrule
DSSAT core      & \texttt{gym\_dssat\_pdi}           & Daily CSM integration via subprocess API \\
Scalar reward   & \texttt{02\_smart\_farm\_env.py}   & Hierarchical scalar reward (REQ001--002) \\
Water-min reward& \texttt{02\_smart\_farm\_env\_water.py} & Reward weights tuned for irrigation saving \\
N-min reward    & \texttt{02\_smart\_farm\_env\_nitrogen.py} & Reward weights tuned for N saving \\
MORL reward     & \texttt{02\_smart\_farm\_env\_pcppo.py} & Terminal objective vector $\mathbf{R}_T$ \\
11-dim adapter  & \texttt{pc\_env.py}                & Box observation for scalar PPO \\
14-dim adapter  & \texttt{pc\_env\_pcppo.py}         & Appends $\boldsymbol{\omega}$ for PC-PPO \\
CAPQL adapter   & \texttt{capql\_env.py}             & PyTorch interface for TD3 training \\
Scalar PPO train& \texttt{04\_pc\_ppo\_quick\_train.py}, \texttt{04b}, \texttt{04c} & Base, Water-min, N-min \\
PC-PPO train    & \texttt{05\_pcppo\_train.py}       & Preference-conditioned PPO \\
CAPQL train     & \texttt{06\_capql\_train.py}       & Concave-augmented Pareto Q-learning \\
Baselines       & \texttt{01\_rule\_based\_baselines.py} & Random, Fixed, Stage, FAO-56 controllers \\
\bottomrule
\end{tabular}
\end{table}


%------------------------------------------------------------------------------
\section{Reward Engineering}
\label{sec:reward}
%------------------------------------------------------------------------------

The reward design is the heart of the proposed solution: it is where the mission
objectives are encoded into learnable feedback. We adopt a \emph{hierarchical}
structure that combines dense daily rewards with a dominant terminal seasonal
reward. The design is presented below as a set of traceable requirements
(REQ000--REQ002), which keeps an explicit link between each reward term and the
mission objective and MoE it serves.

%=========================================================
\subsection{REQ000: Optimise Irrigation and Fertilisation Performance}
%=========================================================

\noindent\textbf{Detailed Requirement.}
The smart-farming system shall optimise irrigation and nitrogen-fertilisation strategies
throughout the crop growing season in order to maximise crop productivity while
minimising water and fertiliser consumption.

\vspace{0.2cm}
\noindent\textbf{Rationale.}
Efficient agricultural management requires balancing productivity and sustainability.
Excessive irrigation wastes water and may increase runoff and drainage losses, while
excessive fertilisation can increase nutrient leaching and environmental impacts.
Therefore, the reinforcement-learning agent must learn resource-efficient management
policies that maintain healthy crop conditions and maximise final crop yield.

\vspace{0.2cm}
\noindent\textbf{Decomposition into Adaptive Requirements.}
To operationalise this objective, REQ000 is decomposed into the following adaptive
requirements:
\begin{itemize}
    \item \textbf{REQ001:} Optimise Daily Reward ($R_{\text{daily}}$).
    \item \textbf{REQ002:} Optimise Seasonal Reward ($R_{\text{season}}$).
\end{itemize}

\vspace{0.2cm}
\noindent\textbf{Associated Measure of Effectiveness.}
\begin{itemize}
    \item Final grain yield.
    \item Total irrigation applied.
    \item Total nitrogen fertiliser applied.
    \item Agronomic Nitrogen Efficiency (ANE).
    \item Harvest Index (HIAD).
\end{itemize}

\vspace{0.2cm}
\noindent\textbf{Traceability.}
REQ001 and REQ002 derive directly from REQ000 by providing short-term and long-term
feedback mechanisms that guide the agent toward sustainable crop management.

%=========================================================
\subsection{REQ001: Optimise Daily Reward ($R_{\text{daily}}$)}
%=========================================================

\noindent\textbf{Detailed Requirement.}
At each simulation day $t$, the system shall compute an intermediate reward that
encourages healthy crop-growth conditions, efficient resource utilisation, and reduced
environmental losses.

\vspace{0.2cm}
\noindent\textbf{Rationale.}
The final crop yield is only observed at harvest. Relying exclusively on harvest
feedback creates a sparse-reward problem that slows learning. Therefore, daily rewards
provide dense feedback that guides the agent toward agronomically desirable decisions
throughout the season.

\vspace{0.2cm}
\noindent\textbf{Associated Measure of Effectiveness.}
\begin{itemize}
    \item Soil-moisture balance.
    \item Water-stress reduction.
    \item Nitrogen-stress reduction.
    \item Nitrogen-uptake promotion.
    \item Reduced irrigation and fertiliser usage.
    \item Reduced runoff, drainage, and nutrient losses.
\end{itemize}

\vspace{0.2cm}
\noindent\textbf{Operational Definition and Indicators.}
\begin{itemize}
  \item \textbf{Water Management Indicators}
  \begin{itemize}
      \item $\mathrm{SWFAC}_t$: Water-stress factor.
      \item $\mathrm{SMR}_t$: Soil-moisture ratio.
  \end{itemize}
  \begin{equation}
  \mathrm{SMR}_t=\frac{\mathrm{SWTD}_t}{\mathrm{SWXD}_t}
  \end{equation}
  where:
  \begin{itemize}
      \item $\mathrm{SWTD}_t$: Total soil water.
      \item $\mathrm{SWXD}_t$: Available water capacity.
  \end{itemize}

  \item \textbf{Nitrogen Management Indicators}
  \begin{itemize}
      \item $\mathrm{NSTRES}_t$: Nitrogen-stress factor.
      \item $\mathrm{TRNU}_t$: Daily nitrogen uptake.
  \end{itemize}

  \item \textbf{Resource Consumption Indicators}
  \begin{itemize}
      \item $\mathrm{IRRC}_t$: Daily irrigation amount.
      \item $\mathrm{FERN}_t$: Daily nitrogen-fertiliser amount.
  \end{itemize}

  \item \textbf{Environmental Loss Indicators}
  \begin{itemize}
      \item $\mathrm{ROFC}_t$: Surface runoff.
      \item $\mathrm{DRNC}_t$: Deep drainage.
      \item $\mathrm{NLEACH}_t$: Nitrate leaching.
      \item $\mathrm{DENIT}_t$: Denitrification losses.
  \end{itemize}
\end{itemize}

\vspace{0.2cm}
\noindent\textbf{Method of Measurement.}
\begin{enumerate}
  \item \textbf{Water Management Reward ($R_{\text{water}}$)}
  \begin{equation}
  R_{\text{water},t}= w_m R_{\text{moisture},t} + w_s\,(1-\mathrm{SWFAC}_t)
  \end{equation}
  where
 \begin{equation}
R_{\text{moisture},t}=
\begin{cases}
+c_1,
& \text{if } \mathrm{SMR}_t \in [\mathrm{SMR}_{\text{low}},\,\mathrm{SMR}_{\text{high}}]
\\[2mm]
-c_2(\mathrm{SMR}_{\text{low}}-\mathrm{SMR}_t)^2,
& \text{if } \mathrm{SMR}_t < \mathrm{SMR}_{\text{low}}
\\[2mm]
-c_3(\mathrm{SMR}_t-\mathrm{SMR}_{\text{high}})^2,
& \text{if } \mathrm{SMR}_t > \mathrm{SMR}_{\text{high}}
\end{cases}
\end{equation}
  This component encourages maintaining soil moisture within an agronomically desirable
  range while minimising water stress.

  \item \textbf{Fertilisation Reward ($R_{\text{fert}}$)}
  \begin{equation}
  R_{\text{fert},t}= w_n\,(1-\mathrm{NSTRES}_t) + w_u\,\mathrm{TRNU}_t
  \end{equation}
  This reward promotes adequate nitrogen availability and efficient crop nitrogen
  uptake.

  \item \textbf{Resource Consumption Penalty ($R_{\text{resource}}$)}
  \begin{equation}
  R_{\text{resource},t}= w_{\text{irr}}\,\mathrm{IRRC}_t + w_{\text{fert}}\,\mathrm{FERN}_t
  \end{equation}
  This component discourages excessive irrigation and fertilisation while allowing
  necessary applications.

  \item \textbf{Environmental Loss Penalty ($R_{\text{losses}}$)}
  \begin{equation}
  R_{\text{losses},t}= w_r\,\mathrm{ROFC}_t + w_d\,\mathrm{DRNC}_t + w_l\,\mathrm{NLEACH}_t + w_{\text{dn}}\,\mathrm{DENIT}_t
  \end{equation}
  This penalty reduces environmentally harmful practices and resource wastage.
\end{enumerate}

\vspace{0.2cm}
\noindent\textbf{Daily Reward Computation.}
\begin{equation}
R_{\text{daily},t}= R_{\text{water},t} + R_{\text{fert},t} - R_{\text{resource},t} - R_{\text{losses},t}
\label{eq:daily_reward}
\end{equation}
The daily reward evaluates the quality of the management decision taken on day $t$
rather than final crop productivity. In the implemented \texttt{SmartFarmSoSEnv},
each sub-term is normalised to $[-1,1]$ or $[0,1]$ before weighting, and the
combined daily signal is scaled by $\alpha=0.01$ so that seasonal objectives
dominate at harvest.

%=========================================================
\subsection{REQ002: Optimise Seasonal Reward ($R_{\text{season}}$)}
%=========================================================

\noindent\textbf{Detailed Requirement.}
At harvest, the system shall compute a terminal reward that evaluates the overall
success of the growing season based on productivity, efficiency, and resource
utilisation.

\vspace{0.2cm}
\noindent\textbf{Rationale.}
The primary mission objective is to maximise crop yield while using water and fertiliser
efficiently. Therefore, the seasonal reward represents the dominant optimisation
objective.

\vspace{0.2cm}
\noindent\textbf{Associated Measure of Effectiveness.}
\begin{itemize}
    \item Grain yield.
    \item Harvest Index.
    \item Agronomic Nitrogen Efficiency.
    \item Seasonal water consumption.
    \item Seasonal fertiliser consumption.
\end{itemize}

\vspace{0.2cm}
\noindent\textbf{Operational Definition and Indicators.}
\begin{itemize}
  \item $\mathrm{GRNWT}_{\text{final}}$: Final grain yield measured at harvest.
  \item Harvest Index:
  \begin{equation}
  \mathrm{HIAD}_{\text{final}}= \frac{\text{Grain Weight}}{\text{Total Above-Ground Biomass}}
  \end{equation}
  \item Agronomic Nitrogen Efficiency:
  \begin{equation}
  \mathrm{ANE}_{\text{final}}= \frac{Y_N-Y_0}{N_{\text{applied}}}
  \end{equation}
  where:
  \begin{itemize}
      \item $Y_N$: Yield with nitrogen application.
      \item $Y_0$: Yield without nitrogen application.
      \item $N_{\text{applied}}$: Total nitrogen applied.
  \end{itemize}
  \item Seasonal irrigation:
  \begin{equation}
  \mathrm{IRRC}_{\text{season}}= \sum_{t=1}^{T} \mathrm{IRRC}_t
  \end{equation}
  \item Seasonal fertilisation:
  \begin{equation}
  \mathrm{FERN}_{\text{season}}= \sum_{t=1}^{T} \mathrm{FERN}_t
  \end{equation}
\end{itemize}

\vspace{0.2cm}
\noindent\textbf{Method of Measurement.}
The terminal reward is computed at harvest as:
\begin{equation}
R_{\text{season}}= w_y\,\mathrm{GRNWT}_{\text{final}} + w_h\,\mathrm{HIAD}_{\text{final}} + w_a\,\mathrm{ANE}_{\text{final}}
            - w_w\,\mathrm{IRRC}_{\text{season}} - w_f\,\mathrm{FERN}_{\text{season}}
\label{eq:season_reward}
\end{equation}
where:
\begin{itemize}
    \item $w_y$: Yield weight (dominant term).
    \item $w_h$: Harvest-index weight.
    \item $w_a$: Agronomic-nitrogen-efficiency weight.
    \item $w_w$: Water-consumption penalty weight.
    \item $w_f$: Fertiliser-consumption penalty weight.
\end{itemize}

\vspace{0.2cm}
\noindent\textbf{Final Transition Reward.}
For non-terminal days:
\begin{equation}
R_t = R_{\text{daily},t}
\end{equation}
For the harvest day:
\begin{equation}
R_T = R_{\text{daily},T} + R_{\text{season}}
\label{eq:transition_scalar}
\end{equation}
This hierarchical reward structure combines dense daily feedback with long-term seasonal
objectives, enabling efficient learning while preserving the primary goal of maximising
sustainable crop productivity.

%=========================================================
\subsection{From Scalar Rewards to a Multi-Reward Vector}
%=========================================================
The requirements above define the reward \emph{terms}; the multi-reward formulation
groups them into the objective channels that the MORL agents optimise jointly. In this
work we consider four objectives and train accordingly:
\begin{equation}
\mathbf{R}_t = \big(\, R^{\text{water}}_t,\; R^{\text{yield}}_t,\; R^{\text{nitrogen}}_t,\; R^{\text{balanced}}_t \,\big)^{\!\top}
\end{equation}
where $R^{\text{water}}$ aggregates the water-management and irrigation-penalty terms,
$R^{\text{yield}}$ is dominated by the seasonal productivity term, $R^{\text{nitrogen}}$
aggregates the fertilisation reward and nitrogen-loss penalties, and
$R^{\text{balanced}}$ is the scalarised composite that trades the three off against one
another. A scalarised return under a preference vector
$\boldsymbol{\omega}\in\Delta$ (the probability simplex) is then
\begin{equation}
G_{\boldsymbol{\omega}} = \sum_{t} \gamma^{t}\, \boldsymbol{\omega}^{\top}\mathbf{R}_t ,
\qquad \boldsymbol{\omega}\succeq 0,\ \ \textstyle\sum_i \omega_i = 1 .
\label{eq:scalarized}
\end{equation}
Single-objective baselines correspond to placing all preference mass on one channel,
while the MORL agents learn across the whole simplex. In the PC-PPO and CAPQL
implementations, the three primary corners map to yield priority, nitrogen-efficiency
priority, and water-efficiency priority; the balanced channel corresponds to uniform
preference sampling over the simplex.


%=========================================================
\subsection{Scalar Reward Variants for Specialised PPO Baselines}
%=========================================================

Beyond the balanced Base reward (\texttt{02\_smart\_farm\_env.py}), two scalar
variants reweight daily and seasonal terms to probe single-resource priorities
without changing the underlying REQ001--REQ002 structure. These variants train
separate PPO agents (Chapter~4); Table~\ref{tab:reward_variants} records the
parameter differences taken from the implementation files.

\begin{table}[ht]
\centering
\caption{Scalar reward variant parameters (Base, Water-min, N-min).}
\label{tab:reward_variants}
\small
\renewcommand{\arraystretch}{1.15}
\begin{tabular}{lccc}
\toprule
\textbf{Parameter} & \textbf{Base} & \textbf{Water-min} & \textbf{N-min} \\
\midrule
\multicolumn{4}{l}{\textit{Daily blend weights}} \\
\quad $W_{\text{WATER}}$    & 0.40 & 0.20 & 0.35 \\
\quad $W_{\text{FERT}}$     & 0.30 & 0.30 & 0.20 \\
\quad $W_{\text{RESOURCE}}$ & 0.15 & 0.35 & 0.30 \\
\quad $W_{\text{LOSSES}}$   & 0.15 & 0.15 & 0.15 \\
\midrule
\multicolumn{4}{l}{\textit{Seasonal blend weights}} \\
\quad $W_{\text{YIELD}}$    & 0.40 & 0.25 & 0.25 \\
\quad $W_{\text{HIAD}}$     & 0.10 & 0.10 & 0.10 \\
\quad $W_{\text{ANE}}$      & 0.30 & 0.30 & 0.40 \\
\quad $W_{\text{PENALTY}}$  & 0.20 & 0.35 & 0.25 \\
\midrule
\multicolumn{4}{l}{\textit{Seasonal penalty thresholds}} \\
\quad IRRIG\_EXCESS (mm)    & 500  & 150  & 500 \\
\quad N\_EXCESS (kg N ha$^{-1}$) & 200 & 200 & 100 \\
\midrule
\multicolumn{4}{l}{\textit{$R_{\text{penalty}}$ split (water / N)}} \\
\quad $\lambda_W : \lambda_N$ & 50:50 & 85:15 & 30:70 \\
\midrule
N-min yield floor on $R_{\text{ane}}$ & --- & --- & 5000 kg ha$^{-1}$ \\
\bottomrule
\end{tabular}
\end{table}

The Water-min variant reduces moisture-chasing and irrigation-excess tolerance;
the N-min variant strengthens N-application penalties and prioritises agronomic
N efficiency, with a yield floor preventing a zero-fertiliser exploit.


%=========================================================
\subsection{Implementation Constants and Normalisation}
%=========================================================

All sub-reward terms are clipped to $[-1,+1]$ or $[0,1]$ before weighting.
The combined daily signal is scaled by $\alpha=\texttt{DAILY\_SCALE}=0.01$.
Moisture-ratio thresholds are $\mathrm{SMR}_{\text{low}}=0.40$ and
$\mathrm{SMR}_{\text{high}}=0.80$. Seasonal yield, harvest index, and ANE are
normalised against \texttt{BASELINE\_YIELD}$=7620$,
\texttt{BASELINE\_HIAD}$=0.45$, and \texttt{BASELINE\_ANE}$=40$~kg grain per
kg~N. In the implemented Base seasonal reward, the positive terms use weights
$W_{\text{YIELD}}=0.40$, $W_{\text{HIAD}}=0.10$, $W_{\text{ANE}}=0.30$ and the
resource penalty uses $W_{\text{PENALTY}}=0.20$, grouped as
$R_{\text{season}} = W_{\text{YIELD}} R_{\text{yield}} + W_{\text{HIAD}} R_{\text{hiad}}
+ W_{\text{ANE}} R_{\text{ane}} - W_{\text{PENALTY}} R_{\text{penalty}}$.


%------------------------------------------------------------------------------
\section{Multi-Objective Learning Methodology}
\label{sec:algos}
%------------------------------------------------------------------------------

Three learner families implement the proposed solution: scalar PPO baselines
(fixed reward shaping), Preference-Conditioned PPO (PC-PPO), and
Concave-Augmented Pareto Q-Learning (CAPQL). Rule-based controllers provide
non-learning reference policies; their evaluation is deferred to Chapter~4.
Training hyperparameters (rollout length, learning rate, \texttt{VecNormalize})
are likewise reported in Chapter~4.

\subsection{Scalar PPO Baselines}

PPO~\cite{schulman2017ppo} maximises the clipped surrogate
\begin{equation}
\mathcal{L}^{\text{CLIP}}(\theta) =
\mathbb{E}_t\!\Big[\, \min\big( \rho_t(\theta)\,\hat{A}_t,\;
\operatorname{clip}(\rho_t(\theta),\,1-\epsilon,\,1+\epsilon)\,\hat{A}_t \big) \Big],
\label{eq:ppoclip}
\end{equation}
where $\hat{A}_t$ is the generalised advantage estimate~\cite{schulman2016gae}
and $\rho_t(\theta)=\pi_\theta(a_t\mid s_t)/\pi_{\theta_{\text{old}}}(a_t\mid s_t)$.
The scalar PPO agent uses the hierarchical reward of
REQ001--REQ002 (Equations~\eqref{eq:daily_reward}--\eqref{eq:transition_scalar}) with an
11-dimensional observation. Additional specialised variants reweight daily
components toward water saving or nitrogen saving (Table~\ref{tab:reward_variants}).
These agents provide rule-free reference policies against which the
multi-objective learners are compared in Chapter~4.
Algorithm~\ref{alg:scalar_ppo} summarises the on-policy training loop.

\begin{algorithm}[ht]
\caption{Scalar PPO training on \texttt{SmartFarmSoSEnv}}
\label{alg:scalar_ppo}
\begin{algorithmic}[1]
\Require Environment $\mathcal{E}$, total steps $N$, rollout length $L$, clip $\epsilon$, GAE $\lambda$, discount $\gamma$
\State Initialise policy $\pi_\theta$ and value network $V_\phi$; $\theta_{\text{old}} \leftarrow \theta$
\While{environment steps $< N$}
  \State Collect rollout $\mathcal{D}=\{(s_t,a_t,r_t)\}_{t=1}^{L}$ with $a_t \sim \pi_{\theta_{\text{old}}}(\cdot \mid s_t)$
  \State Compute returns $\hat{R}_t$ and advantages $\hat{A}_t$ using GAE($\gamma,\lambda$)
  \For{$k = 1$ to \texttt{n\_epochs}}
    \For{each minibatch $\mathcal{B} \subset \mathcal{D}$}
      \State $\rho_t(\theta) \leftarrow \pi_\theta(a_t \mid s_t)\,/\,\pi_{\theta_{\text{old}}}(a_t \mid s_t)$
      \State Update $\theta$ by ascending $\mathcal{L}^{\text{CLIP}}(\theta)$ on $\mathcal{B}$
      \State Update $\phi$ by minimising the value-function loss on $\mathcal{B}$
    \EndFor
  \EndFor
  \State $\theta_{\text{old}} \leftarrow \theta$
\EndWhile
\State \textbf{return} trained policy $\pi_\theta$
\end{algorithmic}
\end{algorithm}


\subsection{Preference-Conditioned PPO (PC-PPO)}
\label{subsec:pcppo}

PC-PPO is \emph{Preference-Conditioned}, not a separate constrained-optimisation
variant: the same PPO clipped objective is used, but the observation includes
$\boldsymbol{\omega}$ and the multi-reward vector is scalarised through
$\boldsymbol{\omega}^{\top}\mathbf{R}_t$ as in Equation~\eqref{eq:scalarized}.
Each episode samples a new preference $\boldsymbol{\omega}\sim\mathrm{Dirichlet}(1,1,1)$,
so one network learns a family of policies over the simplex. At evaluation,
the same model is queried at yield-, N-efficiency-, and water-efficiency-focused
corners $(1,0,0)$, $(0,1,0)$, $(0,0,1)$ and at a balanced profile
$(\tfrac{1}{3},\tfrac{1}{3},\tfrac{1}{3})$, matching the three-dimensional
terminal vector $\mathbf{R}_T=[R_{\text{yield}}, R_{\text{ane}}, R_{\text{water\_eff}}]^\top$
returned by \texttt{02\_smart\_farm\_env\_pcppo.py}.

PC-PPO uses the 14-dimensional observation
$[s_t;\boldsymbol{\omega}]$ with \texttt{VecNormalize} on observations and rewards.
Its main strength is stable on-policy learning on the costly DSSAT simulator while
still covering multiple farmer priorities with a single model.
Algorithm~\ref{alg:pcppo} details the episode-level preference sampling and
PPO update procedure.

\begin{algorithm}[ht]
\caption{Preference-Conditioned PPO (PC-PPO)}
\label{alg:pcppo}
\begin{algorithmic}[1]
\Require Environment $\mathcal{E}$, total steps $N$, rollout length $L$
\State Initialise $\pi_\theta$, $V_\phi$; $\theta_{\text{old}} \leftarrow \theta$
\While{environment steps $< N$}
  \State Sample $\boldsymbol{\omega} \sim \mathrm{Dirichlet}(\mathbf{1})$ for the upcoming episode
  \State Reset $\mathcal{E}$; $t \leftarrow 0$; \textbf{done} $\leftarrow$ \textbf{false}
  \State $\mathcal{D} \leftarrow \emptyset$
  \While{$|\mathcal{D}| < L$ \textbf{and not} \textbf{done}}
    \State $\tilde{s}_t \leftarrow [s_t;\,\boldsymbol{\omega}]$
    \State $a_t \sim \pi_{\theta_{\text{old}}}(\cdot \mid \tilde{s}_t)$
    \State Execute $(s_{t+1}, r_t, \textbf{done}) \leftarrow \mathcal{E}.\textsc{Step}(a_t)$
    \State Store $(\tilde{s}_t, a_t, r_t)$ in $\mathcal{D}$; $t \leftarrow t + 1$
  \EndWhile
  \State Compute GAE advantages $\{\hat{A}_t\}$ and returns $\{\hat{R}_t\}$ on $\mathcal{D}$
  \For{$k = 1$ to \texttt{n\_epochs}}
    \For{each minibatch $\mathcal{B} \subset \mathcal{D}$}
      \State Update $\theta, \phi$ with clipped PPO loss (Equation~\eqref{eq:ppoclip})
    \EndFor
  \EndFor
  \State $\theta_{\text{old}} \leftarrow \theta$
\EndWhile
\State \textbf{return} $\pi_\theta$
\end{algorithmic}
\end{algorithm}


\subsection{Concave-Augmented Pareto Q-Learning (CAPQL)}
\label{subsec:capql}

CAPQL~\cite{lu2023capql} takes an off-policy route based on TD3. A twin
preference-conditioned critic learns vector-valued returns, while the actor
outputs deterministic daily actions. The terminal target uses linear
scalarisation plus a concave augmentation term:
\begin{equation}
  y = R_{\text{daily}} + \mathbb{1}_{t=T}\!\left(
    \boldsymbol{\omega}^{\top}\mathbf{R}_T
    + \lambda \cdot \mathrm{mean}_i\big[\log(R_i + 1 + \varepsilon)\big]
  \right),
  \label{eq:capql_target}
\end{equation}
with $\lambda=0.5$ and $\varepsilon=0.1$ in our implementation. Hindsight
preference relabeling resamples $\boldsymbol{\omega}$ for minibatch
transitions, improving coverage of the objective space without rerunning
DSSAT episodes.

CAPQL is attractive when many preference values must be explored from limited
simulation budget: one network can be queried at any $\boldsymbol{\omega}$ after
training.
Algorithm~\ref{alg:capql} summarises the TD3 interaction loop, replay storage,
hindsight preference relabeling, and concave-augmented critic--actor updates.

\begin{algorithm}[ht]
\caption{Concave-Augmented Pareto Q-Learning (CAPQL)}
\label{alg:capql}
\begin{algorithmic}[1]
\Require Environment $\mathcal{E}$, replay buffer $\mathcal{B}$, total steps $N$, batch size $B$, policy delay $d$
\State Initialise actor $\mu_\theta$, critics $Q_{\phi_1}, Q_{\phi_2}$, and target networks $\mu_{\bar\theta}, Q_{\bar\phi}$
\For{$n = 1$ to $N$}
  \If{episode start}
    \State Sample $\boldsymbol{\omega} \sim \mathrm{Dirichlet}(\mathbf{1})$; reset $\mathcal{E}$
  \EndIf
  \State $\tilde{s}_t \leftarrow [s_t;\,\boldsymbol{\omega}]$
  \State $a_t \leftarrow \mu_\theta(\tilde{s}_t) + \text{exploration noise}$; step $\mathcal{E}$
  \State Store $(s_t, a_t, R_{\text{daily},t}, \mathbf{R}_T\,\mathbb{1}_{t=T}, s_{t+1}, \textbf{done}, \boldsymbol{\omega})$ in $\mathcal{B}$
  \If{$|\mathcal{B}| \ge B$ \textbf{and} $n > \texttt{warmup}$}
    \State Sample minibatch; relabel $\boldsymbol{\omega}' \sim \mathrm{Dirichlet}(\mathbf{1})$ per transition
    \State $\tilde{s} \leftarrow [s;\boldsymbol{\omega}']$,\; $\tilde{s}' \leftarrow [s';\boldsymbol{\omega}']$
    \State $r_{\text{aug}} \leftarrow R_{\text{daily}} + \textbf{done}\cdot\big(\boldsymbol{\omega}'^{\top}\mathbf{R}_T + \lambda\,\mathrm{mean}_i\log(R_i{+}1{+}\varepsilon)\big)$
    \State $a' \leftarrow \mathrm{clip}\big(\mu_{\bar\theta}(\tilde{s}') + \text{smoothing noise}\big)$
    \State $y \leftarrow r_{\text{aug}} + \gamma(1-\textbf{done})\min_{j \in \{1,2\}} Q_{\bar\phi_j}(\tilde{s}', a')$
    \State Update $\phi_1, \phi_2$ by minimising $\|Q_{\phi_j}(\tilde{s},a) - y\|^2$
    \If{$n \bmod d = 0$}
      \State Update $\theta$ by ascending $\mathbb{E}\big[\min_j Q_{\phi_j}(\tilde{s}, \mu_\theta(\tilde{s}))\big]$
      \State Soft-update $\bar\theta \leftarrow \tau\theta + (1{-}\tau)\bar\theta$;\;
            $\bar\phi \leftarrow \tau\phi + (1{-}\tau)\bar\phi$
    \EndIf
  \EndIf
\EndFor
\State \textbf{return} $\mu_\theta$
\end{algorithmic}
\end{algorithm}


\subsection{Justification of Algorithmic Choices}

\begin{table}[ht]
\centering
\caption{Rationale for the selected learning algorithms.}
\label{tab:algos}
\small
\renewcommand{\arraystretch}{1.3}
\begin{tabular}{>{\raggedright\arraybackslash}p{2.8cm}
                >{\raggedright\arraybackslash}p{2.6cm}
                >{\raggedright\arraybackslash}p{6.3cm}}
\toprule
\textbf{Method} & \textbf{Family} & \textbf{Why chosen here} \\
\midrule
Scalar PPO & On-policy, single-objective &
Stable baseline on \texttt{gym-DSSAT}; hierarchical reward matches REQ000--REQ002;
provides fixed-reward reference policies for comparison. \\
PC-PPO & On-policy, preference-conditioned &
Same PPO core as the baseline; one model serves all preference corners;
observation carries $\boldsymbol{\omega}$; strong training stability. \\
CAPQL & Off-policy, preference-conditioned &
Sample-efficient exploration of the Pareto set; concave augmentation mitigates
plain linear scalarisation; TD3 twin critics reduce Q overestimation. \\
\bottomrule
\end{tabular}
\end{table}

In short, scalar PPO defines what fixed hierarchical rewards can encode
(Algorithm~\ref{alg:scalar_ppo}); PC-PPO extends the same PPO core to a
preference-conditioned policy family (Algorithm~\ref{alg:pcppo}); CAPQL
complements this with an off-policy Pareto-learning formulation
(Algorithm~\ref{alg:capql}).


%------------------------------------------------------------------------------
\section{Integrated Proposed Architecture}
\label{sec:overall}
%------------------------------------------------------------------------------

Figure~\ref{fig:overall} assembles the elements of this chapter into the complete
closed-loop learning system. The DSSAT core, wrapped by \texttt{gym-DSSAT} and
\texttt{SmartFarmSoSEnv}, exposes the daily crop state. The reward module
converts DSSAT outputs into daily and seasonal signals (Section~\ref{sec:reward}).
The agent---scalar PPO, PC-PPO, or CAPQL---consumes state and reward and emits
irrigation and nitrogen actions, closing the loop over one maize season.

%---- FIGURE : Overall proposed model ----------------------------------------
\begin{figure}[ht]
\centering
\resizebox{\textwidth}{!}{%
\begin{tikzpicture}[node distance=1.0cm and 1.8cm]

  \node[core, draw=sosblue, fill=sosblue!10, minimum width=3.4cm, minimum height=1.6cm]
       (env) {gym-DSSAT\\Environment\\\scriptsize (DSSAT CSM core)};

  \node[core, draw=sosgreen, fill=sosgreen!10, below=2.2cm of env,
        minimum width=3.4cm, minimum height=1.6cm] (reward)
        {Reward Module\\\scriptsize $R_{\text{daily}}+R_{\text{season}}$};

  \node[core, draw=sosorange, fill=soscream, right=4.2cm of env,
        minimum width=3.6cm, minimum height=1.6cm] (agent)
        {Multi-Objective RL Agent\\\scriptsize PPO $\,\cdot\,$ PC-PPO $\,\cdot\,$ CAPQL};

  \node[block, draw=sosgreen, fill=sosgreen!8, below=2.2cm of agent,
        minimum width=3.6cm, minimum height=1.5cm, font=\scriptsize, align=center]
        (vec) {Seasonal vector $\mathbf{R}_T$\\yield $\cdot$ N-efficiency\\water efficiency};

  \draw[flow] (env.east) -- node[above,font=\scriptsize]{state $s_t$} (agent.west);
  \draw[ctrl] (agent.south) |- node[below right,font=\scriptsize,pos=0.25]
        {action $a_t=(\mathrm{FERN}_t,\mathrm{IRRC}_t)$} (env.south);
  \draw[flow] (env.south) -- node[left,font=\scriptsize]{DSSAT outputs} (reward.north);
  \draw[flow] (reward.east) -- (vec.west);
  \draw[flow] (vec.north) -- node[right,font=\scriptsize]{$\mathbf{R}_T$} (agent.south);

  \node[block, draw=sosgrey, fill=white, right=1.2cm of agent, font=\scriptsize,
        align=center, minimum width=2.6cm] (eval)
        {Evaluation\\\scriptsize vs.\ rule-based\\and scalar PPO baselines};
  \draw[flow, sosgrey] (agent.east) -- (eval.west);

  \begin{scope}[on background layer]
    \node[draw=sosgrey, dashed, rounded corners, inner sep=8pt,
          fit=(env)(reward)(agent)(vec),
          label={[sosgrey,font=\scriptsize\bfseries]above:Closed RL loop (one growing season)}] {};
  \end{scope}

\end{tikzpicture}%
}
\caption{Overall proposed learning architecture. The DSSAT core supplies daily
  crop state; the reward module computes hierarchical daily and seasonal
  signals; the agent issues irrigation and nitrogen actions, closing the loop
  over one maize season.}
\label{fig:overall}
\end{figure}


\subsection{Traceability from Requirements to Implementation}

Table~\ref{tab:trace} summarises how REQ000--REQ002 connect mission objectives,
MoEs, reward terms, and learner families. This traceability matrix is the
design counterpart to the performance tables reported in Chapter~4.

\begin{table}[ht]
\centering
\caption{Traceability: requirements, MoEs, rewards, and learners.}
\label{tab:trace}
\small
\renewcommand{\arraystretch}{1.2}
\begin{tabular}{>{\raggedright\arraybackslash}p{1.6cm}
                >{\raggedright\arraybackslash}p{2.8cm}
                >{\raggedright\arraybackslash}p{2.6cm}
                >{\raggedright\arraybackslash}p{2.4cm}
                >{\raggedright\arraybackslash}p{2.2cm}}
\toprule
\textbf{Req.} & \textbf{MoE / objective} & \textbf{Reward term} & \textbf{Module} & \textbf{Learner} \\
\midrule
REQ000 & Yield, N, water & $R_{\text{daily}}+R_{\text{season}}$ &
\texttt{02\_smart\_farm\_env*} & PPO, PC-PPO, CAPQL \\
REQ001 & SMR, SWFAC, NSTRES & $R_{\text{water}}, R_{\text{fert}}$ &
\texttt{02\_smart\_farm\_env.py} & All \\
REQ001 & Resource use, losses & $R_{\text{resource}}, R_{\text{losses}}$ &
\texttt{02\_smart\_farm\_env.py} & All \\
REQ002 & GRNWT, HIAD, ANE & $R_{\text{yield}}, R_{\text{hiad}}, R_{\text{ane}}$ &
\texttt{02\_smart\_farm\_env.py} & PPO \\
REQ002 & Season totals & $R_{\text{penalty}}$ &
\texttt{02\_smart\_farm\_env*.py} & PPO variants \\
MORL   & Preference trade-offs & $\mathbf{R}_T$, $\boldsymbol{\omega}$ &
\texttt{02\_smart\_farm\_env\_pcppo.py} & PC-PPO, CAPQL \\
\bottomrule
\end{tabular}
\end{table}


\subsection{Comparison Baselines (Design)}

For evaluation (Chapter~4), four rule-based controllers run inside the same
\texttt{SmartFarmSoSEnv} wrapper: \emph{Random} (unconstrained daily inputs),
\emph{Fixed} (preset N splits and periodic irrigation), \emph{Stage-based}
(management keyed to vegetative development), and \emph{FAO-56} (soil-water-balance
irrigation with a three-split N plan). Using a common environment ensures that
MoEs and reward components are directly comparable across learned and rule-based
policies.


%------------------------------------------------------------------------------
\section{Conclusion}
\label{sec:ch3_conclusion}
%------------------------------------------------------------------------------

This chapter presented the proposed solution end to end: a smart-farming
System-of-Systems formalised around the DSSAT Cropping System Model; an MDP
formulation with explicit observation, action, and software layers;
Measures of Effectiveness tied to mission objectives (Table~\ref{tab:moe});
a hierarchical reward design (REQ000--REQ002) with scalar variants
(Table~\ref{tab:reward_variants}) and a multi-objective extension for
preference-conditioned learning; and three learner families---scalar PPO,
PC-PPO, and CAPQL---integrated in the closed-loop architecture of
Figure~\ref{fig:overall}. The next chapter reports the experimental setup,
training dynamics, and quantitative comparison of all learned and rule-based
policies.


% =============================================================================
% If using BibTeX instead, delete this block and add entries to references.bib
% =============================================================================
\begin{thebibliography}{99}

\bibitem{maier1998}
M.~W. Maier,
``Architecting principles for systems-of-systems,''
\emph{Systems Engineering}, vol.~1, no.~4, pp.~267--284, 1998.

\bibitem{jones2003dssat}
J.~W. Jones, G.~Hoogenboom, C.~H. Porter, K.~J. Boote, W.~D. Batchelor, L.~A. Hunt,
P.~W. Wilkens, U.~Singh, A.~J. Gijsman, and J.~T. Ritchie,
``The DSSAT cropping system model,''
\emph{European Journal of Agronomy}, vol.~18, no.~3--4, pp.~235--265, 2003.

\bibitem{hoogenboom2019}
G.~Hoogenboom, C.~H. Porter, K.~J. Boote, V.~Shelia, P.~W. Wilkens, U.~Singh,
J.~W. White, S.~Asseng, J.~I. Lizaso, et~al.,
``The DSSAT crop modeling ecosystem,'' in
\emph{Advances in Crop Modelling for a Sustainable Agriculture}, Burleigh Dodds Science
Publishing, pp.~173--216, 2019.

\bibitem{gautron2022gymdssat}
R.~Gautron, E.~J. Padr\'on, P.~Preux, J.~Bigot, O.-A. Maillard, and D.~Emukpere,
``gym-DSSAT: a crop model turned into a Reinforcement Learning environment,''
Inria Research Report RR-9460, 2022.

\bibitem{gautron2022review}
R.~Gautron, O.-A. Maillard, P.~Preux, M.~Corbeels, and R.~Sabbadin,
``Reinforcement learning for crop management support: Review, prospects and challenges,''
\emph{Computers and Electronics in Agriculture}, vol.~200, art.~107182, 2022.

\bibitem{schulman2017ppo}
J.~Schulman, F.~Wolski, P.~Dhariwal, A.~Radford, and O.~Klimov,
``Proximal policy optimization algorithms,''
\texttt{arXiv:1707.06347}, 2017.

\bibitem{schulman2016gae}
J.~Schulman, P.~Moritz, S.~Levine, M.~Jordan, and P.~Abbeel,
``High-dimensional continuous control using generalized advantage estimation,'' in
\emph{Proc.\ ICLR}, 2016.

\bibitem{lu2023capql}
H.~Lu, D.~Herman, and Y.~Yu,
``Multi-objective reinforcement learning: Convexity, stationarity and Pareto optimality,''
in \emph{Proc.\ ICLR}, 2023.

\end{thebibliography}

% Required: amsmath, booktabs, graphicx, tikz, backgrounds, algorithm, algpseudocode
%   \usetikzlibrary{positioning, arrows.meta, fit, calc, backgrounds}
% =============================================================================

%---- Colour palette for the diagrams ----------------------------------------
\definecolor{sosblue}{RGB}{38,84,124}
\definecolor{sosgreen}{RGB}{61,133,92}
\definecolor{sosorange}{RGB}{200,120,40}
\definecolor{sosgrey}{RGB}{90,90,90}
\definecolor{soslight}{RGB}{234,240,245}
\definecolor{soscream}{RGB}{248,244,232}

%---- Reusable TikZ node styles ----------------------------------------------
\tikzset{
  sensor/.style   = {rectangle, rounded corners=2pt, draw=sosblue, very thick,
                     fill=soslight, align=center, minimum width=2.5cm,
                     minimum height=1cm, font=\small},
  actuator/.style = {rectangle, rounded corners=2pt, draw=sosorange, very thick,
                     fill=soscream, align=center, minimum width=2.5cm,
                     minimum height=1cm, font=\small},
  core/.style     = {rectangle, rounded corners=4pt, draw=sosgreen, line width=1.2pt,
                     fill=sosgreen!12, align=center, minimum width=3.2cm,
                     minimum height=1.3cm, font=\small\bfseries},
  mission/.style  = {ellipse, draw=sosgrey, very thick, fill=white, align=center,
                     font=\small\itshape, minimum width=3cm},
  flow/.style     = {-{Stealth[length=2.5mm]}, line width=0.9pt, sosgrey},
  ctrl/.style     = {-{Stealth[length=2.5mm]}, line width=0.9pt, sosorange, dashed},
  block/.style    = {rectangle, rounded corners=2pt, draw=sosblue, thick,
                     fill=soslight, align=center, font=\small,
                     minimum width=2.6cm, minimum height=0.9cm},
}


\chapter{Proposed Solution}
\label{chap:proposed}


%------------------------------------------------------------------------------
\section{Introduction}
\label{sec:ch3_intro}
%------------------------------------------------------------------------------

Chapter~\ref{ch:background} positioned DSSAT as the crop-simulation subsystem
inside a smart-farming System-of-Systems. This chapter turns that view into a
concrete decision framework.

The problem is no longer a single irrigation rule. It is a \emph{multi-objective}
management task in which yield, water use, and nitrogen efficiency must be
balanced over a full growing season. The proposed framework combines three
elements:
\begin{enumerate}
  \item a DSSAT-based daily simulation environment (\texttt{gym-DSSAT});
  \item a hierarchical reward design with preference conditioning for
        multi-objective learning;
  \item learning algorithms---scalar PPO baselines, Preference-Conditioned PPO
        (PC-PPO), and Concave-Augmented Pareto Q-Learning (CAPQL).
\end{enumerate}

The chapter is organised as follows. Section~\ref{sec:sos} formalises the
control problem and the smart-farming System-of-Systems architecture.
Section~\ref{sec:env} describes the simulation platform, observation and action
spaces, and software stack. Section~\ref{sec:reward} presents the hierarchical
reward design and scalar reward variants. Section~\ref{sec:algos} defines the
learning algorithms. Section~\ref{sec:overall} assembles the integrated
architecture and traceability from requirements to implementation. Experimental
training protocols and numerical results are reported in the next chapter.


%------------------------------------------------------------------------------
\section{Problem Formulation and System-of-Systems Architecture}
\label{sec:sos}
%------------------------------------------------------------------------------

\subsection{Overview and Rationale}

A modern smart farm is not a single monolithic system but a \emph{System of
Systems}: a collection of heterogeneous, independently operated subsystems---
weather data, soil and crop observations, irrigation and fertilisation
controllers, and crop-modelling platforms---integrated to pursue missions such
as maximising yield, conserving water, and improving nitrogen-use efficiency.
Following Maier~\cite{maier1998}, a true SoS is distinguished from a merely
complex monolithic system by five traits: operational independence, managerial
independence, geographic distribution, emergent behaviour, and evolutionary
development. The smart farm exhibits all five, which is precisely what makes
mission-level decision-making difficult: an irrigation pulse today changes soil
moisture, crop stress, and nitrogen dynamics over the rest of the season.
Section~\ref{sec:env} formalises this task as a Markov decision process; here
we focus on the architectural view.


\subsection{Components of the Smart-Farming SoS}

The constituent systems are grouped by their role in the closed decision loop:
observe field conditions, act on the crop, and decide daily management. At the
centre of this loop sits the \textbf{crop-modelling platform (DSSAT)}, which
integrates weather, soil, and management inputs and produces crop growth, stress
indicators, and environmental fluxes. Monitoring systems read from it; control
systems write into it; the decision layer issues the commands that drive both.

\vspace{0.25cm}
\noindent\textbf{Monitoring systems.}
These constituents supply the observation stream the agent uses each day:
\begin{itemize}
  \item \textbf{Weather data} --- supplies exogenous drivers (temperature,
        rainfall, solar radiation) that the agent observes but cannot control.
  \item \textbf{Soil and crop state} --- reported through DSSAT variables such
        as soil-water content, water-stress factor, nitrogen-stress factor,
        leaf area, and grain weight.
\end{itemize}

\vspace{0.25cm}
\noindent\textbf{Control systems.}
These constituents execute management actions issued by the decision layer:
\begin{itemize}
  \item \textbf{Irrigation controller} --- the actuator through which the agent
        applies a daily water amount $a_t^{W}$ (mm/day).
  \item \textbf{Fertilisation unit} --- the actuator through which the agent
        applies a daily nitrogen amount $a_t^{N}$ (kg N/ha/day).
\end{itemize}

\vspace{0.25cm}
\noindent\textbf{Decision and mission layer.}
Mission objectives---yield, water efficiency, and nitrogen efficiency---condition
the agent through the preference vector $\boldsymbol{\omega}$. The decision
constituent that closes the loop is:
\begin{itemize}
  \item \textbf{Mission manager / RL agent} --- reads observations, respects the
        farmer's preference vector, and issues daily irrigation and fertilisation
        commands to the control systems.
\end{itemize}

Figure~\ref{fig:sos} depicts the architecture. Monitoring systems feed
observations into the mission manager (solid arrows); the mission manager issues
control commands to the actuators (dashed arrows); DSSAT closes the loop by
producing the next crop and soil state.


%---- FIGURE : SoS architecture ------------------------------------------------
\begin{figure}[ht]
\centering
\includegraphics[width=\textwidth]{sos_architecture.png}
\caption{Architecture of the proposed smart-farming System-of-Systems. Icons denote
constituent roles; solid arrows carry observations and simulated state, dashed
orange arrows carry management commands. DSSAT integrates inputs and closes the
daily decision loop.}
\label{fig:sos}
\end{figure}


\subsection{System-of-Systems Characteristics in Our Setting}

Mapping Maier's five traits~\cite{maier1998} onto our design clarifies why a
learning-based formulation is appropriate. \emph{Operational independence}
holds because weather data, DSSAT, and actuators can each be described on
their own. \emph{Managerial independence} holds because these parts are
developed and maintained separately yet must cooperate in practice.
\emph{Geographic distribution} is intrinsic to field deployment.
\emph{Emergent behaviour}---yield and resource efficiency arising from the
interaction of weather, soil, crop, and management---is exactly what the agent
must learn to shape. \emph{Evolutionary development} motivates a framework in
which the crop model, reward design, or learning algorithm can be upgraded
without rebuilding the entire platform.


\subsection{Mission Objectives and Measures of Effectiveness}
\label{sec:moe}

Measures of Effectiveness (MoEs) translate mission objectives into quantities
that can be logged in simulation and encoded in the reward function.
Table~\ref{tab:moe} links each MoE to its mission objective and to the reward
term that implements it. \emph{Seasonal} MoEs evaluate end-of-season success;
\emph{daily} indicators supply the dense feedback used in reward shaping
(Section~\ref{sec:reward}).

\begin{table}[ht]
\centering
\caption{Measures of Effectiveness and their alignment with mission objectives.}
\label{tab:moe}
\small
\renewcommand{\arraystretch}{1.25}
\begin{tabular}{>{\raggedright\arraybackslash}p{4.4cm}
                >{\raggedright\arraybackslash}p{3.2cm}
                >{\raggedright\arraybackslash}p{4.6cm}}
\toprule
\textbf{Measure of Effectiveness} & \textbf{Mission objective} & \textbf{Indicator / reward link} \\
\midrule
Final grain yield ($\mathrm{GRNWT}_{\text{final}}$) & Yield optimisation & $w_y\,\mathrm{GRNWT}_{\text{final}}$ \\
Harvest index ($\mathrm{HIAD}_{\text{final}}$) & Yield quality & $w_h\,\mathrm{HIAD}_{\text{final}}$ \\
Agronomic N efficiency ($\mathrm{ANE}_{\text{final}}$) & N-use efficiency & $w_a\,\mathrm{ANE}_{\text{final}}$ \\
Total irrigation ($\mathrm{IRRC}_{\text{season}}$) & Water-use efficiency & $w_w\,\mathrm{IRRC}_{\text{season}}$ \\
Total nitrogen applied ($\mathrm{FERN}_{\text{season}}$) & N-use efficiency & $w_f\,\mathrm{FERN}_{\text{season}}$ \\
Soil-moisture balance ($\mathrm{SMR}_t$) & Crop health / water stress & $R_{\text{water},t}$ \\
Water-stress factor ($\mathrm{SWFAC}_t$) & Crop health & $R_{\text{water},t}$ \\
Nitrogen-stress factor ($\mathrm{NSTRES}_t$) & Crop health & $R_{\text{fert},t}$ \\
Cumulative N \& water losses & Environmental impact & $R_{\text{losses},t}$ \\
\bottomrule
\end{tabular}
\end{table}


%------------------------------------------------------------------------------
\section{Simulation Platform and Environment Design}
\label{sec:env}
%------------------------------------------------------------------------------

\subsection{The DSSAT Cropping System Model}

The Decision Support System for Agrotechnology Transfer (DSSAT) is a suite of
mechanistic crop models developed over more than three decades and widely
validated by the agronomic community~\cite{jones2003dssat,hoogenboom2019}. At
its heart is the \emph{Cropping System Model} (CSM), which advances the state of
a land unit on a daily time step. Given weather, soil, genetics, and management
inputs, the CSM couples soil-water balance, soil--plant--atmosphere exchange,
and plant growth modules. For maize, this work uses CERES-Maize. Its outputs
include grain yield, biomass, harvest index, and water and nitrogen fluxes such
as runoff, leaching, and denitrification. Figure~\ref{fig:dssat} sketches the
modular structure.


\begin{figure}[ht]
\centering
\begin{tikzpicture}[node distance=0.7cm and 1.2cm]
  \node[core, fill=sosblue!10, draw=sosblue] (w) {Weather};
  \node[core, fill=sosblue!10, draw=sosblue, below=of w] (s) {Soil profile};
  \node[core, fill=sosgreen!15, draw=sosgreen, right=2.2cm of w, yshift=-0.4cm] (csm)
        {DSSAT-CSM\\daily integrator};
  \node[core, fill=soscream, draw=sosorange, right=2.2cm of csm, yshift=0.9cm] (crop)
        {CERES-Maize};
  \node[core, fill=soscream, draw=sosorange, right=2.2cm of csm, yshift=-0.9cm] (mgmt)
        {Management\\(N + irrigation)};
  \node[block, below=1.6cm of csm, minimum width=4.2cm] (out)
        {Outputs: yield, stress, losses, state vector};

  \draw[flow] (w) -- (csm);
  \draw[flow] (s) -- (csm);
  \draw[flow] (crop) -- (csm);
  \draw[flow] (mgmt) -- (csm);
  \draw[flow] (csm) -- (out);
\end{tikzpicture}
\caption{Schematic of the DSSAT Cropping System Model used in this work.}
\label{fig:dssat}
\end{figure}


\subsection{From DSSAT to an RL Environment: gym-DSSAT}

To use DSSAT for reinforcement learning, the simulator must expose a step-wise
interface in which an agent observes a state, applies an action, and receives
the resulting next state and reward. We rely on \texttt{gym-DSSAT}~\cite{gautron2022gymdssat},
which wraps the Fortran simulator in a Gymnasium-compatible API allowing daily
interaction throughout the growing season~\cite{gautron2022review}. A
\texttt{SmartFarmSoSEnv} layer maps DSSAT outputs to the observation vector,
translates the agent's action into DSSAT management events
(\texttt{anfer}, \texttt{amir}), and computes the reward defined in
Section~\ref{sec:reward}.


\subsection{Decision-Process Formulation}

Formally, the control task is a Markov Decision Process
$\langle \mathcal{S}, \mathcal{A}, P, R, \gamma \rangle$. In the
multi-objective setting, the terminal step also exposes a reward vector
$\mathbf{R}_T$ that is scalarised through a preference vector
$\boldsymbol{\omega}$ (Section~\ref{sec:reward}).

\begin{itemize}
  \item \textbf{State / observation $s_t$:} an 11-dimensional vector of
        normalised crop and SoS features: biomass, grain weight, days after
        planting, vegetative stage, leaf area index, cumulative nitrogen and
        irrigation, mean soil water, energy budget, communication quality, and
        fraction of healthy virtual sensors. PC-PPO and CAPQL append the
        3-dimensional preference vector, giving a 14-dimensional observation.
  \item \textbf{Action $a_t$:} the daily management decision
        $a_t = (\mathrm{FERN}_t, \mathrm{IRRC}_t) = (a_t^{N}, a_t^{W})$ with
        $a_t^{N}\in[0,200]$~kg N/ha/day and $a_t^{W}\in[0,50]$~mm/day.
  \item \textbf{Transition $P(s_{t+1}\mid s_t,a_t)$:} realised by the DSSAT
        daily integration; stochasticity enters through random weather resampling
        across episodes.
  \item \textbf{Horizon:} one growing season (${\approx}161$ days); the episode
        terminates at harvest, when the seasonal reward is granted.
\end{itemize}

Irrigation and fertilisation are coupled because both alter the same soil-water
and nitrogen pools that drive crop stress and final yield. A sequential
formulation is therefore more faithful than optimising either input in isolation.


\subsection{Observation and Action Spaces}

Table~\ref{tab:obs} lists the 11-dimensional observation vector exposed to
scalar PPO agents through the \texttt{pc\_env.py} adapter. Each component is
linearly scaled to $[0,1]$ from agronomically meaningful bounds. PC-PPO and
CAPQL append the three-dimensional preference vector
$\boldsymbol{\omega}=(\omega_y,\omega_a,\omega_w)$, yielding a
14-dimensional input.

\begin{table}[ht]
\centering
\caption{Observation vector for scalar PPO (\texttt{pc\_env.py}).}
\label{tab:obs}
\small
\renewcommand{\arraystretch}{1.2}
\begin{tabular}{clll}
\toprule
\textbf{Idx} & \textbf{DSSAT / SoS variable} & \textbf{Scale} & \textbf{Role} \\
\midrule
0  & \texttt{topwt}  & $[0, 20\,000]$ kg ha$^{-1}$ & Above-ground biomass \\
1  & \texttt{grnwt}  & $[0, 12\,000]$ kg ha$^{-1}$ & Grain weight \\
2  & \texttt{dap}    & $[0, 200]$ days             & Days after planting \\
3  & \texttt{vstage} & $[0, 18]$                   & Vegetative stage \\
4  & \texttt{xlai}   & $[0, 7]$                    & Leaf area index \\
5  & Cumulative N    & $[0, 300]$ kg ha$^{-1}$     & Seasonal fertiliser \\
6  & Cumulative W    & $[0, 1000]$ mm              & Seasonal irrigation \\
7  & Mean soil water & $[0, 1]$                    & Profile moisture \\
8  & Energy budget   & $[0, 100]$                  & SoS energy state \\
9  & Comm.\ quality  & $[0, 1]$                    & Link reliability \\
10 & Sensors alive   & $[0, 1]$                    & Virtual sensor health \\
\bottomrule
\end{tabular}
\end{table}

The daily action vector maps to DSSAT management events:
$a_t=(a_t^{N},a_t^{W})=(\texttt{anfer},\texttt{amir})$ with
$a_t^{N}\in[0,200]$~kg N ha$^{-1}$ day$^{-1}$ and
$a_t^{W}\in[0,50]$~mm day$^{-1}$.


\subsection{Episode Protocol and Calibration Anchors}

Each episode simulates one maize growing season of approximately 161 days.
Weather is resampled randomly across episodes (\texttt{random\_weather=True}),
injecting the main source of transition stochasticity. Yield calibration is
anchored to the stage-based baseline at ${\approx}7\,620$~kg ha$^{-1}$
(\texttt{BASELINE\_YIELD}); normalisation upper bounds follow
\texttt{MAX\_EXPECTED\_YIELD}$=1.5\times$ baseline. Fault injection
(\texttt{enable\_faults=False} in the main experiments) is disabled so that
learning focuses on crop-management trade-offs; resilience extensions are left
for future work.


\subsection{Software Architecture and Module Map}

The implementation stacks four layers: the DSSAT core accessed through
\texttt{gym-DSSAT}; a \texttt{SmartFarmSoSEnv} reward wrapper; a
Gymnasium-compatible adapter; and the training script. Table~\ref{tab:software}
maps each module to its role in the proposed solution.

\begin{table}[ht]
\centering
\caption{Software modules implementing the proposed solution.}
\label{tab:software}
\small
\renewcommand{\arraystretch}{1.25}
\begin{tabular}{>{\raggedright\arraybackslash}p{3.6cm}
                >{\raggedright\arraybackslash}p{4.2cm}
                >{\raggedright\arraybackslash}p{5.5cm}}
\toprule
\textbf{Layer} & \textbf{Module} & \textbf{Function} \\
\midrule
DSSAT core      & \texttt{gym\_dssat\_pdi}           & Daily CSM integration via subprocess API \\
Scalar reward   & \texttt{02\_smart\_farm\_env.py}   & Hierarchical scalar reward (REQ001--002) \\
Water-min reward& \texttt{02\_smart\_farm\_env\_water.py} & Reward weights tuned for irrigation saving \\
N-min reward    & \texttt{02\_smart\_farm\_env\_nitrogen.py} & Reward weights tuned for N saving \\
MORL reward     & \texttt{02\_smart\_farm\_env\_pcppo.py} & Terminal objective vector $\mathbf{R}_T$ \\
11-dim adapter  & \texttt{pc\_env.py}                & Box observation for scalar PPO \\
14-dim adapter  & \texttt{pc\_env\_pcppo.py}         & Appends $\boldsymbol{\omega}$ for PC-PPO \\
CAPQL adapter   & \texttt{capql\_env.py}             & PyTorch interface for TD3 training \\
Scalar PPO train& \texttt{04\_pc\_ppo\_quick\_train.py}, \texttt{04b}, \texttt{04c} & Base, Water-min, N-min \\
PC-PPO train    & \texttt{05\_pcppo\_train.py}       & Preference-conditioned PPO \\
CAPQL train     & \texttt{06\_capql\_train.py}       & Concave-augmented Pareto Q-learning \\
Baselines       & \texttt{01\_rule\_based\_baselines.py} & Random, Fixed, Stage, FAO-56 controllers \\
\bottomrule
\end{tabular}
\end{table}


%------------------------------------------------------------------------------
\section{Reward Engineering}
\label{sec:reward}
%------------------------------------------------------------------------------

The reward design is the heart of the proposed solution: it is where the mission
objectives are encoded into learnable feedback. We adopt a \emph{hierarchical}
structure that combines dense daily rewards with a dominant terminal seasonal
reward. The design is presented below as a set of traceable requirements
(REQ000--REQ002), which keeps an explicit link between each reward term and the
mission objective and MoE it serves.

%=========================================================
\subsection{REQ000: Optimise Irrigation and Fertilisation Performance}
%=========================================================

\noindent\textbf{Detailed Requirement.}
The smart-farming system shall optimise irrigation and nitrogen-fertilisation strategies
throughout the crop growing season in order to maximise crop productivity while
minimising water and fertiliser consumption.

\vspace{0.2cm}
\noindent\textbf{Rationale.}
Efficient agricultural management requires balancing productivity and sustainability.
Excessive irrigation wastes water and may increase runoff and drainage losses, while
excessive fertilisation can increase nutrient leaching and environmental impacts.
Therefore, the reinforcement-learning agent must learn resource-efficient management
policies that maintain healthy crop conditions and maximise final crop yield.

\vspace{0.2cm}
\noindent\textbf{Decomposition into Adaptive Requirements.}
To operationalise this objective, REQ000 is decomposed into the following adaptive
requirements:
\begin{itemize}
    \item \textbf{REQ001:} Optimise Daily Reward ($R_{\text{daily}}$).
    \item \textbf{REQ002:} Optimise Seasonal Reward ($R_{\text{season}}$).
\end{itemize}

\vspace{0.2cm}
\noindent\textbf{Associated Measure of Effectiveness.}
\begin{itemize}
    \item Final grain yield.
    \item Total irrigation applied.
    \item Total nitrogen fertiliser applied.
    \item Agronomic Nitrogen Efficiency (ANE).
    \item Harvest Index (HIAD).
\end{itemize}

\vspace{0.2cm}
\noindent\textbf{Traceability.}
REQ001 and REQ002 derive directly from REQ000 by providing short-term and long-term
feedback mechanisms that guide the agent toward sustainable crop management.

%=========================================================
\subsection{REQ001: Optimise Daily Reward ($R_{\text{daily}}$)}
%=========================================================

\noindent\textbf{Detailed Requirement.}
At each simulation day $t$, the system shall compute an intermediate reward that
encourages healthy crop-growth conditions, efficient resource utilisation, and reduced
environmental losses.

\vspace{0.2cm}
\noindent\textbf{Rationale.}
The final crop yield is only observed at harvest. Relying exclusively on harvest
feedback creates a sparse-reward problem that slows learning. Therefore, daily rewards
provide dense feedback that guides the agent toward agronomically desirable decisions
throughout the season.

\vspace{0.2cm}
\noindent\textbf{Associated Measure of Effectiveness.}
\begin{itemize}
    \item Soil-moisture balance.
    \item Water-stress reduction.
    \item Nitrogen-stress reduction.
    \item Nitrogen-uptake promotion.
    \item Reduced irrigation and fertiliser usage.
    \item Reduced runoff, drainage, and nutrient losses.
\end{itemize}

\vspace{0.2cm}
\noindent\textbf{Operational Definition and Indicators.}
\begin{itemize}
  \item \textbf{Water Management Indicators}
  \begin{itemize}
      \item $\mathrm{SWFAC}_t$: Water-stress factor.
      \item $\mathrm{SMR}_t$: Soil-moisture ratio.
  \end{itemize}
  \begin{equation}
  \mathrm{SMR}_t=\frac{\mathrm{SWTD}_t}{\mathrm{SWXD}_t}
  \end{equation}
  where:
  \begin{itemize}
      \item $\mathrm{SWTD}_t$: Total soil water.
      \item $\mathrm{SWXD}_t$: Available water capacity.
  \end{itemize}

  \item \textbf{Nitrogen Management Indicators}
  \begin{itemize}
      \item $\mathrm{NSTRES}_t$: Nitrogen-stress factor.
      \item $\mathrm{TRNU}_t$: Daily nitrogen uptake.
  \end{itemize}

  \item \textbf{Resource Consumption Indicators}
  \begin{itemize}
      \item $\mathrm{IRRC}_t$: Daily irrigation amount.
      \item $\mathrm{FERN}_t$: Daily nitrogen-fertiliser amount.
  \end{itemize}

  \item \textbf{Environmental Loss Indicators}
  \begin{itemize}
      \item $\mathrm{ROFC}_t$: Surface runoff.
      \item $\mathrm{DRNC}_t$: Deep drainage.
      \item $\mathrm{NLEACH}_t$: Nitrate leaching.
      \item $\mathrm{DENIT}_t$: Denitrification losses.
  \end{itemize}
\end{itemize}

\vspace{0.2cm}
\noindent\textbf{Method of Measurement.}
\begin{enumerate}
  \item \textbf{Water Management Reward ($R_{\text{water}}$)}
  \begin{equation}
  R_{\text{water},t}= w_m R_{\text{moisture},t} + w_s\,(1-\mathrm{SWFAC}_t)
  \end{equation}
  where
 \begin{equation}
R_{\text{moisture},t}=
\begin{cases}
+c_1,
& \text{if } \mathrm{SMR}_t \in [\mathrm{SMR}_{\text{low}},\,\mathrm{SMR}_{\text{high}}]
\\[2mm]
-c_2(\mathrm{SMR}_{\text{low}}-\mathrm{SMR}_t)^2,
& \text{if } \mathrm{SMR}_t < \mathrm{SMR}_{\text{low}}
\\[2mm]
-c_3(\mathrm{SMR}_t-\mathrm{SMR}_{\text{high}})^2,
& \text{if } \mathrm{SMR}_t > \mathrm{SMR}_{\text{high}}
\end{cases}
\end{equation}
  This component encourages maintaining soil moisture within an agronomically desirable
  range while minimising water stress.

  \item \textbf{Fertilisation Reward ($R_{\text{fert}}$)}
  \begin{equation}
  R_{\text{fert},t}= w_n\,(1-\mathrm{NSTRES}_t) + w_u\,\mathrm{TRNU}_t
  \end{equation}
  This reward promotes adequate nitrogen availability and efficient crop nitrogen
  uptake.

  \item \textbf{Resource Consumption Penalty ($R_{\text{resource}}$)}
  \begin{equation}
  R_{\text{resource},t}= w_{\text{irr}}\,\mathrm{IRRC}_t + w_{\text{fert}}\,\mathrm{FERN}_t
  \end{equation}
  This component discourages excessive irrigation and fertilisation while allowing
  necessary applications.

  \item \textbf{Environmental Loss Penalty ($R_{\text{losses}}$)}
  \begin{equation}
  R_{\text{losses},t}= w_r\,\mathrm{ROFC}_t + w_d\,\mathrm{DRNC}_t + w_l\,\mathrm{NLEACH}_t + w_{\text{dn}}\,\mathrm{DENIT}_t
  \end{equation}
  This penalty reduces environmentally harmful practices and resource wastage.
\end{enumerate}

\vspace{0.2cm}
\noindent\textbf{Daily Reward Computation.}
\begin{equation}
R_{\text{daily},t}= R_{\text{water},t} + R_{\text{fert},t} - R_{\text{resource},t} - R_{\text{losses},t}
\label{eq:daily_reward}
\end{equation}
The daily reward evaluates the quality of the management decision taken on day $t$
rather than final crop productivity. In the implemented \texttt{SmartFarmSoSEnv},
each sub-term is normalised to $[-1,1]$ or $[0,1]$ before weighting, and the
combined daily signal is scaled by $\alpha=0.01$ so that seasonal objectives
dominate at harvest.

%=========================================================
\subsection{REQ002: Optimise Seasonal Reward ($R_{\text{season}}$)}
%=========================================================

\noindent\textbf{Detailed Requirement.}
At harvest, the system shall compute a terminal reward that evaluates the overall
success of the growing season based on productivity, efficiency, and resource
utilisation.

\vspace{0.2cm}
\noindent\textbf{Rationale.}
The primary mission objective is to maximise crop yield while using water and fertiliser
efficiently. Therefore, the seasonal reward represents the dominant optimisation
objective.

\vspace{0.2cm}
\noindent\textbf{Associated Measure of Effectiveness.}
\begin{itemize}
    \item Grain yield.
    \item Harvest Index.
    \item Agronomic Nitrogen Efficiency.
    \item Seasonal water consumption.
    \item Seasonal fertiliser consumption.
\end{itemize}

\vspace{0.2cm}
\noindent\textbf{Operational Definition and Indicators.}
\begin{itemize}
  \item $\mathrm{GRNWT}_{\text{final}}$: Final grain yield measured at harvest.
  \item Harvest Index:
  \begin{equation}
  \mathrm{HIAD}_{\text{final}}= \frac{\text{Grain Weight}}{\text{Total Above-Ground Biomass}}
  \end{equation}
  \item Agronomic Nitrogen Efficiency:
  \begin{equation}
  \mathrm{ANE}_{\text{final}}= \frac{Y_N-Y_0}{N_{\text{applied}}}
  \end{equation}
  where:
  \begin{itemize}
      \item $Y_N$: Yield with nitrogen application.
      \item $Y_0$: Yield without nitrogen application.
      \item $N_{\text{applied}}$: Total nitrogen applied.
  \end{itemize}
  \item Seasonal irrigation:
  \begin{equation}
  \mathrm{IRRC}_{\text{season}}= \sum_{t=1}^{T} \mathrm{IRRC}_t
  \end{equation}
  \item Seasonal fertilisation:
  \begin{equation}
  \mathrm{FERN}_{\text{season}}= \sum_{t=1}^{T} \mathrm{FERN}_t
  \end{equation}
\end{itemize}

\vspace{0.2cm}
\noindent\textbf{Method of Measurement.}
The terminal reward is computed at harvest as:
\begin{equation}
R_{\text{season}}= w_y\,\mathrm{GRNWT}_{\text{final}} + w_h\,\mathrm{HIAD}_{\text{final}} + w_a\,\mathrm{ANE}_{\text{final}}
            - w_w\,\mathrm{IRRC}_{\text{season}} - w_f\,\mathrm{FERN}_{\text{season}}
\label{eq:season_reward}
\end{equation}
where:
\begin{itemize}
    \item $w_y$: Yield weight (dominant term).
    \item $w_h$: Harvest-index weight.
    \item $w_a$: Agronomic-nitrogen-efficiency weight.
    \item $w_w$: Water-consumption penalty weight.
    \item $w_f$: Fertiliser-consumption penalty weight.
\end{itemize}

\vspace{0.2cm}
\noindent\textbf{Final Transition Reward.}
For non-terminal days:
\begin{equation}
R_t = R_{\text{daily},t}
\end{equation}
For the harvest day:
\begin{equation}
R_T = R_{\text{daily},T} + R_{\text{season}}
\label{eq:transition_scalar}
\end{equation}
This hierarchical reward structure combines dense daily feedback with long-term seasonal
objectives, enabling efficient learning while preserving the primary goal of maximising
sustainable crop productivity.

%=========================================================
\subsection{From Scalar Rewards to a Multi-Reward Vector}
%=========================================================
The requirements above define the reward \emph{terms}; the multi-reward formulation
groups them into the objective channels that the MORL agents optimise jointly. In this
work we consider four objectives and train accordingly:
\begin{equation}
\mathbf{R}_t = \big(\, R^{\text{water}}_t,\; R^{\text{yield}}_t,\; R^{\text{nitrogen}}_t,\; R^{\text{balanced}}_t \,\big)^{\!\top}
\end{equation}
where $R^{\text{water}}$ aggregates the water-management and irrigation-penalty terms,
$R^{\text{yield}}$ is dominated by the seasonal productivity term, $R^{\text{nitrogen}}$
aggregates the fertilisation reward and nitrogen-loss penalties, and
$R^{\text{balanced}}$ is the scalarised composite that trades the three off against one
another. A scalarised return under a preference vector
$\boldsymbol{\omega}\in\Delta$ (the probability simplex) is then
\begin{equation}
G_{\boldsymbol{\omega}} = \sum_{t} \gamma^{t}\, \boldsymbol{\omega}^{\top}\mathbf{R}_t ,
\qquad \boldsymbol{\omega}\succeq 0,\ \ \textstyle\sum_i \omega_i = 1 .
\label{eq:scalarized}
\end{equation}
Single-objective baselines correspond to placing all preference mass on one channel,
while the MORL agents learn across the whole simplex. In the PC-PPO and CAPQL
implementations, the three primary corners map to yield priority, nitrogen-efficiency
priority, and water-efficiency priority; the balanced channel corresponds to uniform
preference sampling over the simplex.


%=========================================================
\subsection{Scalar Reward Variants for Specialised PPO Baselines}
%=========================================================

Beyond the balanced Base reward (\texttt{02\_smart\_farm\_env.py}), two scalar
variants reweight daily and seasonal terms to probe single-resource priorities
without changing the underlying REQ001--REQ002 structure. These variants train
separate PPO agents (Chapter~4); Table~\ref{tab:reward_variants} records the
parameter differences taken from the implementation files.

\begin{table}[ht]
\centering
\caption{Scalar reward variant parameters (Base, Water-min, N-min).}
\label{tab:reward_variants}
\small
\renewcommand{\arraystretch}{1.15}
\begin{tabular}{lccc}
\toprule
\textbf{Parameter} & \textbf{Base} & \textbf{Water-min} & \textbf{N-min} \\
\midrule
\multicolumn{4}{l}{\textit{Daily blend weights}} \\
\quad $W_{\text{WATER}}$    & 0.40 & 0.20 & 0.35 \\
\quad $W_{\text{FERT}}$     & 0.30 & 0.30 & 0.20 \\
\quad $W_{\text{RESOURCE}}$ & 0.15 & 0.35 & 0.30 \\
\quad $W_{\text{LOSSES}}$   & 0.15 & 0.15 & 0.15 \\
\midrule
\multicolumn{4}{l}{\textit{Seasonal blend weights}} \\
\quad $W_{\text{YIELD}}$    & 0.40 & 0.25 & 0.25 \\
\quad $W_{\text{HIAD}}$     & 0.10 & 0.10 & 0.10 \\
\quad $W_{\text{ANE}}$      & 0.30 & 0.30 & 0.40 \\
\quad $W_{\text{PENALTY}}$  & 0.20 & 0.35 & 0.25 \\
\midrule
\multicolumn{4}{l}{\textit{Seasonal penalty thresholds}} \\
\quad IRRIG\_EXCESS (mm)    & 500  & 150  & 500 \\
\quad N\_EXCESS (kg N ha$^{-1}$) & 200 & 200 & 100 \\
\midrule
\multicolumn{4}{l}{\textit{$R_{\text{penalty}}$ split (water / N)}} \\
\quad $\lambda_W : \lambda_N$ & 50:50 & 85:15 & 30:70 \\
\midrule
N-min yield floor on $R_{\text{ane}}$ & --- & --- & 5000 kg ha$^{-1}$ \\
\bottomrule
\end{tabular}
\end{table}

The Water-min variant reduces moisture-chasing and irrigation-excess tolerance;
the N-min variant strengthens N-application penalties and prioritises agronomic
N efficiency, with a yield floor preventing a zero-fertiliser exploit.


%=========================================================
\subsection{Implementation Constants and Normalisation}
%=========================================================

All sub-reward terms are clipped to $[-1,+1]$ or $[0,1]$ before weighting.
The combined daily signal is scaled by $\alpha=\texttt{DAILY\_SCALE}=0.01$.
Moisture-ratio thresholds are $\mathrm{SMR}_{\text{low}}=0.40$ and
$\mathrm{SMR}_{\text{high}}=0.80$. Seasonal yield, harvest index, and ANE are
normalised against \texttt{BASELINE\_YIELD}$=7620$,
\texttt{BASELINE\_HIAD}$=0.45$, and \texttt{BASELINE\_ANE}$=40$~kg grain per
kg~N. In the implemented Base seasonal reward, the positive terms use weights
$W_{\text{YIELD}}=0.40$, $W_{\text{HIAD}}=0.10$, $W_{\text{ANE}}=0.30$ and the
resource penalty uses $W_{\text{PENALTY}}=0.20$, grouped as
$R_{\text{season}} = W_{\text{YIELD}} R_{\text{yield}} + W_{\text{HIAD}} R_{\text{hiad}}
+ W_{\text{ANE}} R_{\text{ane}} - W_{\text{PENALTY}} R_{\text{penalty}}$.


%------------------------------------------------------------------------------
\section{Multi-Objective Learning Methodology}
\label{sec:algos}
%------------------------------------------------------------------------------

Three learner families implement the proposed solution: scalar PPO baselines
(fixed reward shaping), Preference-Conditioned PPO (PC-PPO), and
Concave-Augmented Pareto Q-Learning (CAPQL). Rule-based controllers provide
non-learning reference policies; their evaluation is deferred to Chapter~4.
Training hyperparameters (rollout length, learning rate, \texttt{VecNormalize})
are likewise reported in Chapter~4.

\subsection{Scalar PPO Baselines}

PPO~\cite{schulman2017ppo} maximises the clipped surrogate
\begin{equation}
\mathcal{L}^{\text{CLIP}}(\theta) =
\mathbb{E}_t\!\Big[\, \min\big( \rho_t(\theta)\,\hat{A}_t,\;
\operatorname{clip}(\rho_t(\theta),\,1-\epsilon,\,1+\epsilon)\,\hat{A}_t \big) \Big],
\label{eq:ppoclip}
\end{equation}
where $\hat{A}_t$ is the generalised advantage estimate~\cite{schulman2016gae}
and $\rho_t(\theta)=\pi_\theta(a_t\mid s_t)/\pi_{\theta_{\text{old}}}(a_t\mid s_t)$.
The scalar PPO agent uses the hierarchical reward of
REQ001--REQ002 (Equations~\eqref{eq:daily_reward}--\eqref{eq:transition_scalar}) with an
11-dimensional observation. Additional specialised variants reweight daily
components toward water saving or nitrogen saving (Table~\ref{tab:reward_variants}).
These agents provide rule-free reference policies against which the
multi-objective learners are compared in Chapter~4.
Algorithm~\ref{alg:scalar_ppo} summarises the on-policy training loop.

\begin{algorithm}[ht]
\caption{Scalar PPO training on \texttt{SmartFarmSoSEnv}}
\label{alg:scalar_ppo}
\begin{algorithmic}[1]
\Require Environment $\mathcal{E}$, total steps $N$, rollout length $L$, clip $\epsilon$, GAE $\lambda$, discount $\gamma$
\State Initialise policy $\pi_\theta$ and value network $V_\phi$; $\theta_{\text{old}} \leftarrow \theta$
\While{environment steps $< N$}
  \State Collect rollout $\mathcal{D}=\{(s_t,a_t,r_t)\}_{t=1}^{L}$ with $a_t \sim \pi_{\theta_{\text{old}}}(\cdot \mid s_t)$
  \State Compute returns $\hat{R}_t$ and advantages $\hat{A}_t$ using GAE($\gamma,\lambda$)
  \For{$k = 1$ to \texttt{n\_epochs}}
    \For{each minibatch $\mathcal{B} \subset \mathcal{D}$}
      \State $\rho_t(\theta) \leftarrow \pi_\theta(a_t \mid s_t)\,/\,\pi_{\theta_{\text{old}}}(a_t \mid s_t)$
      \State Update $\theta$ by ascending $\mathcal{L}^{\text{CLIP}}(\theta)$ on $\mathcal{B}$
      \State Update $\phi$ by minimising the value-function loss on $\mathcal{B}$
    \EndFor
  \EndFor
  \State $\theta_{\text{old}} \leftarrow \theta$
\EndWhile
\State \textbf{return} trained policy $\pi_\theta$
\end{algorithmic}
\end{algorithm}


\subsection{Preference-Conditioned PPO (PC-PPO)}
\label{subsec:pcppo}

PC-PPO is \emph{Preference-Conditioned}, not a separate constrained-optimisation
variant: the same PPO clipped objective is used, but the observation includes
$\boldsymbol{\omega}$ and the multi-reward vector is scalarised through
$\boldsymbol{\omega}^{\top}\mathbf{R}_t$ as in Equation~\eqref{eq:scalarized}.
Each episode samples a new preference $\boldsymbol{\omega}\sim\mathrm{Dirichlet}(1,1,1)$,
so one network learns a family of policies over the simplex. At evaluation,
the same model is queried at yield-, N-efficiency-, and water-efficiency-focused
corners $(1,0,0)$, $(0,1,0)$, $(0,0,1)$ and at a balanced profile
$(\tfrac{1}{3},\tfrac{1}{3},\tfrac{1}{3})$, matching the three-dimensional
terminal vector $\mathbf{R}_T=[R_{\text{yield}}, R_{\text{ane}}, R_{\text{water\_eff}}]^\top$
returned by \texttt{02\_smart\_farm\_env\_pcppo.py}.

PC-PPO uses the 14-dimensional observation
$[s_t;\boldsymbol{\omega}]$ with \texttt{VecNormalize} on observations and rewards.
Its main strength is stable on-policy learning on the costly DSSAT simulator while
still covering multiple farmer priorities with a single model.
Algorithm~\ref{alg:pcppo} details the episode-level preference sampling and
PPO update procedure.

\begin{algorithm}[ht]
\caption{Preference-Conditioned PPO (PC-PPO)}
\label{alg:pcppo}
\begin{algorithmic}[1]
\Require Environment $\mathcal{E}$, total steps $N$, rollout length $L$
\State Initialise $\pi_\theta$, $V_\phi$; $\theta_{\text{old}} \leftarrow \theta$
\While{environment steps $< N$}
  \State Sample $\boldsymbol{\omega} \sim \mathrm{Dirichlet}(\mathbf{1})$ for the upcoming episode
  \State Reset $\mathcal{E}$; $t \leftarrow 0$; \textbf{done} $\leftarrow$ \textbf{false}
  \State $\mathcal{D} \leftarrow \emptyset$
  \While{$|\mathcal{D}| < L$ \textbf{and not} \textbf{done}}
    \State $\tilde{s}_t \leftarrow [s_t;\,\boldsymbol{\omega}]$
    \State $a_t \sim \pi_{\theta_{\text{old}}}(\cdot \mid \tilde{s}_t)$
    \State Execute $(s_{t+1}, r_t, \textbf{done}) \leftarrow \mathcal{E}.\textsc{Step}(a_t)$
    \State Store $(\tilde{s}_t, a_t, r_t)$ in $\mathcal{D}$; $t \leftarrow t + 1$
  \EndWhile
  \State Compute GAE advantages $\{\hat{A}_t\}$ and returns $\{\hat{R}_t\}$ on $\mathcal{D}$
  \For{$k = 1$ to \texttt{n\_epochs}}
    \For{each minibatch $\mathcal{B} \subset \mathcal{D}$}
      \State Update $\theta, \phi$ with clipped PPO loss (Equation~\eqref{eq:ppoclip})
    \EndFor
  \EndFor
  \State $\theta_{\text{old}} \leftarrow \theta$
\EndWhile
\State \textbf{return} $\pi_\theta$
\end{algorithmic}
\end{algorithm}


\subsection{Concave-Augmented Pareto Q-Learning (CAPQL)}
\label{subsec:capql}

CAPQL~\cite{lu2023capql} takes an off-policy route based on TD3. A twin
preference-conditioned critic learns vector-valued returns, while the actor
outputs deterministic daily actions. The terminal target uses linear
scalarisation plus a concave augmentation term:
\begin{equation}
  y = R_{\text{daily}} + \mathbb{1}_{t=T}\!\left(
    \boldsymbol{\omega}^{\top}\mathbf{R}_T
    + \lambda \cdot \mathrm{mean}_i\big[\log(R_i + 1 + \varepsilon)\big]
  \right),
  \label{eq:capql_target}
\end{equation}
with $\lambda=0.5$ and $\varepsilon=0.1$ in our implementation. Hindsight
preference relabeling resamples $\boldsymbol{\omega}$ for minibatch
transitions, improving coverage of the objective space without rerunning
DSSAT episodes.

CAPQL is attractive when many preference values must be explored from limited
simulation budget: one network can be queried at any $\boldsymbol{\omega}$ after
training.
Algorithm~\ref{alg:capql} summarises the TD3 interaction loop, replay storage,
hindsight preference relabeling, and concave-augmented critic--actor updates.

\begin{algorithm}[ht]
\caption{Concave-Augmented Pareto Q-Learning (CAPQL)}
\label{alg:capql}
\begin{algorithmic}[1]
\Require Environment $\mathcal{E}$, replay buffer $\mathcal{B}$, total steps $N$, batch size $B$, policy delay $d$
\State Initialise actor $\mu_\theta$, critics $Q_{\phi_1}, Q_{\phi_2}$, and target networks $\mu_{\bar\theta}, Q_{\bar\phi}$
\For{$n = 1$ to $N$}
  \If{episode start}
    \State Sample $\boldsymbol{\omega} \sim \mathrm{Dirichlet}(\mathbf{1})$; reset $\mathcal{E}$
  \EndIf
  \State $\tilde{s}_t \leftarrow [s_t;\,\boldsymbol{\omega}]$
  \State $a_t \leftarrow \mu_\theta(\tilde{s}_t) + \text{exploration noise}$; step $\mathcal{E}$
  \State Store $(s_t, a_t, R_{\text{daily},t}, \mathbf{R}_T\,\mathbb{1}_{t=T}, s_{t+1}, \textbf{done}, \boldsymbol{\omega})$ in $\mathcal{B}$
  \If{$|\mathcal{B}| \ge B$ \textbf{and} $n > \texttt{warmup}$}
    \State Sample minibatch; relabel $\boldsymbol{\omega}' \sim \mathrm{Dirichlet}(\mathbf{1})$ per transition
    \State $\tilde{s} \leftarrow [s;\boldsymbol{\omega}']$,\; $\tilde{s}' \leftarrow [s';\boldsymbol{\omega}']$
    \State $r_{\text{aug}} \leftarrow R_{\text{daily}} + \textbf{done}\cdot\big(\boldsymbol{\omega}'^{\top}\mathbf{R}_T + \lambda\,\mathrm{mean}_i\log(R_i{+}1{+}\varepsilon)\big)$
    \State $a' \leftarrow \mathrm{clip}\big(\mu_{\bar\theta}(\tilde{s}') + \text{smoothing noise}\big)$
    \State $y \leftarrow r_{\text{aug}} + \gamma(1-\textbf{done})\min_{j \in \{1,2\}} Q_{\bar\phi_j}(\tilde{s}', a')$
    \State Update $\phi_1, \phi_2$ by minimising $\|Q_{\phi_j}(\tilde{s},a) - y\|^2$
    \If{$n \bmod d = 0$}
      \State Update $\theta$ by ascending $\mathbb{E}\big[\min_j Q_{\phi_j}(\tilde{s}, \mu_\theta(\tilde{s}))\big]$
      \State Soft-update $\bar\theta \leftarrow \tau\theta + (1{-}\tau)\bar\theta$;\;
            $\bar\phi \leftarrow \tau\phi + (1{-}\tau)\bar\phi$
    \EndIf
  \EndIf
\EndFor
\State \textbf{return} $\mu_\theta$
\end{algorithmic}
\end{algorithm}


\subsection{Justification of Algorithmic Choices}

\begin{table}[ht]
\centering
\caption{Rationale for the selected learning algorithms.}
\label{tab:algos}
\small
\renewcommand{\arraystretch}{1.3}
\begin{tabular}{>{\raggedright\arraybackslash}p{2.8cm}
                >{\raggedright\arraybackslash}p{2.6cm}
                >{\raggedright\arraybackslash}p{6.3cm}}
\toprule
\textbf{Method} & \textbf{Family} & \textbf{Why chosen here} \\
\midrule
Scalar PPO & On-policy, single-objective &
Stable baseline on \texttt{gym-DSSAT}; hierarchical reward matches REQ000--REQ002;
provides fixed-reward reference policies for comparison. \\
PC-PPO & On-policy, preference-conditioned &
Same PPO core as the baseline; one model serves all preference corners;
observation carries $\boldsymbol{\omega}$; strong training stability. \\
CAPQL & Off-policy, preference-conditioned &
Sample-efficient exploration of the Pareto set; concave augmentation mitigates
plain linear scalarisation; TD3 twin critics reduce Q overestimation. \\
\bottomrule
\end{tabular}
\end{table}

In short, scalar PPO defines what fixed hierarchical rewards can encode
(Algorithm~\ref{alg:scalar_ppo}); PC-PPO extends the same PPO core to a
preference-conditioned policy family (Algorithm~\ref{alg:pcppo}); CAPQL
complements this with an off-policy Pareto-learning formulation
(Algorithm~\ref{alg:capql}).


%------------------------------------------------------------------------------
\section{Integrated Proposed Architecture}
\label{sec:overall}
%------------------------------------------------------------------------------

Figure~\ref{fig:overall} assembles the elements of this chapter into the complete
closed-loop learning system. The DSSAT core, wrapped by \texttt{gym-DSSAT} and
\texttt{SmartFarmSoSEnv}, exposes the daily crop state. The reward module
converts DSSAT outputs into daily and seasonal signals (Section~\ref{sec:reward}).
The agent---scalar PPO, PC-PPO, or CAPQL---consumes state and reward and emits
irrigation and nitrogen actions, closing the loop over one maize season.

%---- FIGURE : Overall proposed model ----------------------------------------
\begin{figure}[ht]
\centering
\resizebox{\textwidth}{!}{%
\begin{tikzpicture}[node distance=1.0cm and 1.8cm]

  \node[core, draw=sosblue, fill=sosblue!10, minimum width=3.4cm, minimum height=1.6cm]
       (env) {gym-DSSAT\\Environment\\\scriptsize (DSSAT CSM core)};

  \node[core, draw=sosgreen, fill=sosgreen!10, below=2.2cm of env,
        minimum width=3.4cm, minimum height=1.6cm] (reward)
        {Reward Module\\\scriptsize $R_{\text{daily}}+R_{\text{season}}$};

  \node[core, draw=sosorange, fill=soscream, right=4.2cm of env,
        minimum width=3.6cm, minimum height=1.6cm] (agent)
        {Multi-Objective RL Agent\\\scriptsize PPO $\,\cdot\,$ PC-PPO $\,\cdot\,$ CAPQL};

  \node[block, draw=sosgreen, fill=sosgreen!8, below=2.2cm of agent,
        minimum width=3.6cm, minimum height=1.5cm, font=\scriptsize, align=center]
        (vec) {Seasonal vector $\mathbf{R}_T$\\yield $\cdot$ N-efficiency\\water efficiency};

  \draw[flow] (env.east) -- node[above,font=\scriptsize]{state $s_t$} (agent.west);
  \draw[ctrl] (agent.south) |- node[below right,font=\scriptsize,pos=0.25]
        {action $a_t=(\mathrm{FERN}_t,\mathrm{IRRC}_t)$} (env.south);
  \draw[flow] (env.south) -- node[left,font=\scriptsize]{DSSAT outputs} (reward.north);
  \draw[flow] (reward.east) -- (vec.west);
  \draw[flow] (vec.north) -- node[right,font=\scriptsize]{$\mathbf{R}_T$} (agent.south);

  \node[block, draw=sosgrey, fill=white, right=1.2cm of agent, font=\scriptsize,
        align=center, minimum width=2.6cm] (eval)
        {Evaluation\\\scriptsize vs.\ rule-based\\and scalar PPO baselines};
  \draw[flow, sosgrey] (agent.east) -- (eval.west);

  \begin{scope}[on background layer]
    \node[draw=sosgrey, dashed, rounded corners, inner sep=8pt,
          fit=(env)(reward)(agent)(vec),
          label={[sosgrey,font=\scriptsize\bfseries]above:Closed RL loop (one growing season)}] {};
  \end{scope}

\end{tikzpicture}%
}
\caption{Overall proposed learning architecture. The DSSAT core supplies daily
  crop state; the reward module computes hierarchical daily and seasonal
  signals; the agent issues irrigation and nitrogen actions, closing the loop
  over one maize season.}
\label{fig:overall}
\end{figure}


\subsection{Traceability from Requirements to Implementation}

Table~\ref{tab:trace} summarises how REQ000--REQ002 connect mission objectives,
MoEs, reward terms, and learner families. This traceability matrix is the
design counterpart to the performance tables reported in Chapter~4.

\begin{table}[ht]
\centering
\caption{Traceability: requirements, MoEs, rewards, and learners.}
\label{tab:trace}
\small
\renewcommand{\arraystretch}{1.2}
\begin{tabular}{>{\raggedright\arraybackslash}p{1.6cm}
                >{\raggedright\arraybackslash}p{2.8cm}
                >{\raggedright\arraybackslash}p{2.6cm}
                >{\raggedright\arraybackslash}p{2.4cm}
                >{\raggedright\arraybackslash}p{2.2cm}}
\toprule
\textbf{Req.} & \textbf{MoE / objective} & \textbf{Reward term} & \textbf{Module} & \textbf{Learner} \\
\midrule
REQ000 & Yield, N, water & $R_{\text{daily}}+R_{\text{season}}$ &
\texttt{02\_smart\_farm\_env*} & PPO, PC-PPO, CAPQL \\
REQ001 & SMR, SWFAC, NSTRES & $R_{\text{water}}, R_{\text{fert}}$ &
\texttt{02\_smart\_farm\_env.py} & All \\
REQ001 & Resource use, losses & $R_{\text{resource}}, R_{\text{losses}}$ &
\texttt{02\_smart\_farm\_env.py} & All \\
REQ002 & GRNWT, HIAD, ANE & $R_{\text{yield}}, R_{\text{hiad}}, R_{\text{ane}}$ &
\texttt{02\_smart\_farm\_env.py} & PPO \\
REQ002 & Season totals & $R_{\text{penalty}}$ &
\texttt{02\_smart\_farm\_env*.py} & PPO variants \\
MORL   & Preference trade-offs & $\mathbf{R}_T$, $\boldsymbol{\omega}$ &
\texttt{02\_smart\_farm\_env\_pcppo.py} & PC-PPO, CAPQL \\
\bottomrule
\end{tabular}
\end{table}


\subsection{Comparison Baselines (Design)}

For evaluation (Chapter~4), four rule-based controllers run inside the same
\texttt{SmartFarmSoSEnv} wrapper: \emph{Random} (unconstrained daily inputs),
\emph{Fixed} (preset N splits and periodic irrigation), \emph{Stage-based}
(management keyed to vegetative development), and \emph{FAO-56} (soil-water-balance
irrigation with a three-split N plan). Using a common environment ensures that
MoEs and reward components are directly comparable across learned and rule-based
policies.


%------------------------------------------------------------------------------
\section{Conclusion}
\label{sec:ch3_conclusion}
%------------------------------------------------------------------------------

This chapter presented the proposed solution end to end: a smart-farming
System-of-Systems formalised around the DSSAT Cropping System Model; an MDP
formulation with explicit observation, action, and software layers;
Measures of Effectiveness tied to mission objectives (Table~\ref{tab:moe});
a hierarchical reward design (REQ000--REQ002) with scalar variants
(Table~\ref{tab:reward_variants}) and a multi-objective extension for
preference-conditioned learning; and three learner families---scalar PPO,
PC-PPO, and CAPQL---integrated in the closed-loop architecture of
Figure~\ref{fig:overall}. The next chapter reports the experimental setup,
training dynamics, and quantitative comparison of all learned and rule-based
policies.


% =============================================================================
% If using BibTeX instead, delete this block and add entries to references.bib
% =============================================================================
\begin{thebibliography}{99}

\bibitem{maier1998}
M.~W. Maier,
``Architecting principles for systems-of-systems,''
\emph{Systems Engineering}, vol.~1, no.~4, pp.~267--284, 1998.

\bibitem{jones2003dssat}
J.~W. Jones, G.~Hoogenboom, C.~H. Porter, K.~J. Boote, W.~D. Batchelor, L.~A. Hunt,
P.~W. Wilkens, U.~Singh, A.~J. Gijsman, and J.~T. Ritchie,
``The DSSAT cropping system model,''
\emph{European Journal of Agronomy}, vol.~18, no.~3--4, pp.~235--265, 2003.

\bibitem{hoogenboom2019}
G.~Hoogenboom, C.~H. Porter, K.~J. Boote, V.~Shelia, P.~W. Wilkens, U.~Singh,
J.~W. White, S.~Asseng, J.~I. Lizaso, et~al.,
``The DSSAT crop modeling ecosystem,'' in
\emph{Advances in Crop Modelling for a Sustainable Agriculture}, Burleigh Dodds Science
Publishing, pp.~173--216, 2019.

\bibitem{gautron2022gymdssat}
R.~Gautron, E.~J. Padr\'on, P.~Preux, J.~Bigot, O.-A. Maillard, and D.~Emukpere,
``gym-DSSAT: a crop model turned into a Reinforcement Learning environment,''
Inria Research Report RR-9460, 2022.

\bibitem{gautron2022review}
R.~Gautron, O.-A. Maillard, P.~Preux, M.~Corbeels, and R.~Sabbadin,
``Reinforcement learning for crop management support: Review, prospects and challenges,''
\emph{Computers and Electronics in Agriculture}, vol.~200, art.~107182, 2022.

\bibitem{schulman2017ppo}
J.~Schulman, F.~Wolski, P.~Dhariwal, A.~Radford, and O.~Klimov,
``Proximal policy optimization algorithms,''
\texttt{arXiv:1707.06347}, 2017.

\bibitem{schulman2016gae}
J.~Schulman, P.~Moritz, S.~Levine, M.~Jordan, and P.~Abbeel,
``High-dimensional continuous control using generalized advantage estimation,'' in
\emph{Proc.\ ICLR}, 2016.

\bibitem{lu2023capql}
H.~Lu, D.~Herman, and Y.~Yu,
``Multi-objective reinforcement learning: Convexity, stationarity and Pareto optimality,''
in \emph{Proc.\ ICLR}, 2023.

\end{thebibliography}

% Required: amsmath, booktabs, graphicx, tikz, backgrounds, algorithm, algpseudocode
%   \usetikzlibrary{positioning, arrows.meta, fit, calc, backgrounds}
% =============================================================================

%---- Colour palette for the diagrams ----------------------------------------
\definecolor{sosblue}{RGB}{38,84,124}
\definecolor{sosgreen}{RGB}{61,133,92}
\definecolor{sosorange}{RGB}{200,120,40}
\definecolor{sosgrey}{RGB}{90,90,90}
\definecolor{soslight}{RGB}{234,240,245}
\definecolor{soscream}{RGB}{248,244,232}

%---- Reusable TikZ node styles ----------------------------------------------
\tikzset{
  sensor/.style   = {rectangle, rounded corners=2pt, draw=sosblue, very thick,
                     fill=soslight, align=center, minimum width=2.5cm,
                     minimum height=1cm, font=\small},
  actuator/.style = {rectangle, rounded corners=2pt, draw=sosorange, very thick,
                     fill=soscream, align=center, minimum width=2.5cm,
                     minimum height=1cm, font=\small},
  core/.style     = {rectangle, rounded corners=4pt, draw=sosgreen, line width=1.2pt,
                     fill=sosgreen!12, align=center, minimum width=3.2cm,
                     minimum height=1.3cm, font=\small\bfseries},
  mission/.style  = {ellipse, draw=sosgrey, very thick, fill=white, align=center,
                     font=\small\itshape, minimum width=3cm},
  flow/.style     = {-{Stealth[length=2.5mm]}, line width=0.9pt, sosgrey},
  ctrl/.style     = {-{Stealth[length=2.5mm]}, line width=0.9pt, sosorange, dashed},
  block/.style    = {rectangle, rounded corners=2pt, draw=sosblue, thick,
                     fill=soslight, align=center, font=\small,
                     minimum width=2.6cm, minimum height=0.9cm},
}


\chapter{Proposed Solution}
\label{chap:proposed}


%------------------------------------------------------------------------------
\section{Introduction}
\label{sec:ch3_intro}
%------------------------------------------------------------------------------

Chapter~\ref{ch:background} positioned DSSAT as the crop-simulation subsystem
inside a smart-farming System-of-Systems. This chapter turns that view into a
concrete decision framework.

The problem is no longer a single irrigation rule. It is a \emph{multi-objective}
management task in which yield, water use, and nitrogen efficiency must be
balanced over a full growing season. The proposed framework combines three
elements:
\begin{enumerate}
  \item a DSSAT-based daily simulation environment (\texttt{gym-DSSAT});
  \item a hierarchical reward design with preference conditioning for
        multi-objective learning;
  \item learning algorithms---scalar PPO baselines, Preference-Conditioned PPO
        (PC-PPO), and Concave-Augmented Pareto Q-Learning (CAPQL).
\end{enumerate}

The chapter is organised as follows. Section~\ref{sec:sos} formalises the
control problem and the smart-farming System-of-Systems architecture.
Section~\ref{sec:env} describes the simulation platform, observation and action
spaces, and software stack. Section~\ref{sec:reward} presents the hierarchical
reward design and scalar reward variants. Section~\ref{sec:algos} defines the
learning algorithms. Section~\ref{sec:overall} assembles the integrated
architecture and traceability from requirements to implementation. Experimental
training protocols and numerical results are reported in the next chapter.


%------------------------------------------------------------------------------
\section{Problem Formulation and System-of-Systems Architecture}
\label{sec:sos}
%------------------------------------------------------------------------------

\subsection{Overview and Rationale}

A modern smart farm is not a single monolithic system but a \emph{System of
Systems}: a collection of heterogeneous, independently operated subsystems---
weather data, soil and crop observations, irrigation and fertilisation
controllers, and crop-modelling platforms---integrated to pursue missions such
as maximising yield, conserving water, and improving nitrogen-use efficiency.
Following Maier~\cite{maier1998}, a true SoS is distinguished from a merely
complex monolithic system by five traits: operational independence, managerial
independence, geographic distribution, emergent behaviour, and evolutionary
development. The smart farm exhibits all five, which is precisely what makes
mission-level decision-making difficult: an irrigation pulse today changes soil
moisture, crop stress, and nitrogen dynamics over the rest of the season.
Section~\ref{sec:env} formalises this task as a Markov decision process; here
we focus on the architectural view.


\subsection{Components of the Smart-Farming SoS}

The constituent systems are grouped by their role in the closed decision loop:
observe field conditions, act on the crop, and decide daily management. At the
centre of this loop sits the \textbf{crop-modelling platform (DSSAT)}, which
integrates weather, soil, and management inputs and produces crop growth, stress
indicators, and environmental fluxes. Monitoring systems read from it; control
systems write into it; the decision layer issues the commands that drive both.

\vspace{0.25cm}
\noindent\textbf{Monitoring systems.}
These constituents supply the observation stream the agent uses each day:
\begin{itemize}
  \item \textbf{Weather data} --- supplies exogenous drivers (temperature,
        rainfall, solar radiation) that the agent observes but cannot control.
  \item \textbf{Soil and crop state} --- reported through DSSAT variables such
        as soil-water content, water-stress factor, nitrogen-stress factor,
        leaf area, and grain weight.
\end{itemize}

\vspace{0.25cm}
\noindent\textbf{Control systems.}
These constituents execute management actions issued by the decision layer:
\begin{itemize}
  \item \textbf{Irrigation controller} --- the actuator through which the agent
        applies a daily water amount $a_t^{W}$ (mm/day).
  \item \textbf{Fertilisation unit} --- the actuator through which the agent
        applies a daily nitrogen amount $a_t^{N}$ (kg N/ha/day).
\end{itemize}

\vspace{0.25cm}
\noindent\textbf{Decision and mission layer.}
Mission objectives---yield, water efficiency, and nitrogen efficiency---condition
the agent through the preference vector $\boldsymbol{\omega}$. The decision
constituent that closes the loop is:
\begin{itemize}
  \item \textbf{Mission manager / RL agent} --- reads observations, respects the
        farmer's preference vector, and issues daily irrigation and fertilisation
        commands to the control systems.
\end{itemize}

Figure~\ref{fig:sos} depicts the architecture. Monitoring systems feed
observations into the mission manager (solid arrows); the mission manager issues
control commands to the actuators (dashed arrows); DSSAT closes the loop by
producing the next crop and soil state.


%---- FIGURE : SoS architecture ------------------------------------------------
\begin{figure}[ht]
\centering
\includegraphics[width=\textwidth]{sos_architecture.png}
\caption{Architecture of the proposed smart-farming System-of-Systems. Icons denote
constituent roles; solid arrows carry observations and simulated state, dashed
orange arrows carry management commands. DSSAT integrates inputs and closes the
daily decision loop.}
\label{fig:sos}
\end{figure}


\subsection{System-of-Systems Characteristics in Our Setting}

Mapping Maier's five traits~\cite{maier1998} onto our design clarifies why a
learning-based formulation is appropriate. \emph{Operational independence}
holds because weather data, DSSAT, and actuators can each be described on
their own. \emph{Managerial independence} holds because these parts are
developed and maintained separately yet must cooperate in practice.
\emph{Geographic distribution} is intrinsic to field deployment.
\emph{Emergent behaviour}---yield and resource efficiency arising from the
interaction of weather, soil, crop, and management---is exactly what the agent
must learn to shape. \emph{Evolutionary development} motivates a framework in
which the crop model, reward design, or learning algorithm can be upgraded
without rebuilding the entire platform.


\subsection{Mission Objectives and Measures of Effectiveness}
\label{sec:moe}

Measures of Effectiveness (MoEs) translate mission objectives into quantities
that can be logged in simulation and encoded in the reward function.
Table~\ref{tab:moe} links each MoE to its mission objective and to the reward
term that implements it. \emph{Seasonal} MoEs evaluate end-of-season success;
\emph{daily} indicators supply the dense feedback used in reward shaping
(Section~\ref{sec:reward}).

\begin{table}[ht]
\centering
\caption{Measures of Effectiveness and their alignment with mission objectives.}
\label{tab:moe}
\small
\renewcommand{\arraystretch}{1.25}
\begin{tabular}{>{\raggedright\arraybackslash}p{4.4cm}
                >{\raggedright\arraybackslash}p{3.2cm}
                >{\raggedright\arraybackslash}p{4.6cm}}
\toprule
\textbf{Measure of Effectiveness} & \textbf{Mission objective} & \textbf{Indicator / reward link} \\
\midrule
Final grain yield ($\mathrm{GRNWT}_{\text{final}}$) & Yield optimisation & $w_y\,\mathrm{GRNWT}_{\text{final}}$ \\
Harvest index ($\mathrm{HIAD}_{\text{final}}$) & Yield quality & $w_h\,\mathrm{HIAD}_{\text{final}}$ \\
Agronomic N efficiency ($\mathrm{ANE}_{\text{final}}$) & N-use efficiency & $w_a\,\mathrm{ANE}_{\text{final}}$ \\
Total irrigation ($\mathrm{IRRC}_{\text{season}}$) & Water-use efficiency & $w_w\,\mathrm{IRRC}_{\text{season}}$ \\
Total nitrogen applied ($\mathrm{FERN}_{\text{season}}$) & N-use efficiency & $w_f\,\mathrm{FERN}_{\text{season}}$ \\
Soil-moisture balance ($\mathrm{SMR}_t$) & Crop health / water stress & $R_{\text{water},t}$ \\
Water-stress factor ($\mathrm{SWFAC}_t$) & Crop health & $R_{\text{water},t}$ \\
Nitrogen-stress factor ($\mathrm{NSTRES}_t$) & Crop health & $R_{\text{fert},t}$ \\
Cumulative N \& water losses & Environmental impact & $R_{\text{losses},t}$ \\
\bottomrule
\end{tabular}
\end{table}


%------------------------------------------------------------------------------
\section{Simulation Platform and Environment Design}
\label{sec:env}
%------------------------------------------------------------------------------

\subsection{The DSSAT Cropping System Model}

The Decision Support System for Agrotechnology Transfer (DSSAT) is a suite of
mechanistic crop models developed over more than three decades and widely
validated by the agronomic community~\cite{jones2003dssat,hoogenboom2019}. At
its heart is the \emph{Cropping System Model} (CSM), which advances the state of
a land unit on a daily time step. Given weather, soil, genetics, and management
inputs, the CSM couples soil-water balance, soil--plant--atmosphere exchange,
and plant growth modules. For maize, this work uses CERES-Maize. Its outputs
include grain yield, biomass, harvest index, and water and nitrogen fluxes such
as runoff, leaching, and denitrification. Figure~\ref{fig:dssat} sketches the
modular structure.


\begin{figure}[ht]
\centering
\begin{tikzpicture}[node distance=0.7cm and 1.2cm]
  \node[core, fill=sosblue!10, draw=sosblue] (w) {Weather};
  \node[core, fill=sosblue!10, draw=sosblue, below=of w] (s) {Soil profile};
  \node[core, fill=sosgreen!15, draw=sosgreen, right=2.2cm of w, yshift=-0.4cm] (csm)
        {DSSAT-CSM\\daily integrator};
  \node[core, fill=soscream, draw=sosorange, right=2.2cm of csm, yshift=0.9cm] (crop)
        {CERES-Maize};
  \node[core, fill=soscream, draw=sosorange, right=2.2cm of csm, yshift=-0.9cm] (mgmt)
        {Management\\(N + irrigation)};
  \node[block, below=1.6cm of csm, minimum width=4.2cm] (out)
        {Outputs: yield, stress, losses, state vector};

  \draw[flow] (w) -- (csm);
  \draw[flow] (s) -- (csm);
  \draw[flow] (crop) -- (csm);
  \draw[flow] (mgmt) -- (csm);
  \draw[flow] (csm) -- (out);
\end{tikzpicture}
\caption{Schematic of the DSSAT Cropping System Model used in this work.}
\label{fig:dssat}
\end{figure}


\subsection{From DSSAT to an RL Environment: gym-DSSAT}

To use DSSAT for reinforcement learning, the simulator must expose a step-wise
interface in which an agent observes a state, applies an action, and receives
the resulting next state and reward. We rely on \texttt{gym-DSSAT}~\cite{gautron2022gymdssat},
which wraps the Fortran simulator in a Gymnasium-compatible API allowing daily
interaction throughout the growing season~\cite{gautron2022review}. A
\texttt{SmartFarmSoSEnv} layer maps DSSAT outputs to the observation vector,
translates the agent's action into DSSAT management events
(\texttt{anfer}, \texttt{amir}), and computes the reward defined in
Section~\ref{sec:reward}.


\subsection{Decision-Process Formulation}

Formally, the control task is a Markov Decision Process
$\langle \mathcal{S}, \mathcal{A}, P, R, \gamma \rangle$. In the
multi-objective setting, the terminal step also exposes a reward vector
$\mathbf{R}_T$ that is scalarised through a preference vector
$\boldsymbol{\omega}$ (Section~\ref{sec:reward}).

\begin{itemize}
  \item \textbf{State / observation $s_t$:} an 11-dimensional vector of
        normalised crop and SoS features: biomass, grain weight, days after
        planting, vegetative stage, leaf area index, cumulative nitrogen and
        irrigation, mean soil water, energy budget, communication quality, and
        fraction of healthy virtual sensors. PC-PPO and CAPQL append the
        3-dimensional preference vector, giving a 14-dimensional observation.
  \item \textbf{Action $a_t$:} the daily management decision
        $a_t = (\mathrm{FERN}_t, \mathrm{IRRC}_t) = (a_t^{N}, a_t^{W})$ with
        $a_t^{N}\in[0,200]$~kg N/ha/day and $a_t^{W}\in[0,50]$~mm/day.
  \item \textbf{Transition $P(s_{t+1}\mid s_t,a_t)$:} realised by the DSSAT
        daily integration; stochasticity enters through random weather resampling
        across episodes.
  \item \textbf{Horizon:} one growing season (${\approx}161$ days); the episode
        terminates at harvest, when the seasonal reward is granted.
\end{itemize}

Irrigation and fertilisation are coupled because both alter the same soil-water
and nitrogen pools that drive crop stress and final yield. A sequential
formulation is therefore more faithful than optimising either input in isolation.


\subsection{Observation and Action Spaces}

Table~\ref{tab:obs} lists the 11-dimensional observation vector exposed to
scalar PPO agents through the \texttt{pc\_env.py} adapter. Each component is
linearly scaled to $[0,1]$ from agronomically meaningful bounds. PC-PPO and
CAPQL append the three-dimensional preference vector
$\boldsymbol{\omega}=(\omega_y,\omega_a,\omega_w)$, yielding a
14-dimensional input.

\begin{table}[ht]
\centering
\caption{Observation vector for scalar PPO (\texttt{pc\_env.py}).}
\label{tab:obs}
\small
\renewcommand{\arraystretch}{1.2}
\begin{tabular}{clll}
\toprule
\textbf{Idx} & \textbf{DSSAT / SoS variable} & \textbf{Scale} & \textbf{Role} \\
\midrule
0  & \texttt{topwt}  & $[0, 20\,000]$ kg ha$^{-1}$ & Above-ground biomass \\
1  & \texttt{grnwt}  & $[0, 12\,000]$ kg ha$^{-1}$ & Grain weight \\
2  & \texttt{dap}    & $[0, 200]$ days             & Days after planting \\
3  & \texttt{vstage} & $[0, 18]$                   & Vegetative stage \\
4  & \texttt{xlai}   & $[0, 7]$                    & Leaf area index \\
5  & Cumulative N    & $[0, 300]$ kg ha$^{-1}$     & Seasonal fertiliser \\
6  & Cumulative W    & $[0, 1000]$ mm              & Seasonal irrigation \\
7  & Mean soil water & $[0, 1]$                    & Profile moisture \\
8  & Energy budget   & $[0, 100]$                  & SoS energy state \\
9  & Comm.\ quality  & $[0, 1]$                    & Link reliability \\
10 & Sensors alive   & $[0, 1]$                    & Virtual sensor health \\
\bottomrule
\end{tabular}
\end{table}

The daily action vector maps to DSSAT management events:
$a_t=(a_t^{N},a_t^{W})=(\texttt{anfer},\texttt{amir})$ with
$a_t^{N}\in[0,200]$~kg N ha$^{-1}$ day$^{-1}$ and
$a_t^{W}\in[0,50]$~mm day$^{-1}$.


\subsection{Episode Protocol and Calibration Anchors}

Each episode simulates one maize growing season of approximately 161 days.
Weather is resampled randomly across episodes (\texttt{random\_weather=True}),
injecting the main source of transition stochasticity. Yield calibration is
anchored to the stage-based baseline at ${\approx}7\,620$~kg ha$^{-1}$
(\texttt{BASELINE\_YIELD}); normalisation upper bounds follow
\texttt{MAX\_EXPECTED\_YIELD}$=1.5\times$ baseline. Fault injection
(\texttt{enable\_faults=False} in the main experiments) is disabled so that
learning focuses on crop-management trade-offs; resilience extensions are left
for future work.


\subsection{Software Architecture and Module Map}

The implementation stacks four layers: the DSSAT core accessed through
\texttt{gym-DSSAT}; a \texttt{SmartFarmSoSEnv} reward wrapper; a
Gymnasium-compatible adapter; and the training script. Table~\ref{tab:software}
maps each module to its role in the proposed solution.

\begin{table}[ht]
\centering
\caption{Software modules implementing the proposed solution.}
\label{tab:software}
\small
\renewcommand{\arraystretch}{1.25}
\begin{tabular}{>{\raggedright\arraybackslash}p{3.6cm}
                >{\raggedright\arraybackslash}p{4.2cm}
                >{\raggedright\arraybackslash}p{5.5cm}}
\toprule
\textbf{Layer} & \textbf{Module} & \textbf{Function} \\
\midrule
DSSAT core      & \texttt{gym\_dssat\_pdi}           & Daily CSM integration via subprocess API \\
Scalar reward   & \texttt{02\_smart\_farm\_env.py}   & Hierarchical scalar reward (REQ001--002) \\
Water-min reward& \texttt{02\_smart\_farm\_env\_water.py} & Reward weights tuned for irrigation saving \\
N-min reward    & \texttt{02\_smart\_farm\_env\_nitrogen.py} & Reward weights tuned for N saving \\
MORL reward     & \texttt{02\_smart\_farm\_env\_pcppo.py} & Terminal objective vector $\mathbf{R}_T$ \\
11-dim adapter  & \texttt{pc\_env.py}                & Box observation for scalar PPO \\
14-dim adapter  & \texttt{pc\_env\_pcppo.py}         & Appends $\boldsymbol{\omega}$ for PC-PPO \\
CAPQL adapter   & \texttt{capql\_env.py}             & PyTorch interface for TD3 training \\
Scalar PPO train& \texttt{04\_pc\_ppo\_quick\_train.py}, \texttt{04b}, \texttt{04c} & Base, Water-min, N-min \\
PC-PPO train    & \texttt{05\_pcppo\_train.py}       & Preference-conditioned PPO \\
CAPQL train     & \texttt{06\_capql\_train.py}       & Concave-augmented Pareto Q-learning \\
Baselines       & \texttt{01\_rule\_based\_baselines.py} & Random, Fixed, Stage, FAO-56 controllers \\
\bottomrule
\end{tabular}
\end{table}


%------------------------------------------------------------------------------
\section{Reward Engineering}
\label{sec:reward}
%------------------------------------------------------------------------------

The reward design is the heart of the proposed solution: it is where the mission
objectives are encoded into learnable feedback. We adopt a \emph{hierarchical}
structure that combines dense daily rewards with a dominant terminal seasonal
reward. The design is presented below as a set of traceable requirements
(REQ000--REQ002), which keeps an explicit link between each reward term and the
mission objective and MoE it serves.

%=========================================================
\subsection{REQ000: Optimise Irrigation and Fertilisation Performance}
%=========================================================

\noindent\textbf{Detailed Requirement.}
The smart-farming system shall optimise irrigation and nitrogen-fertilisation strategies
throughout the crop growing season in order to maximise crop productivity while
minimising water and fertiliser consumption.

\vspace{0.2cm}
\noindent\textbf{Rationale.}
Efficient agricultural management requires balancing productivity and sustainability.
Excessive irrigation wastes water and may increase runoff and drainage losses, while
excessive fertilisation can increase nutrient leaching and environmental impacts.
Therefore, the reinforcement-learning agent must learn resource-efficient management
policies that maintain healthy crop conditions and maximise final crop yield.

\vspace{0.2cm}
\noindent\textbf{Decomposition into Adaptive Requirements.}
To operationalise this objective, REQ000 is decomposed into the following adaptive
requirements:
\begin{itemize}
    \item \textbf{REQ001:} Optimise Daily Reward ($R_{\text{daily}}$).
    \item \textbf{REQ002:} Optimise Seasonal Reward ($R_{\text{season}}$).
\end{itemize}

\vspace{0.2cm}
\noindent\textbf{Associated Measure of Effectiveness.}
\begin{itemize}
    \item Final grain yield.
    \item Total irrigation applied.
    \item Total nitrogen fertiliser applied.
    \item Agronomic Nitrogen Efficiency (ANE).
    \item Harvest Index (HIAD).
\end{itemize}

\vspace{0.2cm}
\noindent\textbf{Traceability.}
REQ001 and REQ002 derive directly from REQ000 by providing short-term and long-term
feedback mechanisms that guide the agent toward sustainable crop management.

%=========================================================
\subsection{REQ001: Optimise Daily Reward ($R_{\text{daily}}$)}
%=========================================================

\noindent\textbf{Detailed Requirement.}
At each simulation day $t$, the system shall compute an intermediate reward that
encourages healthy crop-growth conditions, efficient resource utilisation, and reduced
environmental losses.

\vspace{0.2cm}
\noindent\textbf{Rationale.}
The final crop yield is only observed at harvest. Relying exclusively on harvest
feedback creates a sparse-reward problem that slows learning. Therefore, daily rewards
provide dense feedback that guides the agent toward agronomically desirable decisions
throughout the season.

\vspace{0.2cm}
\noindent\textbf{Associated Measure of Effectiveness.}
\begin{itemize}
    \item Soil-moisture balance.
    \item Water-stress reduction.
    \item Nitrogen-stress reduction.
    \item Nitrogen-uptake promotion.
    \item Reduced irrigation and fertiliser usage.
    \item Reduced runoff, drainage, and nutrient losses.
\end{itemize}

\vspace{0.2cm}
\noindent\textbf{Operational Definition and Indicators.}
\begin{itemize}
  \item \textbf{Water Management Indicators}
  \begin{itemize}
      \item $\mathrm{SWFAC}_t$: Water-stress factor.
      \item $\mathrm{SMR}_t$: Soil-moisture ratio.
  \end{itemize}
  \begin{equation}
  \mathrm{SMR}_t=\frac{\mathrm{SWTD}_t}{\mathrm{SWXD}_t}
  \end{equation}
  where:
  \begin{itemize}
      \item $\mathrm{SWTD}_t$: Total soil water.
      \item $\mathrm{SWXD}_t$: Available water capacity.
  \end{itemize}

  \item \textbf{Nitrogen Management Indicators}
  \begin{itemize}
      \item $\mathrm{NSTRES}_t$: Nitrogen-stress factor.
      \item $\mathrm{TRNU}_t$: Daily nitrogen uptake.
  \end{itemize}

  \item \textbf{Resource Consumption Indicators}
  \begin{itemize}
      \item $\mathrm{IRRC}_t$: Daily irrigation amount.
      \item $\mathrm{FERN}_t$: Daily nitrogen-fertiliser amount.
  \end{itemize}

  \item \textbf{Environmental Loss Indicators}
  \begin{itemize}
      \item $\mathrm{ROFC}_t$: Surface runoff.
      \item $\mathrm{DRNC}_t$: Deep drainage.
      \item $\mathrm{NLEACH}_t$: Nitrate leaching.
      \item $\mathrm{DENIT}_t$: Denitrification losses.
  \end{itemize}
\end{itemize}

\vspace{0.2cm}
\noindent\textbf{Method of Measurement.}
\begin{enumerate}
  \item \textbf{Water Management Reward ($R_{\text{water}}$)}
  \begin{equation}
  R_{\text{water},t}= w_m R_{\text{moisture},t} + w_s\,(1-\mathrm{SWFAC}_t)
  \end{equation}
  where
 \begin{equation}
R_{\text{moisture},t}=
\begin{cases}
+c_1,
& \text{if } \mathrm{SMR}_t \in [\mathrm{SMR}_{\text{low}},\,\mathrm{SMR}_{\text{high}}]
\\[2mm]
-c_2(\mathrm{SMR}_{\text{low}}-\mathrm{SMR}_t)^2,
& \text{if } \mathrm{SMR}_t < \mathrm{SMR}_{\text{low}}
\\[2mm]
-c_3(\mathrm{SMR}_t-\mathrm{SMR}_{\text{high}})^2,
& \text{if } \mathrm{SMR}_t > \mathrm{SMR}_{\text{high}}
\end{cases}
\end{equation}
  This component encourages maintaining soil moisture within an agronomically desirable
  range while minimising water stress.

  \item \textbf{Fertilisation Reward ($R_{\text{fert}}$)}
  \begin{equation}
  R_{\text{fert},t}= w_n\,(1-\mathrm{NSTRES}_t) + w_u\,\mathrm{TRNU}_t
  \end{equation}
  This reward promotes adequate nitrogen availability and efficient crop nitrogen
  uptake.

  \item \textbf{Resource Consumption Penalty ($R_{\text{resource}}$)}
  \begin{equation}
  R_{\text{resource},t}= w_{\text{irr}}\,\mathrm{IRRC}_t + w_{\text{fert}}\,\mathrm{FERN}_t
  \end{equation}
  This component discourages excessive irrigation and fertilisation while allowing
  necessary applications.

  \item \textbf{Environmental Loss Penalty ($R_{\text{losses}}$)}
  \begin{equation}
  R_{\text{losses},t}= w_r\,\mathrm{ROFC}_t + w_d\,\mathrm{DRNC}_t + w_l\,\mathrm{NLEACH}_t + w_{\text{dn}}\,\mathrm{DENIT}_t
  \end{equation}
  This penalty reduces environmentally harmful practices and resource wastage.
\end{enumerate}

\vspace{0.2cm}
\noindent\textbf{Daily Reward Computation.}
\begin{equation}
R_{\text{daily},t}= R_{\text{water},t} + R_{\text{fert},t} - R_{\text{resource},t} - R_{\text{losses},t}
\label{eq:daily_reward}
\end{equation}
The daily reward evaluates the quality of the management decision taken on day $t$
rather than final crop productivity. In the implemented \texttt{SmartFarmSoSEnv},
each sub-term is normalised to $[-1,1]$ or $[0,1]$ before weighting, and the
combined daily signal is scaled by $\alpha=0.01$ so that seasonal objectives
dominate at harvest.

%=========================================================
\subsection{REQ002: Optimise Seasonal Reward ($R_{\text{season}}$)}
%=========================================================

\noindent\textbf{Detailed Requirement.}
At harvest, the system shall compute a terminal reward that evaluates the overall
success of the growing season based on productivity, efficiency, and resource
utilisation.

\vspace{0.2cm}
\noindent\textbf{Rationale.}
The primary mission objective is to maximise crop yield while using water and fertiliser
efficiently. Therefore, the seasonal reward represents the dominant optimisation
objective.

\vspace{0.2cm}
\noindent\textbf{Associated Measure of Effectiveness.}
\begin{itemize}
    \item Grain yield.
    \item Harvest Index.
    \item Agronomic Nitrogen Efficiency.
    \item Seasonal water consumption.
    \item Seasonal fertiliser consumption.
\end{itemize}

\vspace{0.2cm}
\noindent\textbf{Operational Definition and Indicators.}
\begin{itemize}
  \item $\mathrm{GRNWT}_{\text{final}}$: Final grain yield measured at harvest.
  \item Harvest Index:
  \begin{equation}
  \mathrm{HIAD}_{\text{final}}= \frac{\text{Grain Weight}}{\text{Total Above-Ground Biomass}}
  \end{equation}
  \item Agronomic Nitrogen Efficiency:
  \begin{equation}
  \mathrm{ANE}_{\text{final}}= \frac{Y_N-Y_0}{N_{\text{applied}}}
  \end{equation}
  where:
  \begin{itemize}
      \item $Y_N$: Yield with nitrogen application.
      \item $Y_0$: Yield without nitrogen application.
      \item $N_{\text{applied}}$: Total nitrogen applied.
  \end{itemize}
  \item Seasonal irrigation:
  \begin{equation}
  \mathrm{IRRC}_{\text{season}}= \sum_{t=1}^{T} \mathrm{IRRC}_t
  \end{equation}
  \item Seasonal fertilisation:
  \begin{equation}
  \mathrm{FERN}_{\text{season}}= \sum_{t=1}^{T} \mathrm{FERN}_t
  \end{equation}
\end{itemize}

\vspace{0.2cm}
\noindent\textbf{Method of Measurement.}
The terminal reward is computed at harvest as:
\begin{equation}
R_{\text{season}}= w_y\,\mathrm{GRNWT}_{\text{final}} + w_h\,\mathrm{HIAD}_{\text{final}} + w_a\,\mathrm{ANE}_{\text{final}}
            - w_w\,\mathrm{IRRC}_{\text{season}} - w_f\,\mathrm{FERN}_{\text{season}}
\label{eq:season_reward}
\end{equation}
where:
\begin{itemize}
    \item $w_y$: Yield weight (dominant term).
    \item $w_h$: Harvest-index weight.
    \item $w_a$: Agronomic-nitrogen-efficiency weight.
    \item $w_w$: Water-consumption penalty weight.
    \item $w_f$: Fertiliser-consumption penalty weight.
\end{itemize}

\vspace{0.2cm}
\noindent\textbf{Final Transition Reward.}
For non-terminal days:
\begin{equation}
R_t = R_{\text{daily},t}
\end{equation}
For the harvest day:
\begin{equation}
R_T = R_{\text{daily},T} + R_{\text{season}}
\label{eq:transition_scalar}
\end{equation}
This hierarchical reward structure combines dense daily feedback with long-term seasonal
objectives, enabling efficient learning while preserving the primary goal of maximising
sustainable crop productivity.

%=========================================================
\subsection{From Scalar Rewards to a Multi-Reward Vector}
%=========================================================
The requirements above define the reward \emph{terms}; the multi-reward formulation
groups them into the objective channels that the MORL agents optimise jointly. In this
work we consider four objectives and train accordingly:
\begin{equation}
\mathbf{R}_t = \big(\, R^{\text{water}}_t,\; R^{\text{yield}}_t,\; R^{\text{nitrogen}}_t,\; R^{\text{balanced}}_t \,\big)^{\!\top}
\end{equation}
where $R^{\text{water}}$ aggregates the water-management and irrigation-penalty terms,
$R^{\text{yield}}$ is dominated by the seasonal productivity term, $R^{\text{nitrogen}}$
aggregates the fertilisation reward and nitrogen-loss penalties, and
$R^{\text{balanced}}$ is the scalarised composite that trades the three off against one
another. A scalarised return under a preference vector
$\boldsymbol{\omega}\in\Delta$ (the probability simplex) is then
\begin{equation}
G_{\boldsymbol{\omega}} = \sum_{t} \gamma^{t}\, \boldsymbol{\omega}^{\top}\mathbf{R}_t ,
\qquad \boldsymbol{\omega}\succeq 0,\ \ \textstyle\sum_i \omega_i = 1 .
\label{eq:scalarized}
\end{equation}
Single-objective baselines correspond to placing all preference mass on one channel,
while the MORL agents learn across the whole simplex. In the PC-PPO and CAPQL
implementations, the three primary corners map to yield priority, nitrogen-efficiency
priority, and water-efficiency priority; the balanced channel corresponds to uniform
preference sampling over the simplex.


%=========================================================
\subsection{Scalar Reward Variants for Specialised PPO Baselines}
%=========================================================

Beyond the balanced Base reward (\texttt{02\_smart\_farm\_env.py}), two scalar
variants reweight daily and seasonal terms to probe single-resource priorities
without changing the underlying REQ001--REQ002 structure. These variants train
separate PPO agents (Chapter~4); Table~\ref{tab:reward_variants} records the
parameter differences taken from the implementation files.

\begin{table}[ht]
\centering
\caption{Scalar reward variant parameters (Base, Water-min, N-min).}
\label{tab:reward_variants}
\small
\renewcommand{\arraystretch}{1.15}
\begin{tabular}{lccc}
\toprule
\textbf{Parameter} & \textbf{Base} & \textbf{Water-min} & \textbf{N-min} \\
\midrule
\multicolumn{4}{l}{\textit{Daily blend weights}} \\
\quad $W_{\text{WATER}}$    & 0.40 & 0.20 & 0.35 \\
\quad $W_{\text{FERT}}$     & 0.30 & 0.30 & 0.20 \\
\quad $W_{\text{RESOURCE}}$ & 0.15 & 0.35 & 0.30 \\
\quad $W_{\text{LOSSES}}$   & 0.15 & 0.15 & 0.15 \\
\midrule
\multicolumn{4}{l}{\textit{Seasonal blend weights}} \\
\quad $W_{\text{YIELD}}$    & 0.40 & 0.25 & 0.25 \\
\quad $W_{\text{HIAD}}$     & 0.10 & 0.10 & 0.10 \\
\quad $W_{\text{ANE}}$      & 0.30 & 0.30 & 0.40 \\
\quad $W_{\text{PENALTY}}$  & 0.20 & 0.35 & 0.25 \\
\midrule
\multicolumn{4}{l}{\textit{Seasonal penalty thresholds}} \\
\quad IRRIG\_EXCESS (mm)    & 500  & 150  & 500 \\
\quad N\_EXCESS (kg N ha$^{-1}$) & 200 & 200 & 100 \\
\midrule
\multicolumn{4}{l}{\textit{$R_{\text{penalty}}$ split (water / N)}} \\
\quad $\lambda_W : \lambda_N$ & 50:50 & 85:15 & 30:70 \\
\midrule
N-min yield floor on $R_{\text{ane}}$ & --- & --- & 5000 kg ha$^{-1}$ \\
\bottomrule
\end{tabular}
\end{table}

The Water-min variant reduces moisture-chasing and irrigation-excess tolerance;
the N-min variant strengthens N-application penalties and prioritises agronomic
N efficiency, with a yield floor preventing a zero-fertiliser exploit.


%=========================================================
\subsection{Implementation Constants and Normalisation}
%=========================================================

All sub-reward terms are clipped to $[-1,+1]$ or $[0,1]$ before weighting.
The combined daily signal is scaled by $\alpha=\texttt{DAILY\_SCALE}=0.01$.
Moisture-ratio thresholds are $\mathrm{SMR}_{\text{low}}=0.40$ and
$\mathrm{SMR}_{\text{high}}=0.80$. Seasonal yield, harvest index, and ANE are
normalised against \texttt{BASELINE\_YIELD}$=7620$,
\texttt{BASELINE\_HIAD}$=0.45$, and \texttt{BASELINE\_ANE}$=40$~kg grain per
kg~N. In the implemented Base seasonal reward, the positive terms use weights
$W_{\text{YIELD}}=0.40$, $W_{\text{HIAD}}=0.10$, $W_{\text{ANE}}=0.30$ and the
resource penalty uses $W_{\text{PENALTY}}=0.20$, grouped as
$R_{\text{season}} = W_{\text{YIELD}} R_{\text{yield}} + W_{\text{HIAD}} R_{\text{hiad}}
+ W_{\text{ANE}} R_{\text{ane}} - W_{\text{PENALTY}} R_{\text{penalty}}$.


%------------------------------------------------------------------------------
\section{Multi-Objective Learning Methodology}
\label{sec:algos}
%------------------------------------------------------------------------------

Three learner families implement the proposed solution: scalar PPO baselines
(fixed reward shaping), Preference-Conditioned PPO (PC-PPO), and
Concave-Augmented Pareto Q-Learning (CAPQL). Rule-based controllers provide
non-learning reference policies; their evaluation is deferred to Chapter~4.
Training hyperparameters (rollout length, learning rate, \texttt{VecNormalize})
are likewise reported in Chapter~4.

\subsection{Scalar PPO Baselines}

PPO~\cite{schulman2017ppo} maximises the clipped surrogate
\begin{equation}
\mathcal{L}^{\text{CLIP}}(\theta) =
\mathbb{E}_t\!\Big[\, \min\big( \rho_t(\theta)\,\hat{A}_t,\;
\operatorname{clip}(\rho_t(\theta),\,1-\epsilon,\,1+\epsilon)\,\hat{A}_t \big) \Big],
\label{eq:ppoclip}
\end{equation}
where $\hat{A}_t$ is the generalised advantage estimate~\cite{schulman2016gae}
and $\rho_t(\theta)=\pi_\theta(a_t\mid s_t)/\pi_{\theta_{\text{old}}}(a_t\mid s_t)$.
The scalar PPO agent uses the hierarchical reward of
REQ001--REQ002 (Equations~\eqref{eq:daily_reward}--\eqref{eq:transition_scalar}) with an
11-dimensional observation. Additional specialised variants reweight daily
components toward water saving or nitrogen saving (Table~\ref{tab:reward_variants}).
These agents provide rule-free reference policies against which the
multi-objective learners are compared in Chapter~4.
Algorithm~\ref{alg:scalar_ppo} summarises the on-policy training loop.

\begin{algorithm}[ht]
\caption{Scalar PPO training on \texttt{SmartFarmSoSEnv}}
\label{alg:scalar_ppo}
\begin{algorithmic}[1]
\Require Environment $\mathcal{E}$, total steps $N$, rollout length $L$, clip $\epsilon$, GAE $\lambda$, discount $\gamma$
\State Initialise policy $\pi_\theta$ and value network $V_\phi$; $\theta_{\text{old}} \leftarrow \theta$
\While{environment steps $< N$}
  \State Collect rollout $\mathcal{D}=\{(s_t,a_t,r_t)\}_{t=1}^{L}$ with $a_t \sim \pi_{\theta_{\text{old}}}(\cdot \mid s_t)$
  \State Compute returns $\hat{R}_t$ and advantages $\hat{A}_t$ using GAE($\gamma,\lambda$)
  \For{$k = 1$ to \texttt{n\_epochs}}
    \For{each minibatch $\mathcal{B} \subset \mathcal{D}$}
      \State $\rho_t(\theta) \leftarrow \pi_\theta(a_t \mid s_t)\,/\,\pi_{\theta_{\text{old}}}(a_t \mid s_t)$
      \State Update $\theta$ by ascending $\mathcal{L}^{\text{CLIP}}(\theta)$ on $\mathcal{B}$
      \State Update $\phi$ by minimising the value-function loss on $\mathcal{B}$
    \EndFor
  \EndFor
  \State $\theta_{\text{old}} \leftarrow \theta$
\EndWhile
\State \textbf{return} trained policy $\pi_\theta$
\end{algorithmic}
\end{algorithm}


\subsection{Preference-Conditioned PPO (PC-PPO)}
\label{subsec:pcppo}

PC-PPO is \emph{Preference-Conditioned}, not a separate constrained-optimisation
variant: the same PPO clipped objective is used, but the observation includes
$\boldsymbol{\omega}$ and the multi-reward vector is scalarised through
$\boldsymbol{\omega}^{\top}\mathbf{R}_t$ as in Equation~\eqref{eq:scalarized}.
Each episode samples a new preference $\boldsymbol{\omega}\sim\mathrm{Dirichlet}(1,1,1)$,
so one network learns a family of policies over the simplex. At evaluation,
the same model is queried at yield-, N-efficiency-, and water-efficiency-focused
corners $(1,0,0)$, $(0,1,0)$, $(0,0,1)$ and at a balanced profile
$(\tfrac{1}{3},\tfrac{1}{3},\tfrac{1}{3})$, matching the three-dimensional
terminal vector $\mathbf{R}_T=[R_{\text{yield}}, R_{\text{ane}}, R_{\text{water\_eff}}]^\top$
returned by \texttt{02\_smart\_farm\_env\_pcppo.py}.

PC-PPO uses the 14-dimensional observation
$[s_t;\boldsymbol{\omega}]$ with \texttt{VecNormalize} on observations and rewards.
Its main strength is stable on-policy learning on the costly DSSAT simulator while
still covering multiple farmer priorities with a single model.
Algorithm~\ref{alg:pcppo} details the episode-level preference sampling and
PPO update procedure.

\begin{algorithm}[ht]
\caption{Preference-Conditioned PPO (PC-PPO)}
\label{alg:pcppo}
\begin{algorithmic}[1]
\Require Environment $\mathcal{E}$, total steps $N$, rollout length $L$
\State Initialise $\pi_\theta$, $V_\phi$; $\theta_{\text{old}} \leftarrow \theta$
\While{environment steps $< N$}
  \State Sample $\boldsymbol{\omega} \sim \mathrm{Dirichlet}(\mathbf{1})$ for the upcoming episode
  \State Reset $\mathcal{E}$; $t \leftarrow 0$; \textbf{done} $\leftarrow$ \textbf{false}
  \State $\mathcal{D} \leftarrow \emptyset$
  \While{$|\mathcal{D}| < L$ \textbf{and not} \textbf{done}}
    \State $\tilde{s}_t \leftarrow [s_t;\,\boldsymbol{\omega}]$
    \State $a_t \sim \pi_{\theta_{\text{old}}}(\cdot \mid \tilde{s}_t)$
    \State Execute $(s_{t+1}, r_t, \textbf{done}) \leftarrow \mathcal{E}.\textsc{Step}(a_t)$
    \State Store $(\tilde{s}_t, a_t, r_t)$ in $\mathcal{D}$; $t \leftarrow t + 1$
  \EndWhile
  \State Compute GAE advantages $\{\hat{A}_t\}$ and returns $\{\hat{R}_t\}$ on $\mathcal{D}$
  \For{$k = 1$ to \texttt{n\_epochs}}
    \For{each minibatch $\mathcal{B} \subset \mathcal{D}$}
      \State Update $\theta, \phi$ with clipped PPO loss (Equation~\eqref{eq:ppoclip})
    \EndFor
  \EndFor
  \State $\theta_{\text{old}} \leftarrow \theta$
\EndWhile
\State \textbf{return} $\pi_\theta$
\end{algorithmic}
\end{algorithm}


\subsection{Concave-Augmented Pareto Q-Learning (CAPQL)}
\label{subsec:capql}

CAPQL~\cite{lu2023capql} takes an off-policy route based on TD3. A twin
preference-conditioned critic learns vector-valued returns, while the actor
outputs deterministic daily actions. The terminal target uses linear
scalarisation plus a concave augmentation term:
\begin{equation}
  y = R_{\text{daily}} + \mathbb{1}_{t=T}\!\left(
    \boldsymbol{\omega}^{\top}\mathbf{R}_T
    + \lambda \cdot \mathrm{mean}_i\big[\log(R_i + 1 + \varepsilon)\big]
  \right),
  \label{eq:capql_target}
\end{equation}
with $\lambda=0.5$ and $\varepsilon=0.1$ in our implementation. Hindsight
preference relabeling resamples $\boldsymbol{\omega}$ for minibatch
transitions, improving coverage of the objective space without rerunning
DSSAT episodes.

CAPQL is attractive when many preference values must be explored from limited
simulation budget: one network can be queried at any $\boldsymbol{\omega}$ after
training.
Algorithm~\ref{alg:capql} summarises the TD3 interaction loop, replay storage,
hindsight preference relabeling, and concave-augmented critic--actor updates.

\begin{algorithm}[ht]
\caption{Concave-Augmented Pareto Q-Learning (CAPQL)}
\label{alg:capql}
\begin{algorithmic}[1]
\Require Environment $\mathcal{E}$, replay buffer $\mathcal{B}$, total steps $N$, batch size $B$, policy delay $d$
\State Initialise actor $\mu_\theta$, critics $Q_{\phi_1}, Q_{\phi_2}$, and target networks $\mu_{\bar\theta}, Q_{\bar\phi}$
\For{$n = 1$ to $N$}
  \If{episode start}
    \State Sample $\boldsymbol{\omega} \sim \mathrm{Dirichlet}(\mathbf{1})$; reset $\mathcal{E}$
  \EndIf
  \State $\tilde{s}_t \leftarrow [s_t;\,\boldsymbol{\omega}]$
  \State $a_t \leftarrow \mu_\theta(\tilde{s}_t) + \text{exploration noise}$; step $\mathcal{E}$
  \State Store $(s_t, a_t, R_{\text{daily},t}, \mathbf{R}_T\,\mathbb{1}_{t=T}, s_{t+1}, \textbf{done}, \boldsymbol{\omega})$ in $\mathcal{B}$
  \If{$|\mathcal{B}| \ge B$ \textbf{and} $n > \texttt{warmup}$}
    \State Sample minibatch; relabel $\boldsymbol{\omega}' \sim \mathrm{Dirichlet}(\mathbf{1})$ per transition
    \State $\tilde{s} \leftarrow [s;\boldsymbol{\omega}']$,\; $\tilde{s}' \leftarrow [s';\boldsymbol{\omega}']$
    \State $r_{\text{aug}} \leftarrow R_{\text{daily}} + \textbf{done}\cdot\big(\boldsymbol{\omega}'^{\top}\mathbf{R}_T + \lambda\,\mathrm{mean}_i\log(R_i{+}1{+}\varepsilon)\big)$
    \State $a' \leftarrow \mathrm{clip}\big(\mu_{\bar\theta}(\tilde{s}') + \text{smoothing noise}\big)$
    \State $y \leftarrow r_{\text{aug}} + \gamma(1-\textbf{done})\min_{j \in \{1,2\}} Q_{\bar\phi_j}(\tilde{s}', a')$
    \State Update $\phi_1, \phi_2$ by minimising $\|Q_{\phi_j}(\tilde{s},a) - y\|^2$
    \If{$n \bmod d = 0$}
      \State Update $\theta$ by ascending $\mathbb{E}\big[\min_j Q_{\phi_j}(\tilde{s}, \mu_\theta(\tilde{s}))\big]$
      \State Soft-update $\bar\theta \leftarrow \tau\theta + (1{-}\tau)\bar\theta$;\;
            $\bar\phi \leftarrow \tau\phi + (1{-}\tau)\bar\phi$
    \EndIf
  \EndIf
\EndFor
\State \textbf{return} $\mu_\theta$
\end{algorithmic}
\end{algorithm}


\subsection{Justification of Algorithmic Choices}

\begin{table}[ht]
\centering
\caption{Rationale for the selected learning algorithms.}
\label{tab:algos}
\small
\renewcommand{\arraystretch}{1.3}
\begin{tabular}{>{\raggedright\arraybackslash}p{2.8cm}
                >{\raggedright\arraybackslash}p{2.6cm}
                >{\raggedright\arraybackslash}p{6.3cm}}
\toprule
\textbf{Method} & \textbf{Family} & \textbf{Why chosen here} \\
\midrule
Scalar PPO & On-policy, single-objective &
Stable baseline on \texttt{gym-DSSAT}; hierarchical reward matches REQ000--REQ002;
provides fixed-reward reference policies for comparison. \\
PC-PPO & On-policy, preference-conditioned &
Same PPO core as the baseline; one model serves all preference corners;
observation carries $\boldsymbol{\omega}$; strong training stability. \\
CAPQL & Off-policy, preference-conditioned &
Sample-efficient exploration of the Pareto set; concave augmentation mitigates
plain linear scalarisation; TD3 twin critics reduce Q overestimation. \\
\bottomrule
\end{tabular}
\end{table}

In short, scalar PPO defines what fixed hierarchical rewards can encode
(Algorithm~\ref{alg:scalar_ppo}); PC-PPO extends the same PPO core to a
preference-conditioned policy family (Algorithm~\ref{alg:pcppo}); CAPQL
complements this with an off-policy Pareto-learning formulation
(Algorithm~\ref{alg:capql}).


%------------------------------------------------------------------------------
\section{Integrated Proposed Architecture}
\label{sec:overall}
%------------------------------------------------------------------------------

Figure~\ref{fig:overall} assembles the elements of this chapter into the complete
closed-loop learning system. The DSSAT core, wrapped by \texttt{gym-DSSAT} and
\texttt{SmartFarmSoSEnv}, exposes the daily crop state. The reward module
converts DSSAT outputs into daily and seasonal signals (Section~\ref{sec:reward}).
The agent---scalar PPO, PC-PPO, or CAPQL---consumes state and reward and emits
irrigation and nitrogen actions, closing the loop over one maize season.

%---- FIGURE : Overall proposed model ----------------------------------------
\begin{figure}[ht]
\centering
\resizebox{\textwidth}{!}{%
\begin{tikzpicture}[node distance=1.0cm and 1.8cm]

  \node[core, draw=sosblue, fill=sosblue!10, minimum width=3.4cm, minimum height=1.6cm]
       (env) {gym-DSSAT\\Environment\\\scriptsize (DSSAT CSM core)};

  \node[core, draw=sosgreen, fill=sosgreen!10, below=2.2cm of env,
        minimum width=3.4cm, minimum height=1.6cm] (reward)
        {Reward Module\\\scriptsize $R_{\text{daily}}+R_{\text{season}}$};

  \node[core, draw=sosorange, fill=soscream, right=4.2cm of env,
        minimum width=3.6cm, minimum height=1.6cm] (agent)
        {Multi-Objective RL Agent\\\scriptsize PPO $\,\cdot\,$ PC-PPO $\,\cdot\,$ CAPQL};

  \node[block, draw=sosgreen, fill=sosgreen!8, below=2.2cm of agent,
        minimum width=3.6cm, minimum height=1.5cm, font=\scriptsize, align=center]
        (vec) {Seasonal vector $\mathbf{R}_T$\\yield $\cdot$ N-efficiency\\water efficiency};

  \draw[flow] (env.east) -- node[above,font=\scriptsize]{state $s_t$} (agent.west);
  \draw[ctrl] (agent.south) |- node[below right,font=\scriptsize,pos=0.25]
        {action $a_t=(\mathrm{FERN}_t,\mathrm{IRRC}_t)$} (env.south);
  \draw[flow] (env.south) -- node[left,font=\scriptsize]{DSSAT outputs} (reward.north);
  \draw[flow] (reward.east) -- (vec.west);
  \draw[flow] (vec.north) -- node[right,font=\scriptsize]{$\mathbf{R}_T$} (agent.south);

  \node[block, draw=sosgrey, fill=white, right=1.2cm of agent, font=\scriptsize,
        align=center, minimum width=2.6cm] (eval)
        {Evaluation\\\scriptsize vs.\ rule-based\\and scalar PPO baselines};
  \draw[flow, sosgrey] (agent.east) -- (eval.west);

  \begin{scope}[on background layer]
    \node[draw=sosgrey, dashed, rounded corners, inner sep=8pt,
          fit=(env)(reward)(agent)(vec),
          label={[sosgrey,font=\scriptsize\bfseries]above:Closed RL loop (one growing season)}] {};
  \end{scope}

\end{tikzpicture}%
}
\caption{Overall proposed learning architecture. The DSSAT core supplies daily
  crop state; the reward module computes hierarchical daily and seasonal
  signals; the agent issues irrigation and nitrogen actions, closing the loop
  over one maize season.}
\label{fig:overall}
\end{figure}


\subsection{Traceability from Requirements to Implementation}

Table~\ref{tab:trace} summarises how REQ000--REQ002 connect mission objectives,
MoEs, reward terms, and learner families. This traceability matrix is the
design counterpart to the performance tables reported in Chapter~4.

\begin{table}[ht]
\centering
\caption{Traceability: requirements, MoEs, rewards, and learners.}
\label{tab:trace}
\small
\renewcommand{\arraystretch}{1.2}
\begin{tabular}{>{\raggedright\arraybackslash}p{1.6cm}
                >{\raggedright\arraybackslash}p{2.8cm}
                >{\raggedright\arraybackslash}p{2.6cm}
                >{\raggedright\arraybackslash}p{2.4cm}
                >{\raggedright\arraybackslash}p{2.2cm}}
\toprule
\textbf{Req.} & \textbf{MoE / objective} & \textbf{Reward term} & \textbf{Module} & \textbf{Learner} \\
\midrule
REQ000 & Yield, N, water & $R_{\text{daily}}+R_{\text{season}}$ &
\texttt{02\_smart\_farm\_env*} & PPO, PC-PPO, CAPQL \\
REQ001 & SMR, SWFAC, NSTRES & $R_{\text{water}}, R_{\text{fert}}$ &
\texttt{02\_smart\_farm\_env.py} & All \\
REQ001 & Resource use, losses & $R_{\text{resource}}, R_{\text{losses}}$ &
\texttt{02\_smart\_farm\_env.py} & All \\
REQ002 & GRNWT, HIAD, ANE & $R_{\text{yield}}, R_{\text{hiad}}, R_{\text{ane}}$ &
\texttt{02\_smart\_farm\_env.py} & PPO \\
REQ002 & Season totals & $R_{\text{penalty}}$ &
\texttt{02\_smart\_farm\_env*.py} & PPO variants \\
MORL   & Preference trade-offs & $\mathbf{R}_T$, $\boldsymbol{\omega}$ &
\texttt{02\_smart\_farm\_env\_pcppo.py} & PC-PPO, CAPQL \\
\bottomrule
\end{tabular}
\end{table}


\subsection{Comparison Baselines (Design)}

For evaluation (Chapter~4), four rule-based controllers run inside the same
\texttt{SmartFarmSoSEnv} wrapper: \emph{Random} (unconstrained daily inputs),
\emph{Fixed} (preset N splits and periodic irrigation), \emph{Stage-based}
(management keyed to vegetative development), and \emph{FAO-56} (soil-water-balance
irrigation with a three-split N plan). Using a common environment ensures that
MoEs and reward components are directly comparable across learned and rule-based
policies.


%------------------------------------------------------------------------------
\section{Conclusion}
\label{sec:ch3_conclusion}
%------------------------------------------------------------------------------

This chapter presented the proposed solution end to end: a smart-farming
System-of-Systems formalised around the DSSAT Cropping System Model; an MDP
formulation with explicit observation, action, and software layers;
Measures of Effectiveness tied to mission objectives (Table~\ref{tab:moe});
a hierarchical reward design (REQ000--REQ002) with scalar variants
(Table~\ref{tab:reward_variants}) and a multi-objective extension for
preference-conditioned learning; and three learner families---scalar PPO,
PC-PPO, and CAPQL---integrated in the closed-loop architecture of
Figure~\ref{fig:overall}. The next chapter reports the experimental setup,
training dynamics, and quantitative comparison of all learned and rule-based
policies.


% =============================================================================
% If using BibTeX instead, delete this block and add entries to references.bib
% =============================================================================
\begin{thebibliography}{99}

\bibitem{maier1998}
M.~W. Maier,
``Architecting principles for systems-of-systems,''
\emph{Systems Engineering}, vol.~1, no.~4, pp.~267--284, 1998.

\bibitem{jones2003dssat}
J.~W. Jones, G.~Hoogenboom, C.~H. Porter, K.~J. Boote, W.~D. Batchelor, L.~A. Hunt,
P.~W. Wilkens, U.~Singh, A.~J. Gijsman, and J.~T. Ritchie,
``The DSSAT cropping system model,''
\emph{European Journal of Agronomy}, vol.~18, no.~3--4, pp.~235--265, 2003.

\bibitem{hoogenboom2019}
G.~Hoogenboom, C.~H. Porter, K.~J. Boote, V.~Shelia, P.~W. Wilkens, U.~Singh,
J.~W. White, S.~Asseng, J.~I. Lizaso, et~al.,
``The DSSAT crop modeling ecosystem,'' in
\emph{Advances in Crop Modelling for a Sustainable Agriculture}, Burleigh Dodds Science
Publishing, pp.~173--216, 2019.

\bibitem{gautron2022gymdssat}
R.~Gautron, E.~J. Padr\'on, P.~Preux, J.~Bigot, O.-A. Maillard, and D.~Emukpere,
``gym-DSSAT: a crop model turned into a Reinforcement Learning environment,''
Inria Research Report RR-9460, 2022.

\bibitem{gautron2022review}
R.~Gautron, O.-A. Maillard, P.~Preux, M.~Corbeels, and R.~Sabbadin,
``Reinforcement learning for crop management support: Review, prospects and challenges,''
\emph{Computers and Electronics in Agriculture}, vol.~200, art.~107182, 2022.

\bibitem{schulman2017ppo}
J.~Schulman, F.~Wolski, P.~Dhariwal, A.~Radford, and O.~Klimov,
``Proximal policy optimization algorithms,''
\texttt{arXiv:1707.06347}, 2017.

\bibitem{schulman2016gae}
J.~Schulman, P.~Moritz, S.~Levine, M.~Jordan, and P.~Abbeel,
``High-dimensional continuous control using generalized advantage estimation,'' in
\emph{Proc.\ ICLR}, 2016.

\bibitem{lu2023capql}
H.~Lu, D.~Herman, and Y.~Yu,
``Multi-objective reinforcement learning: Convexity, stationarity and Pareto optimality,''
in \emph{Proc.\ ICLR}, 2023.

\end{thebibliography}

% Required: amsmath, booktabs, graphicx, tikz, backgrounds, algorithm, algpseudocode
%   \usetikzlibrary{positioning, arrows.meta, fit, calc, backgrounds}
% =============================================================================

%---- Colour palette for the diagrams ----------------------------------------
\definecolor{sosblue}{RGB}{38,84,124}
\definecolor{sosgreen}{RGB}{61,133,92}
\definecolor{sosorange}{RGB}{200,120,40}
\definecolor{sosgrey}{RGB}{90,90,90}
\definecolor{soslight}{RGB}{234,240,245}
\definecolor{soscream}{RGB}{248,244,232}

%---- Reusable TikZ node styles ----------------------------------------------
\tikzset{
  sensor/.style   = {rectangle, rounded corners=2pt, draw=sosblue, very thick,
                     fill=soslight, align=center, minimum width=2.5cm,
                     minimum height=1cm, font=\small},
  actuator/.style = {rectangle, rounded corners=2pt, draw=sosorange, very thick,
                     fill=soscream, align=center, minimum width=2.5cm,
                     minimum height=1cm, font=\small},
  core/.style     = {rectangle, rounded corners=4pt, draw=sosgreen, line width=1.2pt,
                     fill=sosgreen!12, align=center, minimum width=3.2cm,
                     minimum height=1.3cm, font=\small\bfseries},
  mission/.style  = {ellipse, draw=sosgrey, very thick, fill=white, align=center,
                     font=\small\itshape, minimum width=3cm},
  flow/.style     = {-{Stealth[length=2.5mm]}, line width=0.9pt, sosgrey},
  ctrl/.style     = {-{Stealth[length=2.5mm]}, line width=0.9pt, sosorange, dashed},
  block/.style    = {rectangle, rounded corners=2pt, draw=sosblue, thick,
                     fill=soslight, align=center, font=\small,
                     minimum width=2.6cm, minimum height=0.9cm},
}


\chapter{Proposed Solution}
\label{chap:proposed}


%------------------------------------------------------------------------------
\section{Introduction}
\label{sec:ch3_intro}
%------------------------------------------------------------------------------

Chapter~\ref{ch:background} positioned DSSAT as the crop-simulation subsystem
inside a smart-farming System-of-Systems. This chapter turns that view into a
concrete decision framework.

The problem is no longer a single irrigation rule. It is a \emph{multi-objective}
management task in which yield, water use, and nitrogen efficiency must be
balanced over a full growing season. The proposed framework combines three
elements:
\begin{enumerate}
  \item a DSSAT-based daily simulation environment (\texttt{gym-DSSAT});
  \item a hierarchical reward design with preference conditioning for
        multi-objective learning;
  \item learning algorithms---scalar PPO baselines, Preference-Conditioned PPO
        (PC-PPO), and Concave-Augmented Pareto Q-Learning (CAPQL).
\end{enumerate}

The chapter is organised as follows. Section~\ref{sec:sos} formalises the
control problem and the smart-farming System-of-Systems architecture.
Section~\ref{sec:env} describes the simulation platform, observation and action
spaces, and software stack. Section~\ref{sec:reward} presents the hierarchical
reward design and scalar reward variants. Section~\ref{sec:algos} defines the
learning algorithms. Section~\ref{sec:overall} assembles the integrated
architecture and traceability from requirements to implementation. Experimental
training protocols and numerical results are reported in the next chapter.


%------------------------------------------------------------------------------
\section{Problem Formulation and System-of-Systems Architecture}
\label{sec:sos}
%------------------------------------------------------------------------------

\subsection{Overview and Rationale}

A modern smart farm is not a single monolithic system but a \emph{System of
Systems}: a collection of heterogeneous, independently operated subsystems---
weather data, soil and crop observations, irrigation and fertilisation
controllers, and crop-modelling platforms---integrated to pursue missions such
as maximising yield, conserving water, and improving nitrogen-use efficiency.
Following Maier~\cite{maier1998}, a true SoS is distinguished from a merely
complex monolithic system by five traits: operational independence, managerial
independence, geographic distribution, emergent behaviour, and evolutionary
development. The smart farm exhibits all five, which is precisely what makes
mission-level decision-making difficult: an irrigation pulse today changes soil
moisture, crop stress, and nitrogen dynamics over the rest of the season.
Section~\ref{sec:env} formalises this task as a Markov decision process; here
we focus on the architectural view.


\subsection{Components of the Smart-Farming SoS}

The constituent systems are grouped by their role in the closed decision loop:
observe field conditions, act on the crop, and decide daily management. At the
centre of this loop sits the \textbf{crop-modelling platform (DSSAT)}, which
integrates weather, soil, and management inputs and produces crop growth, stress
indicators, and environmental fluxes. Monitoring systems read from it; control
systems write into it; the decision layer issues the commands that drive both.

\vspace{0.25cm}
\noindent\textbf{Monitoring systems.}
These constituents supply the observation stream the agent uses each day:
\begin{itemize}
  \item \textbf{Weather data} --- supplies exogenous drivers (temperature,
        rainfall, solar radiation) that the agent observes but cannot control.
  \item \textbf{Soil and crop state} --- reported through DSSAT variables such
        as soil-water content, water-stress factor, nitrogen-stress factor,
        leaf area, and grain weight.
\end{itemize}

\vspace{0.25cm}
\noindent\textbf{Control systems.}
These constituents execute management actions issued by the decision layer:
\begin{itemize}
  \item \textbf{Irrigation controller} --- the actuator through which the agent
        applies a daily water amount $a_t^{W}$ (mm/day).
  \item \textbf{Fertilisation unit} --- the actuator through which the agent
        applies a daily nitrogen amount $a_t^{N}$ (kg N/ha/day).
\end{itemize}

\vspace{0.25cm}
\noindent\textbf{Decision and mission layer.}
Mission objectives---yield, water efficiency, and nitrogen efficiency---condition
the agent through the preference vector $\boldsymbol{\omega}$. The decision
constituent that closes the loop is:
\begin{itemize}
  \item \textbf{Mission manager / RL agent} --- reads observations, respects the
        farmer's preference vector, and issues daily irrigation and fertilisation
        commands to the control systems.
\end{itemize}

Figure~\ref{fig:sos} depicts the architecture. Monitoring systems feed
observations into the mission manager (solid arrows); the mission manager issues
control commands to the actuators (dashed arrows); DSSAT closes the loop by
producing the next crop and soil state.


%---- FIGURE : SoS architecture ------------------------------------------------
\begin{figure}[ht]
\centering
\includegraphics[width=\textwidth]{sos_architecture.png}
\caption{Architecture of the proposed smart-farming System-of-Systems. Icons denote
constituent roles; solid arrows carry observations and simulated state, dashed
orange arrows carry management commands. DSSAT integrates inputs and closes the
daily decision loop.}
\label{fig:sos}
\end{figure}


\subsection{System-of-Systems Characteristics in Our Setting}

Mapping Maier's five traits~\cite{maier1998} onto our design clarifies why a
learning-based formulation is appropriate. \emph{Operational independence}
holds because weather data, DSSAT, and actuators can each be described on
their own. \emph{Managerial independence} holds because these parts are
developed and maintained separately yet must cooperate in practice.
\emph{Geographic distribution} is intrinsic to field deployment.
\emph{Emergent behaviour}---yield and resource efficiency arising from the
interaction of weather, soil, crop, and management---is exactly what the agent
must learn to shape. \emph{Evolutionary development} motivates a framework in
which the crop model, reward design, or learning algorithm can be upgraded
without rebuilding the entire platform.


\subsection{Mission Objectives and Measures of Effectiveness}
\label{sec:moe}

Measures of Effectiveness (MoEs) translate mission objectives into quantities
that can be logged in simulation and encoded in the reward function.
Table~\ref{tab:moe} links each MoE to its mission objective and to the reward
term that implements it. \emph{Seasonal} MoEs evaluate end-of-season success;
\emph{daily} indicators supply the dense feedback used in reward shaping
(Section~\ref{sec:reward}).

\begin{table}[ht]
\centering
\caption{Measures of Effectiveness and their alignment with mission objectives.}
\label{tab:moe}
\small
\renewcommand{\arraystretch}{1.25}
\begin{tabular}{>{\raggedright\arraybackslash}p{4.4cm}
                >{\raggedright\arraybackslash}p{3.2cm}
                >{\raggedright\arraybackslash}p{4.6cm}}
\toprule
\textbf{Measure of Effectiveness} & \textbf{Mission objective} & \textbf{Indicator / reward link} \\
\midrule
Final grain yield ($\mathrm{GRNWT}_{\text{final}}$) & Yield optimisation & $w_y\,\mathrm{GRNWT}_{\text{final}}$ \\
Harvest index ($\mathrm{HIAD}_{\text{final}}$) & Yield quality & $w_h\,\mathrm{HIAD}_{\text{final}}$ \\
Agronomic N efficiency ($\mathrm{ANE}_{\text{final}}$) & N-use efficiency & $w_a\,\mathrm{ANE}_{\text{final}}$ \\
Total irrigation ($\mathrm{IRRC}_{\text{season}}$) & Water-use efficiency & $w_w\,\mathrm{IRRC}_{\text{season}}$ \\
Total nitrogen applied ($\mathrm{FERN}_{\text{season}}$) & N-use efficiency & $w_f\,\mathrm{FERN}_{\text{season}}$ \\
Soil-moisture balance ($\mathrm{SMR}_t$) & Crop health / water stress & $R_{\text{water},t}$ \\
Water-stress factor ($\mathrm{SWFAC}_t$) & Crop health & $R_{\text{water},t}$ \\
Nitrogen-stress factor ($\mathrm{NSTRES}_t$) & Crop health & $R_{\text{fert},t}$ \\
Cumulative N \& water losses & Environmental impact & $R_{\text{losses},t}$ \\
\bottomrule
\end{tabular}
\end{table}


%------------------------------------------------------------------------------
\section{Simulation Platform and Environment Design}
\label{sec:env}
%------------------------------------------------------------------------------

\subsection{The DSSAT Cropping System Model}

The Decision Support System for Agrotechnology Transfer (DSSAT) is a suite of
mechanistic crop models developed over more than three decades and widely
validated by the agronomic community~\cite{jones2003dssat,hoogenboom2019}. At
its heart is the \emph{Cropping System Model} (CSM), which advances the state of
a land unit on a daily time step. Given weather, soil, genetics, and management
inputs, the CSM couples soil-water balance, soil--plant--atmosphere exchange,
and plant growth modules. For maize, this work uses CERES-Maize. Its outputs
include grain yield, biomass, harvest index, and water and nitrogen fluxes such
as runoff, leaching, and denitrification. Figure~\ref{fig:dssat} sketches the
modular structure.


\begin{figure}[ht]
\centering
\begin{tikzpicture}[node distance=0.7cm and 1.2cm]
  \node[core, fill=sosblue!10, draw=sosblue] (w) {Weather};
  \node[core, fill=sosblue!10, draw=sosblue, below=of w] (s) {Soil profile};
  \node[core, fill=sosgreen!15, draw=sosgreen, right=2.2cm of w, yshift=-0.4cm] (csm)
        {DSSAT-CSM\\daily integrator};
  \node[core, fill=soscream, draw=sosorange, right=2.2cm of csm, yshift=0.9cm] (crop)
        {CERES-Maize};
  \node[core, fill=soscream, draw=sosorange, right=2.2cm of csm, yshift=-0.9cm] (mgmt)
        {Management\\(N + irrigation)};
  \node[block, below=1.6cm of csm, minimum width=4.2cm] (out)
        {Outputs: yield, stress, losses, state vector};

  \draw[flow] (w) -- (csm);
  \draw[flow] (s) -- (csm);
  \draw[flow] (crop) -- (csm);
  \draw[flow] (mgmt) -- (csm);
  \draw[flow] (csm) -- (out);
\end{tikzpicture}
\caption{Schematic of the DSSAT Cropping System Model used in this work.}
\label{fig:dssat}
\end{figure}


\subsection{From DSSAT to an RL Environment: gym-DSSAT}

To use DSSAT for reinforcement learning, the simulator must expose a step-wise
interface in which an agent observes a state, applies an action, and receives
the resulting next state and reward. We rely on \texttt{gym-DSSAT}~\cite{gautron2022gymdssat},
which wraps the Fortran simulator in a Gymnasium-compatible API allowing daily
interaction throughout the growing season~\cite{gautron2022review}. A
\texttt{SmartFarmSoSEnv} layer maps DSSAT outputs to the observation vector,
translates the agent's action into DSSAT management events
(\texttt{anfer}, \texttt{amir}), and computes the reward defined in
Section~\ref{sec:reward}.


\subsection{Decision-Process Formulation}

Formally, the control task is a Markov Decision Process
$\langle \mathcal{S}, \mathcal{A}, P, R, \gamma \rangle$. In the
multi-objective setting, the terminal step also exposes a reward vector
$\mathbf{R}_T$ that is scalarised through a preference vector
$\boldsymbol{\omega}$ (Section~\ref{sec:reward}).

\begin{itemize}
  \item \textbf{State / observation $s_t$:} an 11-dimensional vector of
        normalised crop and SoS features: biomass, grain weight, days after
        planting, vegetative stage, leaf area index, cumulative nitrogen and
        irrigation, mean soil water, energy budget, communication quality, and
        fraction of healthy virtual sensors. PC-PPO and CAPQL append the
        3-dimensional preference vector, giving a 14-dimensional observation.
  \item \textbf{Action $a_t$:} the daily management decision
        $a_t = (\mathrm{FERN}_t, \mathrm{IRRC}_t) = (a_t^{N}, a_t^{W})$ with
        $a_t^{N}\in[0,200]$~kg N/ha/day and $a_t^{W}\in[0,50]$~mm/day.
  \item \textbf{Transition $P(s_{t+1}\mid s_t,a_t)$:} realised by the DSSAT
        daily integration; stochasticity enters through random weather resampling
        across episodes.
  \item \textbf{Horizon:} one growing season (${\approx}161$ days); the episode
        terminates at harvest, when the seasonal reward is granted.
\end{itemize}

Irrigation and fertilisation are coupled because both alter the same soil-water
and nitrogen pools that drive crop stress and final yield. A sequential
formulation is therefore more faithful than optimising either input in isolation.


\subsection{Observation and Action Spaces}

Table~\ref{tab:obs} lists the 11-dimensional observation vector exposed to
scalar PPO agents through the \texttt{pc\_env.py} adapter. Each component is
linearly scaled to $[0,1]$ from agronomically meaningful bounds. PC-PPO and
CAPQL append the three-dimensional preference vector
$\boldsymbol{\omega}=(\omega_y,\omega_a,\omega_w)$, yielding a
14-dimensional input.

\begin{table}[ht]
\centering
\caption{Observation vector for scalar PPO (\texttt{pc\_env.py}).}
\label{tab:obs}
\small
\renewcommand{\arraystretch}{1.2}
\begin{tabular}{clll}
\toprule
\textbf{Idx} & \textbf{DSSAT / SoS variable} & \textbf{Scale} & \textbf{Role} \\
\midrule
0  & \texttt{topwt}  & $[0, 20\,000]$ kg ha$^{-1}$ & Above-ground biomass \\
1  & \texttt{grnwt}  & $[0, 12\,000]$ kg ha$^{-1}$ & Grain weight \\
2  & \texttt{dap}    & $[0, 200]$ days             & Days after planting \\
3  & \texttt{vstage} & $[0, 18]$                   & Vegetative stage \\
4  & \texttt{xlai}   & $[0, 7]$                    & Leaf area index \\
5  & Cumulative N    & $[0, 300]$ kg ha$^{-1}$     & Seasonal fertiliser \\
6  & Cumulative W    & $[0, 1000]$ mm              & Seasonal irrigation \\
7  & Mean soil water & $[0, 1]$                    & Profile moisture \\
8  & Energy budget   & $[0, 100]$                  & SoS energy state \\
9  & Comm.\ quality  & $[0, 1]$                    & Link reliability \\
10 & Sensors alive   & $[0, 1]$                    & Virtual sensor health \\
\bottomrule
\end{tabular}
\end{table}

The daily action vector maps to DSSAT management events:
$a_t=(a_t^{N},a_t^{W})=(\texttt{anfer},\texttt{amir})$ with
$a_t^{N}\in[0,200]$~kg N ha$^{-1}$ day$^{-1}$ and
$a_t^{W}\in[0,50]$~mm day$^{-1}$.


\subsection{Episode Protocol and Calibration Anchors}

Each episode simulates one maize growing season of approximately 161 days.
Weather is resampled randomly across episodes (\texttt{random\_weather=True}),
injecting the main source of transition stochasticity. Yield calibration is
anchored to the stage-based baseline at ${\approx}7\,620$~kg ha$^{-1}$
(\texttt{BASELINE\_YIELD}); normalisation upper bounds follow
\texttt{MAX\_EXPECTED\_YIELD}$=1.5\times$ baseline. Fault injection
(\texttt{enable\_faults=False} in the main experiments) is disabled so that
learning focuses on crop-management trade-offs; resilience extensions are left
for future work.


\subsection{Software Architecture and Module Map}

The implementation stacks four layers: the DSSAT core accessed through
\texttt{gym-DSSAT}; a \texttt{SmartFarmSoSEnv} reward wrapper; a
Gymnasium-compatible adapter; and the training script. Table~\ref{tab:software}
maps each module to its role in the proposed solution.

\begin{table}[ht]
\centering
\caption{Software modules implementing the proposed solution.}
\label{tab:software}
\small
\renewcommand{\arraystretch}{1.25}
\begin{tabular}{>{\raggedright\arraybackslash}p{3.6cm}
                >{\raggedright\arraybackslash}p{4.2cm}
                >{\raggedright\arraybackslash}p{5.5cm}}
\toprule
\textbf{Layer} & \textbf{Module} & \textbf{Function} \\
\midrule
DSSAT core      & \texttt{gym\_dssat\_pdi}           & Daily CSM integration via subprocess API \\
Scalar reward   & \texttt{02\_smart\_farm\_env.py}   & Hierarchical scalar reward (REQ001--002) \\
Water-min reward& \texttt{02\_smart\_farm\_env\_water.py} & Reward weights tuned for irrigation saving \\
N-min reward    & \texttt{02\_smart\_farm\_env\_nitrogen.py} & Reward weights tuned for N saving \\
MORL reward     & \texttt{02\_smart\_farm\_env\_pcppo.py} & Terminal objective vector $\mathbf{R}_T$ \\
11-dim adapter  & \texttt{pc\_env.py}                & Box observation for scalar PPO \\
14-dim adapter  & \texttt{pc\_env\_pcppo.py}         & Appends $\boldsymbol{\omega}$ for PC-PPO \\
CAPQL adapter   & \texttt{capql\_env.py}             & PyTorch interface for TD3 training \\
Scalar PPO train& \texttt{04\_pc\_ppo\_quick\_train.py}, \texttt{04b}, \texttt{04c} & Base, Water-min, N-min \\
PC-PPO train    & \texttt{05\_pcppo\_train.py}       & Preference-conditioned PPO \\
CAPQL train     & \texttt{06\_capql\_train.py}       & Concave-augmented Pareto Q-learning \\
Baselines       & \texttt{01\_rule\_based\_baselines.py} & Random, Fixed, Stage, FAO-56 controllers \\
\bottomrule
\end{tabular}
\end{table}


%------------------------------------------------------------------------------
\section{Reward Engineering}
\label{sec:reward}
%------------------------------------------------------------------------------

The reward design is the heart of the proposed solution: it is where the mission
objectives are encoded into learnable feedback. We adopt a \emph{hierarchical}
structure that combines dense daily rewards with a dominant terminal seasonal
reward. The design is presented below as a set of traceable requirements
(REQ000--REQ002), which keeps an explicit link between each reward term and the
mission objective and MoE it serves.

%=========================================================
\subsection{REQ000: Optimise Irrigation and Fertilisation Performance}
%=========================================================

\noindent\textbf{Detailed Requirement.}
The smart-farming system shall optimise irrigation and nitrogen-fertilisation strategies
throughout the crop growing season in order to maximise crop productivity while
minimising water and fertiliser consumption.

\vspace{0.2cm}
\noindent\textbf{Rationale.}
Efficient agricultural management requires balancing productivity and sustainability.
Excessive irrigation wastes water and may increase runoff and drainage losses, while
excessive fertilisation can increase nutrient leaching and environmental impacts.
Therefore, the reinforcement-learning agent must learn resource-efficient management
policies that maintain healthy crop conditions and maximise final crop yield.

\vspace{0.2cm}
\noindent\textbf{Decomposition into Adaptive Requirements.}
To operationalise this objective, REQ000 is decomposed into the following adaptive
requirements:
\begin{itemize}
    \item \textbf{REQ001:} Optimise Daily Reward ($R_{\text{daily}}$).
    \item \textbf{REQ002:} Optimise Seasonal Reward ($R_{\text{season}}$).
\end{itemize}

\vspace{0.2cm}
\noindent\textbf{Associated Measure of Effectiveness.}
\begin{itemize}
    \item Final grain yield.
    \item Total irrigation applied.
    \item Total nitrogen fertiliser applied.
    \item Agronomic Nitrogen Efficiency (ANE).
    \item Harvest Index (HIAD).
\end{itemize}

\vspace{0.2cm}
\noindent\textbf{Traceability.}
REQ001 and REQ002 derive directly from REQ000 by providing short-term and long-term
feedback mechanisms that guide the agent toward sustainable crop management.

%=========================================================
\subsection{REQ001: Optimise Daily Reward ($R_{\text{daily}}$)}
%=========================================================

\noindent\textbf{Detailed Requirement.}
At each simulation day $t$, the system shall compute an intermediate reward that
encourages healthy crop-growth conditions, efficient resource utilisation, and reduced
environmental losses.

\vspace{0.2cm}
\noindent\textbf{Rationale.}
The final crop yield is only observed at harvest. Relying exclusively on harvest
feedback creates a sparse-reward problem that slows learning. Therefore, daily rewards
provide dense feedback that guides the agent toward agronomically desirable decisions
throughout the season.

\vspace{0.2cm}
\noindent\textbf{Associated Measure of Effectiveness.}
\begin{itemize}
    \item Soil-moisture balance.
    \item Water-stress reduction.
    \item Nitrogen-stress reduction.
    \item Nitrogen-uptake promotion.
    \item Reduced irrigation and fertiliser usage.
    \item Reduced runoff, drainage, and nutrient losses.
\end{itemize}

\vspace{0.2cm}
\noindent\textbf{Operational Definition and Indicators.}
\begin{itemize}
  \item \textbf{Water Management Indicators}
  \begin{itemize}
      \item $\mathrm{SWFAC}_t$: Water-stress factor.
      \item $\mathrm{SMR}_t$: Soil-moisture ratio.
  \end{itemize}
  \begin{equation}
  \mathrm{SMR}_t=\frac{\mathrm{SWTD}_t}{\mathrm{SWXD}_t}
  \end{equation}
  where:
  \begin{itemize}
      \item $\mathrm{SWTD}_t$: Total soil water.
      \item $\mathrm{SWXD}_t$: Available water capacity.
  \end{itemize}

  \item \textbf{Nitrogen Management Indicators}
  \begin{itemize}
      \item $\mathrm{NSTRES}_t$: Nitrogen-stress factor.
      \item $\mathrm{TRNU}_t$: Daily nitrogen uptake.
  \end{itemize}

  \item \textbf{Resource Consumption Indicators}
  \begin{itemize}
      \item $\mathrm{IRRC}_t$: Daily irrigation amount.
      \item $\mathrm{FERN}_t$: Daily nitrogen-fertiliser amount.
  \end{itemize}

  \item \textbf{Environmental Loss Indicators}
  \begin{itemize}
      \item $\mathrm{ROFC}_t$: Surface runoff.
      \item $\mathrm{DRNC}_t$: Deep drainage.
      \item $\mathrm{NLEACH}_t$: Nitrate leaching.
      \item $\mathrm{DENIT}_t$: Denitrification losses.
  \end{itemize}
\end{itemize}

\vspace{0.2cm}
\noindent\textbf{Method of Measurement.}
\begin{enumerate}
  \item \textbf{Water Management Reward ($R_{\text{water}}$)}
  \begin{equation}
  R_{\text{water},t}= w_m R_{\text{moisture},t} + w_s\,(1-\mathrm{SWFAC}_t)
  \end{equation}
  where
 \begin{equation}
R_{\text{moisture},t}=
\begin{cases}
+c_1,
& \text{if } \mathrm{SMR}_t \in [\mathrm{SMR}_{\text{low}},\,\mathrm{SMR}_{\text{high}}]
\\[2mm]
-c_2(\mathrm{SMR}_{\text{low}}-\mathrm{SMR}_t)^2,
& \text{if } \mathrm{SMR}_t < \mathrm{SMR}_{\text{low}}
\\[2mm]
-c_3(\mathrm{SMR}_t-\mathrm{SMR}_{\text{high}})^2,
& \text{if } \mathrm{SMR}_t > \mathrm{SMR}_{\text{high}}
\end{cases}
\end{equation}
  This component encourages maintaining soil moisture within an agronomically desirable
  range while minimising water stress.

  \item \textbf{Fertilisation Reward ($R_{\text{fert}}$)}
  \begin{equation}
  R_{\text{fert},t}= w_n\,(1-\mathrm{NSTRES}_t) + w_u\,\mathrm{TRNU}_t
  \end{equation}
  This reward promotes adequate nitrogen availability and efficient crop nitrogen
  uptake.

  \item \textbf{Resource Consumption Penalty ($R_{\text{resource}}$)}
  \begin{equation}
  R_{\text{resource},t}= w_{\text{irr}}\,\mathrm{IRRC}_t + w_{\text{fert}}\,\mathrm{FERN}_t
  \end{equation}
  This component discourages excessive irrigation and fertilisation while allowing
  necessary applications.

  \item \textbf{Environmental Loss Penalty ($R_{\text{losses}}$)}
  \begin{equation}
  R_{\text{losses},t}= w_r\,\mathrm{ROFC}_t + w_d\,\mathrm{DRNC}_t + w_l\,\mathrm{NLEACH}_t + w_{\text{dn}}\,\mathrm{DENIT}_t
  \end{equation}
  This penalty reduces environmentally harmful practices and resource wastage.
\end{enumerate}

\vspace{0.2cm}
\noindent\textbf{Daily Reward Computation.}
\begin{equation}
R_{\text{daily},t}= R_{\text{water},t} + R_{\text{fert},t} - R_{\text{resource},t} - R_{\text{losses},t}
\label{eq:daily_reward}
\end{equation}
The daily reward evaluates the quality of the management decision taken on day $t$
rather than final crop productivity. In the implemented \texttt{SmartFarmSoSEnv},
each sub-term is normalised to $[-1,1]$ or $[0,1]$ before weighting, and the
combined daily signal is scaled by $\alpha=0.01$ so that seasonal objectives
dominate at harvest.

%=========================================================
\subsection{REQ002: Optimise Seasonal Reward ($R_{\text{season}}$)}
%=========================================================

\noindent\textbf{Detailed Requirement.}
At harvest, the system shall compute a terminal reward that evaluates the overall
success of the growing season based on productivity, efficiency, and resource
utilisation.

\vspace{0.2cm}
\noindent\textbf{Rationale.}
The primary mission objective is to maximise crop yield while using water and fertiliser
efficiently. Therefore, the seasonal reward represents the dominant optimisation
objective.

\vspace{0.2cm}
\noindent\textbf{Associated Measure of Effectiveness.}
\begin{itemize}
    \item Grain yield.
    \item Harvest Index.
    \item Agronomic Nitrogen Efficiency.
    \item Seasonal water consumption.
    \item Seasonal fertiliser consumption.
\end{itemize}

\vspace{0.2cm}
\noindent\textbf{Operational Definition and Indicators.}
\begin{itemize}
  \item $\mathrm{GRNWT}_{\text{final}}$: Final grain yield measured at harvest.
  \item Harvest Index:
  \begin{equation}
  \mathrm{HIAD}_{\text{final}}= \frac{\text{Grain Weight}}{\text{Total Above-Ground Biomass}}
  \end{equation}
  \item Agronomic Nitrogen Efficiency:
  \begin{equation}
  \mathrm{ANE}_{\text{final}}= \frac{Y_N-Y_0}{N_{\text{applied}}}
  \end{equation}
  where:
  \begin{itemize}
      \item $Y_N$: Yield with nitrogen application.
      \item $Y_0$: Yield without nitrogen application.
      \item $N_{\text{applied}}$: Total nitrogen applied.
  \end{itemize}
  \item Seasonal irrigation:
  \begin{equation}
  \mathrm{IRRC}_{\text{season}}= \sum_{t=1}^{T} \mathrm{IRRC}_t
  \end{equation}
  \item Seasonal fertilisation:
  \begin{equation}
  \mathrm{FERN}_{\text{season}}= \sum_{t=1}^{T} \mathrm{FERN}_t
  \end{equation}
\end{itemize}

\vspace{0.2cm}
\noindent\textbf{Method of Measurement.}
The terminal reward is computed at harvest as:
\begin{equation}
R_{\text{season}}= w_y\,\mathrm{GRNWT}_{\text{final}} + w_h\,\mathrm{HIAD}_{\text{final}} + w_a\,\mathrm{ANE}_{\text{final}}
            - w_w\,\mathrm{IRRC}_{\text{season}} - w_f\,\mathrm{FERN}_{\text{season}}
\label{eq:season_reward}
\end{equation}
where:
\begin{itemize}
    \item $w_y$: Yield weight (dominant term).
    \item $w_h$: Harvest-index weight.
    \item $w_a$: Agronomic-nitrogen-efficiency weight.
    \item $w_w$: Water-consumption penalty weight.
    \item $w_f$: Fertiliser-consumption penalty weight.
\end{itemize}

\vspace{0.2cm}
\noindent\textbf{Final Transition Reward.}
For non-terminal days:
\begin{equation}
R_t = R_{\text{daily},t}
\end{equation}
For the harvest day:
\begin{equation}
R_T = R_{\text{daily},T} + R_{\text{season}}
\label{eq:transition_scalar}
\end{equation}
This hierarchical reward structure combines dense daily feedback with long-term seasonal
objectives, enabling efficient learning while preserving the primary goal of maximising
sustainable crop productivity.

%=========================================================
\subsection{From Scalar Rewards to a Multi-Reward Vector}
%=========================================================
The requirements above define the reward \emph{terms}; the multi-reward formulation
groups them into the objective channels that the MORL agents optimise jointly. In this
work we consider four objectives and train accordingly:
\begin{equation}
\mathbf{R}_t = \big(\, R^{\text{water}}_t,\; R^{\text{yield}}_t,\; R^{\text{nitrogen}}_t,\; R^{\text{balanced}}_t \,\big)^{\!\top}
\end{equation}
where $R^{\text{water}}$ aggregates the water-management and irrigation-penalty terms,
$R^{\text{yield}}$ is dominated by the seasonal productivity term, $R^{\text{nitrogen}}$
aggregates the fertilisation reward and nitrogen-loss penalties, and
$R^{\text{balanced}}$ is the scalarised composite that trades the three off against one
another. A scalarised return under a preference vector
$\boldsymbol{\omega}\in\Delta$ (the probability simplex) is then
\begin{equation}
G_{\boldsymbol{\omega}} = \sum_{t} \gamma^{t}\, \boldsymbol{\omega}^{\top}\mathbf{R}_t ,
\qquad \boldsymbol{\omega}\succeq 0,\ \ \textstyle\sum_i \omega_i = 1 .
\label{eq:scalarized}
\end{equation}
Single-objective baselines correspond to placing all preference mass on one channel,
while the MORL agents learn across the whole simplex. In the PC-PPO and CAPQL
implementations, the three primary corners map to yield priority, nitrogen-efficiency
priority, and water-efficiency priority; the balanced channel corresponds to uniform
preference sampling over the simplex.


%=========================================================
\subsection{Scalar Reward Variants for Specialised PPO Baselines}
%=========================================================

Beyond the balanced Base reward (\texttt{02\_smart\_farm\_env.py}), two scalar
variants reweight daily and seasonal terms to probe single-resource priorities
without changing the underlying REQ001--REQ002 structure. These variants train
separate PPO agents (Chapter~4); Table~\ref{tab:reward_variants} records the
parameter differences taken from the implementation files.

\begin{table}[ht]
\centering
\caption{Scalar reward variant parameters (Base, Water-min, N-min).}
\label{tab:reward_variants}
\small
\renewcommand{\arraystretch}{1.15}
\begin{tabular}{lccc}
\toprule
\textbf{Parameter} & \textbf{Base} & \textbf{Water-min} & \textbf{N-min} \\
\midrule
\multicolumn{4}{l}{\textit{Daily blend weights}} \\
\quad $W_{\text{WATER}}$    & 0.40 & 0.20 & 0.35 \\
\quad $W_{\text{FERT}}$     & 0.30 & 0.30 & 0.20 \\
\quad $W_{\text{RESOURCE}}$ & 0.15 & 0.35 & 0.30 \\
\quad $W_{\text{LOSSES}}$   & 0.15 & 0.15 & 0.15 \\
\midrule
\multicolumn{4}{l}{\textit{Seasonal blend weights}} \\
\quad $W_{\text{YIELD}}$    & 0.40 & 0.25 & 0.25 \\
\quad $W_{\text{HIAD}}$     & 0.10 & 0.10 & 0.10 \\
\quad $W_{\text{ANE}}$      & 0.30 & 0.30 & 0.40 \\
\quad $W_{\text{PENALTY}}$  & 0.20 & 0.35 & 0.25 \\
\midrule
\multicolumn{4}{l}{\textit{Seasonal penalty thresholds}} \\
\quad IRRIG\_EXCESS (mm)    & 500  & 150  & 500 \\
\quad N\_EXCESS (kg N ha$^{-1}$) & 200 & 200 & 100 \\
\midrule
\multicolumn{4}{l}{\textit{$R_{\text{penalty}}$ split (water / N)}} \\
\quad $\lambda_W : \lambda_N$ & 50:50 & 85:15 & 30:70 \\
\midrule
N-min yield floor on $R_{\text{ane}}$ & --- & --- & 5000 kg ha$^{-1}$ \\
\bottomrule
\end{tabular}
\end{table}

The Water-min variant reduces moisture-chasing and irrigation-excess tolerance;
the N-min variant strengthens N-application penalties and prioritises agronomic
N efficiency, with a yield floor preventing a zero-fertiliser exploit.


%=========================================================
\subsection{Implementation Constants and Normalisation}
%=========================================================

All sub-reward terms are clipped to $[-1,+1]$ or $[0,1]$ before weighting.
The combined daily signal is scaled by $\alpha=\texttt{DAILY\_SCALE}=0.01$.
Moisture-ratio thresholds are $\mathrm{SMR}_{\text{low}}=0.40$ and
$\mathrm{SMR}_{\text{high}}=0.80$. Seasonal yield, harvest index, and ANE are
normalised against \texttt{BASELINE\_YIELD}$=7620$,
\texttt{BASELINE\_HIAD}$=0.45$, and \texttt{BASELINE\_ANE}$=40$~kg grain per
kg~N. In the implemented Base seasonal reward, the positive terms use weights
$W_{\text{YIELD}}=0.40$, $W_{\text{HIAD}}=0.10$, $W_{\text{ANE}}=0.30$ and the
resource penalty uses $W_{\text{PENALTY}}=0.20$, grouped as
$R_{\text{season}} = W_{\text{YIELD}} R_{\text{yield}} + W_{\text{HIAD}} R_{\text{hiad}}
+ W_{\text{ANE}} R_{\text{ane}} - W_{\text{PENALTY}} R_{\text{penalty}}$.


%------------------------------------------------------------------------------
\section{Multi-Objective Learning Methodology}
\label{sec:algos}
%------------------------------------------------------------------------------

Three learner families implement the proposed solution: scalar PPO baselines
(fixed reward shaping), Preference-Conditioned PPO (PC-PPO), and
Concave-Augmented Pareto Q-Learning (CAPQL). Rule-based controllers provide
non-learning reference policies; their evaluation is deferred to Chapter~4.
Training hyperparameters (rollout length, learning rate, \texttt{VecNormalize})
are likewise reported in Chapter~4.

\subsection{Scalar PPO Baselines}

PPO~\cite{schulman2017ppo} maximises the clipped surrogate
\begin{equation}
\mathcal{L}^{\text{CLIP}}(\theta) =
\mathbb{E}_t\!\Big[\, \min\big( \rho_t(\theta)\,\hat{A}_t,\;
\operatorname{clip}(\rho_t(\theta),\,1-\epsilon,\,1+\epsilon)\,\hat{A}_t \big) \Big],
\label{eq:ppoclip}
\end{equation}
where $\hat{A}_t$ is the generalised advantage estimate~\cite{schulman2016gae}
and $\rho_t(\theta)=\pi_\theta(a_t\mid s_t)/\pi_{\theta_{\text{old}}}(a_t\mid s_t)$.
The scalar PPO agent uses the hierarchical reward of
REQ001--REQ002 (Equations~\eqref{eq:daily_reward}--\eqref{eq:transition_scalar}) with an
11-dimensional observation. Additional specialised variants reweight daily
components toward water saving or nitrogen saving (Table~\ref{tab:reward_variants}).
These agents provide rule-free reference policies against which the
multi-objective learners are compared in Chapter~4.
Algorithm~\ref{alg:scalar_ppo} summarises the on-policy training loop.

\begin{algorithm}[ht]
\caption{Scalar PPO training on \texttt{SmartFarmSoSEnv}}
\label{alg:scalar_ppo}
\begin{algorithmic}[1]
\Require Environment $\mathcal{E}$, total steps $N$, rollout length $L$, clip $\epsilon$, GAE $\lambda$, discount $\gamma$
\State Initialise policy $\pi_\theta$ and value network $V_\phi$; $\theta_{\text{old}} \leftarrow \theta$
\While{environment steps $< N$}
  \State Collect rollout $\mathcal{D}=\{(s_t,a_t,r_t)\}_{t=1}^{L}$ with $a_t \sim \pi_{\theta_{\text{old}}}(\cdot \mid s_t)$
  \State Compute returns $\hat{R}_t$ and advantages $\hat{A}_t$ using GAE($\gamma,\lambda$)
  \For{$k = 1$ to \texttt{n\_epochs}}
    \For{each minibatch $\mathcal{B} \subset \mathcal{D}$}
      \State $\rho_t(\theta) \leftarrow \pi_\theta(a_t \mid s_t)\,/\,\pi_{\theta_{\text{old}}}(a_t \mid s_t)$
      \State Update $\theta$ by ascending $\mathcal{L}^{\text{CLIP}}(\theta)$ on $\mathcal{B}$
      \State Update $\phi$ by minimising the value-function loss on $\mathcal{B}$
    \EndFor
  \EndFor
  \State $\theta_{\text{old}} \leftarrow \theta$
\EndWhile
\State \textbf{return} trained policy $\pi_\theta$
\end{algorithmic}
\end{algorithm}


\subsection{Preference-Conditioned PPO (PC-PPO)}
\label{subsec:pcppo}

PC-PPO is \emph{Preference-Conditioned}, not a separate constrained-optimisation
variant: the same PPO clipped objective is used, but the observation includes
$\boldsymbol{\omega}$ and the multi-reward vector is scalarised through
$\boldsymbol{\omega}^{\top}\mathbf{R}_t$ as in Equation~\eqref{eq:scalarized}.
Each episode samples a new preference $\boldsymbol{\omega}\sim\mathrm{Dirichlet}(1,1,1)$,
so one network learns a family of policies over the simplex. At evaluation,
the same model is queried at yield-, N-efficiency-, and water-efficiency-focused
corners $(1,0,0)$, $(0,1,0)$, $(0,0,1)$ and at a balanced profile
$(\tfrac{1}{3},\tfrac{1}{3},\tfrac{1}{3})$, matching the three-dimensional
terminal vector $\mathbf{R}_T=[R_{\text{yield}}, R_{\text{ane}}, R_{\text{water\_eff}}]^\top$
returned by \texttt{02\_smart\_farm\_env\_pcppo.py}.

PC-PPO uses the 14-dimensional observation
$[s_t;\boldsymbol{\omega}]$ with \texttt{VecNormalize} on observations and rewards.
Its main strength is stable on-policy learning on the costly DSSAT simulator while
still covering multiple farmer priorities with a single model.
Algorithm~\ref{alg:pcppo} details the episode-level preference sampling and
PPO update procedure.

\begin{algorithm}[ht]
\caption{Preference-Conditioned PPO (PC-PPO)}
\label{alg:pcppo}
\begin{algorithmic}[1]
\Require Environment $\mathcal{E}$, total steps $N$, rollout length $L$
\State Initialise $\pi_\theta$, $V_\phi$; $\theta_{\text{old}} \leftarrow \theta$
\While{environment steps $< N$}
  \State Sample $\boldsymbol{\omega} \sim \mathrm{Dirichlet}(\mathbf{1})$ for the upcoming episode
  \State Reset $\mathcal{E}$; $t \leftarrow 0$; \textbf{done} $\leftarrow$ \textbf{false}
  \State $\mathcal{D} \leftarrow \emptyset$
  \While{$|\mathcal{D}| < L$ \textbf{and not} \textbf{done}}
    \State $\tilde{s}_t \leftarrow [s_t;\,\boldsymbol{\omega}]$
    \State $a_t \sim \pi_{\theta_{\text{old}}}(\cdot \mid \tilde{s}_t)$
    \State Execute $(s_{t+1}, r_t, \textbf{done}) \leftarrow \mathcal{E}.\textsc{Step}(a_t)$
    \State Store $(\tilde{s}_t, a_t, r_t)$ in $\mathcal{D}$; $t \leftarrow t + 1$
  \EndWhile
  \State Compute GAE advantages $\{\hat{A}_t\}$ and returns $\{\hat{R}_t\}$ on $\mathcal{D}$
  \For{$k = 1$ to \texttt{n\_epochs}}
    \For{each minibatch $\mathcal{B} \subset \mathcal{D}$}
      \State Update $\theta, \phi$ with clipped PPO loss (Equation~\eqref{eq:ppoclip})
    \EndFor
  \EndFor
  \State $\theta_{\text{old}} \leftarrow \theta$
\EndWhile
\State \textbf{return} $\pi_\theta$
\end{algorithmic}
\end{algorithm}


\subsection{Concave-Augmented Pareto Q-Learning (CAPQL)}
\label{subsec:capql}

CAPQL~\cite{lu2023capql} takes an off-policy route based on TD3. A twin
preference-conditioned critic learns vector-valued returns, while the actor
outputs deterministic daily actions. The terminal target uses linear
scalarisation plus a concave augmentation term:
\begin{equation}
  y = R_{\text{daily}} + \mathbb{1}_{t=T}\!\left(
    \boldsymbol{\omega}^{\top}\mathbf{R}_T
    + \lambda \cdot \mathrm{mean}_i\big[\log(R_i + 1 + \varepsilon)\big]
  \right),
  \label{eq:capql_target}
\end{equation}
with $\lambda=0.5$ and $\varepsilon=0.1$ in our implementation. Hindsight
preference relabeling resamples $\boldsymbol{\omega}$ for minibatch
transitions, improving coverage of the objective space without rerunning
DSSAT episodes.

CAPQL is attractive when many preference values must be explored from limited
simulation budget: one network can be queried at any $\boldsymbol{\omega}$ after
training.
Algorithm~\ref{alg:capql} summarises the TD3 interaction loop, replay storage,
hindsight preference relabeling, and concave-augmented critic--actor updates.

\begin{algorithm}[ht]
\caption{Concave-Augmented Pareto Q-Learning (CAPQL)}
\label{alg:capql}
\begin{algorithmic}[1]
\Require Environment $\mathcal{E}$, replay buffer $\mathcal{B}$, total steps $N$, batch size $B$, policy delay $d$
\State Initialise actor $\mu_\theta$, critics $Q_{\phi_1}, Q_{\phi_2}$, and target networks $\mu_{\bar\theta}, Q_{\bar\phi}$
\For{$n = 1$ to $N$}
  \If{episode start}
    \State Sample $\boldsymbol{\omega} \sim \mathrm{Dirichlet}(\mathbf{1})$; reset $\mathcal{E}$
  \EndIf
  \State $\tilde{s}_t \leftarrow [s_t;\,\boldsymbol{\omega}]$
  \State $a_t \leftarrow \mu_\theta(\tilde{s}_t) + \text{exploration noise}$; step $\mathcal{E}$
  \State Store $(s_t, a_t, R_{\text{daily},t}, \mathbf{R}_T\,\mathbb{1}_{t=T}, s_{t+1}, \textbf{done}, \boldsymbol{\omega})$ in $\mathcal{B}$
  \If{$|\mathcal{B}| \ge B$ \textbf{and} $n > \texttt{warmup}$}
    \State Sample minibatch; relabel $\boldsymbol{\omega}' \sim \mathrm{Dirichlet}(\mathbf{1})$ per transition
    \State $\tilde{s} \leftarrow [s;\boldsymbol{\omega}']$,\; $\tilde{s}' \leftarrow [s';\boldsymbol{\omega}']$
    \State $r_{\text{aug}} \leftarrow R_{\text{daily}} + \textbf{done}\cdot\big(\boldsymbol{\omega}'^{\top}\mathbf{R}_T + \lambda\,\mathrm{mean}_i\log(R_i{+}1{+}\varepsilon)\big)$
    \State $a' \leftarrow \mathrm{clip}\big(\mu_{\bar\theta}(\tilde{s}') + \text{smoothing noise}\big)$
    \State $y \leftarrow r_{\text{aug}} + \gamma(1-\textbf{done})\min_{j \in \{1,2\}} Q_{\bar\phi_j}(\tilde{s}', a')$
    \State Update $\phi_1, \phi_2$ by minimising $\|Q_{\phi_j}(\tilde{s},a) - y\|^2$
    \If{$n \bmod d = 0$}
      \State Update $\theta$ by ascending $\mathbb{E}\big[\min_j Q_{\phi_j}(\tilde{s}, \mu_\theta(\tilde{s}))\big]$
      \State Soft-update $\bar\theta \leftarrow \tau\theta + (1{-}\tau)\bar\theta$;\;
            $\bar\phi \leftarrow \tau\phi + (1{-}\tau)\bar\phi$
    \EndIf
  \EndIf
\EndFor
\State \textbf{return} $\mu_\theta$
\end{algorithmic}
\end{algorithm}


\subsection{Justification of Algorithmic Choices}

\begin{table}[ht]
\centering
\caption{Rationale for the selected learning algorithms.}
\label{tab:algos}
\small
\renewcommand{\arraystretch}{1.3}
\begin{tabular}{>{\raggedright\arraybackslash}p{2.8cm}
                >{\raggedright\arraybackslash}p{2.6cm}
                >{\raggedright\arraybackslash}p{6.3cm}}
\toprule
\textbf{Method} & \textbf{Family} & \textbf{Why chosen here} \\
\midrule
Scalar PPO & On-policy, single-objective &
Stable baseline on \texttt{gym-DSSAT}; hierarchical reward matches REQ000--REQ002;
provides fixed-reward reference policies for comparison. \\
PC-PPO & On-policy, preference-conditioned &
Same PPO core as the baseline; one model serves all preference corners;
observation carries $\boldsymbol{\omega}$; strong training stability. \\
CAPQL & Off-policy, preference-conditioned &
Sample-efficient exploration of the Pareto set; concave augmentation mitigates
plain linear scalarisation; TD3 twin critics reduce Q overestimation. \\
\bottomrule
\end{tabular}
\end{table}

In short, scalar PPO defines what fixed hierarchical rewards can encode
(Algorithm~\ref{alg:scalar_ppo}); PC-PPO extends the same PPO core to a
preference-conditioned policy family (Algorithm~\ref{alg:pcppo}); CAPQL
complements this with an off-policy Pareto-learning formulation
(Algorithm~\ref{alg:capql}).


%------------------------------------------------------------------------------
\section{Integrated Proposed Architecture}
\label{sec:overall}
%------------------------------------------------------------------------------

Figure~\ref{fig:overall} assembles the elements of this chapter into the complete
closed-loop learning system. The DSSAT core, wrapped by \texttt{gym-DSSAT} and
\texttt{SmartFarmSoSEnv}, exposes the daily crop state. The reward module
converts DSSAT outputs into daily and seasonal signals (Section~\ref{sec:reward}).
The agent---scalar PPO, PC-PPO, or CAPQL---consumes state and reward and emits
irrigation and nitrogen actions, closing the loop over one maize season.

%---- FIGURE : Overall proposed model ----------------------------------------
\begin{figure}[ht]
\centering
\resizebox{\textwidth}{!}{%
\begin{tikzpicture}[node distance=1.0cm and 1.8cm]

  \node[core, draw=sosblue, fill=sosblue!10, minimum width=3.4cm, minimum height=1.6cm]
       (env) {gym-DSSAT\\Environment\\\scriptsize (DSSAT CSM core)};

  \node[core, draw=sosgreen, fill=sosgreen!10, below=2.2cm of env,
        minimum width=3.4cm, minimum height=1.6cm] (reward)
        {Reward Module\\\scriptsize $R_{\text{daily}}+R_{\text{season}}$};

  \node[core, draw=sosorange, fill=soscream, right=4.2cm of env,
        minimum width=3.6cm, minimum height=1.6cm] (agent)
        {Multi-Objective RL Agent\\\scriptsize PPO $\,\cdot\,$ PC-PPO $\,\cdot\,$ CAPQL};

  \node[block, draw=sosgreen, fill=sosgreen!8, below=2.2cm of agent,
        minimum width=3.6cm, minimum height=1.5cm, font=\scriptsize, align=center]
        (vec) {Seasonal vector $\mathbf{R}_T$\\yield $\cdot$ N-efficiency\\water efficiency};

  \draw[flow] (env.east) -- node[above,font=\scriptsize]{state $s_t$} (agent.west);
  \draw[ctrl] (agent.south) |- node[below right,font=\scriptsize,pos=0.25]
        {action $a_t=(\mathrm{FERN}_t,\mathrm{IRRC}_t)$} (env.south);
  \draw[flow] (env.south) -- node[left,font=\scriptsize]{DSSAT outputs} (reward.north);
  \draw[flow] (reward.east) -- (vec.west);
  \draw[flow] (vec.north) -- node[right,font=\scriptsize]{$\mathbf{R}_T$} (agent.south);

  \node[block, draw=sosgrey, fill=white, right=1.2cm of agent, font=\scriptsize,
        align=center, minimum width=2.6cm] (eval)
        {Evaluation\\\scriptsize vs.\ rule-based\\and scalar PPO baselines};
  \draw[flow, sosgrey] (agent.east) -- (eval.west);

  \begin{scope}[on background layer]
    \node[draw=sosgrey, dashed, rounded corners, inner sep=8pt,
          fit=(env)(reward)(agent)(vec),
          label={[sosgrey,font=\scriptsize\bfseries]above:Closed RL loop (one growing season)}] {};
  \end{scope}

\end{tikzpicture}%
}
\caption{Overall proposed learning architecture. The DSSAT core supplies daily
  crop state; the reward module computes hierarchical daily and seasonal
  signals; the agent issues irrigation and nitrogen actions, closing the loop
  over one maize season.}
\label{fig:overall}
\end{figure}


\subsection{Traceability from Requirements to Implementation}

Table~\ref{tab:trace} summarises how REQ000--REQ002 connect mission objectives,
MoEs, reward terms, and learner families. This traceability matrix is the
design counterpart to the performance tables reported in Chapter~4.

\begin{table}[ht]
\centering
\caption{Traceability: requirements, MoEs, rewards, and learners.}
\label{tab:trace}
\small
\renewcommand{\arraystretch}{1.2}
\begin{tabular}{>{\raggedright\arraybackslash}p{1.6cm}
                >{\raggedright\arraybackslash}p{2.8cm}
                >{\raggedright\arraybackslash}p{2.6cm}
                >{\raggedright\arraybackslash}p{2.4cm}
                >{\raggedright\arraybackslash}p{2.2cm}}
\toprule
\textbf{Req.} & \textbf{MoE / objective} & \textbf{Reward term} & \textbf{Module} & \textbf{Learner} \\
\midrule
REQ000 & Yield, N, water & $R_{\text{daily}}+R_{\text{season}}$ &
\texttt{02\_smart\_farm\_env*} & PPO, PC-PPO, CAPQL \\
REQ001 & SMR, SWFAC, NSTRES & $R_{\text{water}}, R_{\text{fert}}$ &
\texttt{02\_smart\_farm\_env.py} & All \\
REQ001 & Resource use, losses & $R_{\text{resource}}, R_{\text{losses}}$ &
\texttt{02\_smart\_farm\_env.py} & All \\
REQ002 & GRNWT, HIAD, ANE & $R_{\text{yield}}, R_{\text{hiad}}, R_{\text{ane}}$ &
\texttt{02\_smart\_farm\_env.py} & PPO \\
REQ002 & Season totals & $R_{\text{penalty}}$ &
\texttt{02\_smart\_farm\_env*.py} & PPO variants \\
MORL   & Preference trade-offs & $\mathbf{R}_T$, $\boldsymbol{\omega}$ &
\texttt{02\_smart\_farm\_env\_pcppo.py} & PC-PPO, CAPQL \\
\bottomrule
\end{tabular}
\end{table}


\subsection{Comparison Baselines (Design)}

For evaluation (Chapter~4), four rule-based controllers run inside the same
\texttt{SmartFarmSoSEnv} wrapper: \emph{Random} (unconstrained daily inputs),
\emph{Fixed} (preset N splits and periodic irrigation), \emph{Stage-based}
(management keyed to vegetative development), and \emph{FAO-56} (soil-water-balance
irrigation with a three-split N plan). Using a common environment ensures that
MoEs and reward components are directly comparable across learned and rule-based
policies.


%------------------------------------------------------------------------------
\section{Conclusion}
\label{sec:ch3_conclusion}
%------------------------------------------------------------------------------

This chapter presented the proposed solution end to end: a smart-farming
System-of-Systems formalised around the DSSAT Cropping System Model; an MDP
formulation with explicit observation, action, and software layers;
Measures of Effectiveness tied to mission objectives (Table~\ref{tab:moe});
a hierarchical reward design (REQ000--REQ002) with scalar variants
(Table~\ref{tab:reward_variants}) and a multi-objective extension for
preference-conditioned learning; and three learner families---scalar PPO,
PC-PPO, and CAPQL---integrated in the closed-loop architecture of
Figure~\ref{fig:overall}. The next chapter reports the experimental setup,
training dynamics, and quantitative comparison of all learned and rule-based
policies.


% =============================================================================
% If using BibTeX instead, delete this block and add entries to references.bib
% =============================================================================
\begin{thebibliography}{99}

\bibitem{maier1998}
M.~W. Maier,
``Architecting principles for systems-of-systems,''
\emph{Systems Engineering}, vol.~1, no.~4, pp.~267--284, 1998.

\bibitem{jones2003dssat}
J.~W. Jones, G.~Hoogenboom, C.~H. Porter, K.~J. Boote, W.~D. Batchelor, L.~A. Hunt,
P.~W. Wilkens, U.~Singh, A.~J. Gijsman, and J.~T. Ritchie,
``The DSSAT cropping system model,''
\emph{European Journal of Agronomy}, vol.~18, no.~3--4, pp.~235--265, 2003.

\bibitem{hoogenboom2019}
G.~Hoogenboom, C.~H. Porter, K.~J. Boote, V.~Shelia, P.~W. Wilkens, U.~Singh,
J.~W. White, S.~Asseng, J.~I. Lizaso, et~al.,
``The DSSAT crop modeling ecosystem,'' in
\emph{Advances in Crop Modelling for a Sustainable Agriculture}, Burleigh Dodds Science
Publishing, pp.~173--216, 2019.

\bibitem{gautron2022gymdssat}
R.~Gautron, E.~J. Padr\'on, P.~Preux, J.~Bigot, O.-A. Maillard, and D.~Emukpere,
``gym-DSSAT: a crop model turned into a Reinforcement Learning environment,''
Inria Research Report RR-9460, 2022.

\bibitem{gautron2022review}
R.~Gautron, O.-A. Maillard, P.~Preux, M.~Corbeels, and R.~Sabbadin,
``Reinforcement learning for crop management support: Review, prospects and challenges,''
\emph{Computers and Electronics in Agriculture}, vol.~200, art.~107182, 2022.

\bibitem{schulman2017ppo}
J.~Schulman, F.~Wolski, P.~Dhariwal, A.~Radford, and O.~Klimov,
``Proximal policy optimization algorithms,''
\texttt{arXiv:1707.06347}, 2017.

\bibitem{schulman2016gae}
J.~Schulman, P.~Moritz, S.~Levine, M.~Jordan, and P.~Abbeel,
``High-dimensional continuous control using generalized advantage estimation,'' in
\emph{Proc.\ ICLR}, 2016.

\bibitem{lu2023capql}
H.~Lu, D.~Herman, and Y.~Yu,
``Multi-objective reinforcement learning: Convexity, stationarity and Pareto optimality,''
in \emph{Proc.\ ICLR}, 2023.

\end{thebibliography}
